% Public paper source; file identities are recorded in MANIFEST.json.
\documentclass[11pt,reqno]{amsart}
\usepackage[T1]{fontenc}
\usepackage[utf8]{inputenc}
\usepackage{lmodern}
\usepackage{amsmath,amssymb}
\usepackage[letterpaper,textwidth=6in,textheight=8.7in,centering]{geometry}
\usepackage[expansion=false]{microtype}
\usepackage{booktabs,array,longtable}
\usepackage{needspace}
\usepackage{fvextra}
\usepackage{xurl}
\usepackage[hidelinks,unicode]{hyperref}
\newtheorem{theorem}{Theorem}
\newtheorem{corollary}[theorem]{Corollary}
\newtheorem{lemma}[theorem]{Lemma}
\newtheorem{proposition}[theorem]{Proposition}
\theoremstyle{definition}
\newtheorem*{definition}{Definition}
\theoremstyle{remark}
\newtheorem*{remark}{Remark}
\providecommand{\tightlist}{\setlength{\itemsep}{0pt}\setlength{\parskip}{0pt}}
\setlength{\emergencystretch}{1em}
\raggedbottom
\hypersetup{pdftitle={Finite graph certificates for composite terms in rational floor sequences},
  pdfauthor={Yuri Odagiri}}
\pdfinfoomitdate=1
\pdftrailerid{}
\pdfsuppressptexinfo=15
\listfiles
\title[Finite graph certificates]{Finite graph certificates for composite terms in rational floor sequences}
\author{Yuri Odagiri}
\address{Yuri Odagiri: Independent researcher}
\email{ixixizko@gmail.com}
\date{}
% BEGIN PUBLIC ADDITION figure-preamble
\newcommand{\DoNotLoadEpstopdf}{}
% Shared paper illustrations. Load in the preamble.
\usepackage{tikz}
\usetikzlibrary{arrows.meta,positioning,calc,patterns,shapes.geometric}
\providecommand{\figtext}[2]{\ifdefined\FigureJapanese#2\else#1\fi}
\definecolor{figblue}{RGB}{30,70,110}
\definecolor{figgray}{RGB}{105,105,105}
\tikzset{
  paperfig/.style={font=\fontsize{9}{11}\selectfont,>=Stealth,line width=.65pt},
  fbox/.style={draw,rounded corners=2pt,align=center,inner sep=5pt,fill=white},
  fnode/.style={draw,circle,minimum size=6mm,inner sep=1pt,fill=white},
  fedge/.style={->},
  femph/.style={draw=figblue,line width=1.05pt},
  flabel/.style={fill=white,inner sep=1.5pt,align=center},
  ftitle/.style={font=\fontsize{9.5}{12}\selectfont\bfseries,align=center}
}

% END PUBLIC ADDITION figure-preamble

% BEGIN PUBLIC ADDITION doi
\thanks{DOI: \url{https://doi.org/10.5281/zenodo.22962354}.}
\hypersetup{pdfsubject={DOI: 10.5281/zenodo.22962354}}
% END PUBLIC ADDITION doi

\begin{document}
% Markdown AST block 0: Finite graph certificates for composite terms in rational floor sequences
% Title appears in front matter.

% Markdown AST block 1: Abstract
\begin{abstract}

% Markdown AST block 2: Para
For every real number \(\xi>0\), we give computer-assisted proofs that
\(\lfloor\xi(7/5)^n\rfloor\) is divisible by at least one of
\(2,3,5,11,13\) for infinitely many positive integers \(n\), and that
\(\lfloor\xi(5/2)^n\rfloor\) is divisible by at least one of
\(2,3,7,11,13,17,19,23,29,31\) for infinitely many positive integers
\(n\). Consequently both sequences contain infinitely many composite
terms. The proofs represent an orbit avoiding the given primes by a walk in
a finite labelled graph whose vertices record residues and fractional
cells and whose edges carry signed carries. Components whose output
depends only on a cyclic phase are excluded by an elementary
nonperiodicity lemma. In the word-labelled construction the edges carry
finite carry words; deterministic paths are contracted exactly, and each
new prime is imposed at every intermediate time of a word. For the ratio
\(7/5\) we also prove an explicit waiting-time bound: if
\(P=\lfloor\xi\rfloor>13\) then some term with index at most
\(428+12L(P)\), \(L(P)=\min\{h\ge0:5^h\ge P+2\}\), is divisible by one
of the five primes. We then determine the reach of this certificate
method: contraction, prime powers, and the primes dividing \(2ab\) do not
enlarge the class of finitely certifiable bases; the method cannot
terminate when \(a\ge2\operatorname{rad}M\); and the product \(R\) of the
auxiliary primes not dividing \(2ab\) must satisfy
\(a\le2R\,J(N_0)\) for the Jacobsthal function \(J\), whence
\(R\ge a^{1-o(1)}\) by an elementary bound. Every adopted finite
computation was performed by two implementations with disjoint scientific
cores, which agree element by element on all compared finite data.
Lean proofs cover both divisibility theorems and their compositeness
corollaries: the one-step proof for \(7/5\) is kernel-checked, and the
word-labelled proofs combine kernel-checked soundness theorems with
\nolinkurl{native_decide} finite evaluations, which add trust in the compiler
and native evaluation.

% Markdown AST block 3: 1. Introduction
\end{abstract}
\maketitle
\enlargethispage{4pt}

\Needspace{6\baselineskip}
\section{Introduction}\label{sec:1}

% Markdown AST block 4: Para
The arithmetic of the integer parts of rational powers presents difficulties
even for simple nonintegral bases. A useful way to prove that a sequence
contains infinitely many composite terms is to exhibit a fixed finite set of
primes such that infinitely many terms have a divisor in that set. Once the
terms exceed all those primes, divisibility implies compositeness.
Following Dubickas \hyperref[ref-dub06]{Dub06, §2}, such a set is called
\emph{unavoidable} with respect to a base \(b\) if it has this property for
\(\lfloor\xi b^n\rfloor\) and every \(\xi>0\). Among integer bases, such
sets exist for \(b\in\{2,3,4,6\}\); none exists for \(b=5\)
(\hyperref[ref-dn]{DN}, as recalled in \hyperref[ref-dub06]{Dub06, §2}) or
for any \(b\ge3\) with \(b-1\) not squarefree
(\hyperref[ref-dub06]{Dub06, Theorem 4}), and none is expected for the
other integer bases (\hyperref[ref-dub06]{Dub06, Conjecture 3}). In this
terminology, the results recalled next give unavoidable sets for the
rational bases \(3/2\), \(4/3\), and \(5/4\), and Theorem 1 gives
unavoidable sets for \(7/5\) and \(5/2\).

% Markdown AST block 5: Para
In 1967, Forman and Shapiro \hyperref[ref-fs]{FS} proved that the integer
parts of the powers of \(3/2\) and \(4/3\) contain infinitely many
composite numbers; see the historical account in
\hyperref[ref-dn]{DN, pp.635--636}. As noted there, their proofs also apply
after multiplication by any fixed positive coefficient \(\xi\).
Dubickas and Novikas treated the additional base \(5/4\) and stated
explicit divisibility results for all three bases: for every \(\xi>0\),
infinitely many terms of each of
\[
\lfloor\xi(3/2)^n\rfloor,\qquad
\lfloor\xi(4/3)^n\rfloor,\qquad
\lfloor\xi(5/4)^n\rfloor
\]
have a divisor in the respective prime sets
\[
\{2,5,7,11\},\qquad \{2,3,5\},\qquad \{2,3,7,11,13\};
\]
see \hyperref[ref-dn]{DN, Theorem 1, p.636} and
\hyperref[ref-n]{N, Theorem 1.1, p.24}.

% Markdown AST block 6: Para
Dubickas and Novikas also proved the
corresponding assertion with primes \(\{2,3,5,11\}\) for the \emph{nearest
integers}
\(\lfloor\xi(7/5)^n+1/2\rfloor\): this is the third assertion of
\hyperref[ref-dn]{DN, Theorem 4, p.637}, also stated as
\hyperref[ref-n]{N, Theorem 1.5(iii), p.26}. For the ratio \(5/2\) they proved
that, for every \(\xi>0\), the \emph{shifted} sequence
\(\lfloor\xi(5/2)^n\rfloor-1\) has infinitely many terms divisible by
\(2\), \(3\) or \(5\); see \hyperref[ref-dn]{DN, Theorem 3, p.637} and
\hyperref[ref-n]{N, Theorem 1.4, p.25}; inside the proof, the same argument is
noted to apply to \(\lfloor\xi(5/2)^n\rfloor-1+30k\) for every fixed
integer \(k\) (\hyperref[ref-dn]{DN, p.646}, \hyperref[ref-n]{N, p.26}).
The rounding and shift conventions matter here. The present results
concern the unshifted integer parts. For \(5/2\) the shift changes the
alphabet of admissible carries, and our certificate uses a different
prime set.

% Markdown AST block 7: Theorem 1 (Main Theorem).
\begin{samepage}
\begin{theorem}[Main Theorem]\leavevmode\label{thm:1}
For every \(\xi\in\mathbb R\) with
\(\xi>0\) and every \(N\in\mathbb N\):

% Markdown AST block 8: OrderedList
\begin{enumerate}
\def\labelenumi{(\roman{enumi})}
\item
  there is an integer \(n\ge\max(N,1)\) such that
  \[
  \gcd\!\left(\left\lfloor\xi(7/5)^n\right\rfloor,4290\right)>1,
  \qquad 4290=2\cdot3\cdot5\cdot11\cdot13;
  \tag{1}\label{eq:1}
  \]
\item
  there is an integer \(n\ge\max(N,1)\) such that
  \[
  \begin{aligned}
  \gcd\!&\left(\left\lfloor\xi(5/2)^n\right\rfloor,40112098026\right)>1,\\
  40112098026&=2\cdot3\cdot7\cdot11\cdot13\cdot17\cdot19\cdot23\cdot29\cdot31.
  \end{aligned}
  \tag{2}\label{eq:2}
  \]
\end{enumerate}
\label{thm-end:1}
\end{theorem}
\end{samepage}

% Markdown AST block 9: Para
The indices in (i) and (ii) are quantified separately; no common index
is asserted. The modulus in (2) does not contain the prime \(5\).

% Markdown AST block 10: Para
Here and below \(\mathbb N=\{0,1,2,\ldots\}\), and for an integer \(q\)
we use \(\gcd(q,M)=\gcd(|q|,M)\). The formulation with arbitrary \(N\)
states precisely that the qualifying positive indices form an infinite set.
A positive integer is \emph{composite} if it equals \(uv\) with integers
\(u,v>1\).

% Markdown AST block 11: Corollary 2.
\begin{samepage}
\begin{corollary}\leavevmode\label{thm:2}
For every \(\xi>0\), each of the sequences
\[
\bigl(\lfloor\xi(7/5)^n\rfloor\bigr)_{n\ge1}\qquad \text{and}\qquad \bigl(\lfloor\xi(5/2)^n\rfloor\bigr)_{n\ge1}
\]
contains infinitely
many composite terms.
\label{thm-end:2}
\end{corollary}
\end{samepage}

% Markdown AST block 12: Para
For the ratio \(7/5\) the certificate also yields an explicit waiting
time. Write
\[
L(P)=\min\{h\ge0:5^h\ge P+2\}.
\]

% Markdown AST block 13: Theorem 3 (Quantitative Theorem).
\begin{samepage}
\begin{theorem}[Quantitative Theorem]\leavevmode\label{thm:3}
Let \(\xi>0\) with
\(P=\lfloor\xi\rfloor>13\). Then there is an integer \(n\) with
\[
\begin{aligned}
0&\le n\le428+12L(P)\\
\text{and}\qquad
\gcd\!&\left(\left\lfloor\xi(7/5)^n\right\rfloor,4290\right)>1,
\end{aligned}
\tag{3}\label{eq:3}
\]
and this term is composite. Consequently, for every \(N\in\mathbb N\)
with \(Q_N=\lfloor\xi(7/5)^N\rfloor>13\), some
\(\lfloor\xi(7/5)^n\rfloor\) with
\[
N\le n\le N+428+12L(Q_N)
\]
is
composite.
\label{thm-end:3}
\end{theorem}
\end{samepage}

% Markdown AST block 14: Para
The related study of fractional parts includes Mahler's question
\hyperref[ref-m68]{M68, p.313} whether some \(\xi>0\) satisfies
\[
0\le\{\xi(3/2)^n\}<1/2\qquad\text{for all }n\ge0.
\]
Flatto, Lagarias, and Pollington
\hyperref[ref-flp]{FLP, Theorem 1.4, p.129} proved that, for coprime integers
\(p>q\ge2\) and every \(\xi>0\),
\[
\limsup_{n\to\infty}\{\xi(p/q)^n\}
-\liminf_{n\to\infty}\{\xi(p/q)^n\}\ge1/p.
\]
Dubickas \hyperref[ref-dub08]{Dub08} obtained further results on whether all
these fractional parts can lie in prescribed intervals or unions of
intervals. These confinement questions provide a background for the use
of fractional cells here. The range bound alone does not yield the
infinitely many even terms that would suffice for compositeness, as
explained in \hyperref[ref-dn]{DN, p.636}.

% Markdown AST block 15: Para
For other real bases, Alkauskas and Dubickas \hyperref[ref-ad04]{AD04}
constructed Pisot numbers whose positive powers all have composite
integer parts, as well as a transcendental number whose integer-part
sequence meets every integer arithmetic progression infinitely often.
Those results construct particular bases; here the base is a specified
rational number and the coefficient \(\xi\) ranges over all positive reals.

% Markdown AST block 16: Para
\textbf{Method.} The proofs combine an existing nonperiodicity obstruction with
explicit finite certificates. Nonperiodicity of the carry sequence
\(c_n=bQ_{n+1}-aQ_n\) is already contained in
\hyperref[ref-dn]{DN, Lemma 2, p.638} and \hyperref[ref-n]{N, Lemma 1.7, pp.28--29}, and
the use of restrictions on operation words to force periodicity also has
precedents in \hyperref[ref-dn]{DN, pp.638--639 and 647}. We include an integral
proof of the obstruction and make the graph construction, the all-edge
phase test, and the tail-preservation arguments explicit. The one-step construction (Section 3) records one carry per edge and refines
the fractional cells; it proves (1). The word-labelled construction (Section 4)
labels edges by finite carry words, contracts deterministic paths without
changing their output, and lifts the graph by one new prime at a time
while checking the residues at every intermediate time of each word; it
proves (2) and gives a second proof of (1). A finite rank certificate on
the complete one-step graph, combined with a finite version of the
nonperiodicity lemma, gives Theorem 3 (Section 6).

% Markdown AST block 17: Para
Akiyama, Frougny, and Sakarovitch \hyperref[ref-afs08]{AFS08} developed
rational-base representations with finite transducers for arithmetic.
Their nonperiodicity results use divisibility by growing powers of the
denominator (\hyperref[ref-afs08]{AFS08, Lemma 8, p.65 and Proposition 26, p.75}).
Their language of integer representations is nonregular
(\hyperref[ref-afs08]{AFS08, Corollary 7, p.65}). Our finite graphs instead
record necessary residue and fractional-cell constraints on a
prime-avoiding orbit: every such orbit gives a walk, while a walk need
not represent an orbit. This distinction allows finite certificates
without requiring a finite automaton for the full representation language.

% Markdown AST block 18: Para
Stephan's preprint \hyperref[ref-s26]{S26}, dated August 15, 2026, proves
the case \(b=7\) of \hyperref[ref-dub06]{Dub06, Conjecture 2}: for every
\(\xi>0\), the sequence \(\lfloor\xi7^n\rfloor\) contains infinitely many
composite terms; by the same certificate, there is no infinite right
truncatable prime in base \(7\). Its method is closely related to the one
used here. The vertices of its graph are the residues modulo
\(M=\prod_{p\le31}p\) that are coprime to \(M\), with an edge for each
appended digit. The certificate states that, after the states not lying on
bi-infinite paths are discarded, no strongly connected component contains
two distinct cycles. Every infinite path is then ultimately periodic, and
an ultimately periodic digit word gives infinitely many composite terms by
a return theorem of Dubickas \hyperref[ref-dub09]{Dub09, Theorem 4},
recalled after Theorem 46 below. The certificate is reached through a
ladder of moduli that adjoins one prime at a time, using projections and
rank functions, and it is verified in Lean, where the finite computation
is evaluated by \nolinkurl{native_decide} (\hyperref[ref-s26]{S26, §§4--6
and Appendix A}).

% Markdown AST block 18 (continued): Para
The present work shares several ingredients with \hyperref[ref-s26]{S26}:
finite residue graphs in which every hypothetical prime-avoiding orbit
gives an infinite walk, the lifting of a certified graph by one new prime
at a time (Section 4.4), rank functions that a checker verifies in
place of a formal proof of graph algorithms (Sections 6.2 and 9.4), and a
Lean verification whose word-labelled finite evaluations are likewise
accepted through \nolinkurl{native_decide} (Section 9.4). The
constructions differ in three respects. First, the graph of
\hyperref[ref-s26]{S26} consists of residues alone. For a nonintegral base
not every carry word is realized by an orbit; the states here also record
fractional cells, and the graphs are outer approximations: every orbit
gives a walk, but a walk need not represent an orbit. Second, both
certificates single out components in which every infinite walk has a
periodic label sequence (simple cycles in \hyperref[ref-s26]{S26};
components passing the phase criterion of Lemma 7 here), but they treat
them in opposite ways. In base \(7\), such components carry actual
prime-avoiding orbits: the core of the graph of \hyperref[ref-s26]{S26}
at the modulus \(M\) still has 371702 states, and
\hyperref[ref-dj22]{DJ22, Theorem 1.2, p.167} implies that every finite
set of primes can be avoided by all terms for a suitable positive
coefficient, so no unavoidable set exists; \hyperref[ref-s26]{S26} keeps
these components and applies the return theorem to their periodic words.
For a nonintegral base, Lemma 4 shows that no orbit remains in such a
component, so it is deleted; the certificates here reduce the graph to
the empty graph and thereby yield the unavoidable sets of Theorem 1. The
result of \hyperref[ref-dj22]{DJ22} does not apply to \(7/5\) or \(5/2\).
Third, the word-labelled contraction with prime conditions imposed at
every intermediate time, the waiting-time bound of Theorem 3, and the
normal-form and obstruction theorems of Sections 7 and 8 are not part of
\hyperref[ref-s26]{S26}. The computation of \hyperref[ref-s26]{S26} is a
methodological comparison, not a dependency of the proofs given here.
The specific contributions here are the finite certificates for the
unshifted \(7/5\) and \(5/2\) sequences with their exact verification,
the word-labelled contraction with prime conditions imposed at every
intermediate time, the waiting-time bound, and the structural results
described next.

% Markdown AST block 19: Para
\textbf{Scope of the method.} Sections 7 and 8 study what this certificate
method can and cannot do. With the operations fixed precisely (initial
complete graph, phase removal, exact contraction, cell refinement,
modulus enlargement, carry restriction), we prove a normal-form theorem:
a finite certificate exists for some choice of parameters if and only if,
for some cell count \(K\) and some finite set \(S\) of primes not dividing
\(2ab\), every cyclic strongly connected component of one explicit graph
passes the phase test. Contraction does not enlarge the certifiable
class, prime powers in the modulus can be dropped, and the primes dividing
\(2ab\) reduce to a restriction of the carry alphabet; the normal form
changes the prime set of the conclusion, which then contains all primes
dividing \(2ab\). Two obstruction theorems then show that the method
does not terminate when \(a\ge2\operatorname{rad}M\), and that it does
not terminate when a certain two-letter carry word admits unit residues
modulo \(M\). The second obstruction gives the necessary condition
\(a\le2R\,J(N_0)\), where \(R\) is the product of the primes of \(M\)
not dividing \(2ab\), \(N_0\) is the product of the odd primes dividing
\(ab\), and \(J\) is the Jacobsthal function; an elementary bound
\(J(N)\le3^{\omega(N)}\) gives \(R\ge a^{1-o(1)}\). These obstruction
theorems are statements about the non-termination of a fixed procedure.
They are not counterexamples to any arithmetic statement, and they do not
assert that any base fails to have infinitely many composite terms.

% Markdown AST block 20: Para
\textbf{Verification.} Every adopted finite computation was carried out by two
implementations whose scientific cores (edge enumeration, component
decomposition, certification, contraction, lifting) were written
separately from the mathematical specification; their complete graph
data were compared element by element. Section 9 describes their
provenance and the comparison. Lean proofs establish both parts of
Theorem 1 and Corollary 2. The one-step proof for \(7/5\) uses a
kernel-checked finite certificate; the word-labelled proofs use
kernel-checked soundness theorems and \nolinkurl{native_decide} finite
evaluations. Section 9.4 gives the precise formalization coverage and
the additional compiler and native-evaluation trust of the latter
proofs.

% Markdown AST block 21: Para
Section 2 proves the nonperiodicity lemma. Section 3 presents the one-step
outer graphs and finite certification, and Section 4 the word-labelled
certificates. Section 5 gives the finite calculations for (1) and (2).
Section 6 proves Theorem 3. Sections 7 and 8 contain the structural and
obstruction results. Section 9 describes the verification and the exact
scope of the formalization. Section 10 states further results whose proofs are given in the
appendices, and Section 11 lists open problems.

% BEGIN PUBLIC ADDITION figure:certificate-flow
The two finite certification procedures share the same tail-preservation argument, summarized below.

\begin{figure}[htbp]
\centering
\begin{tikzpicture}[paperfig]
\node[fbox,text width=11.5cm] (orbit) at (0,0)
 {\figtext{Assume a true orbit eventually avoids the chosen primes}{真の軌道が選んだ素数を最終的にすべて回避すると仮定}\\
 $Q_n=\lfloor\xi(a/b)^n\rfloor,\qquad c_n=bQ_{n+1}-aQ_n$};
\node[fbox,text width=10cm] (graph) at (0,-1.45)
 {\figtext{Finite outer graph: every such orbit gives a walk}{有限外側グラフ：そのような軌道はすべてウォークを与える}\\
 \figtext{congruence + cell compatibility + carry alphabet}{合同条件＋セル整合条件＋キャリー集合}};
\draw[fedge] (orbit)--(graph);
\node[fbox,text width=10cm] (remove) at (0,-3)
 {\figtext{Decompose into SCCs; test all internal letters}{SCC分解し、内部辺の全文字について位相判定}\\
 \figtext{Remove certified cyclic SCCs and acyclic parts}{認証された巡回SCCと非巡回部分を除去}};
\draw[fedge] (graph)--(remove);
\node[fbox,text width=5cm] (ref) at (-3.2,-5.2)
 {\textbf{\figtext{One-step method}{一段階法}}\\[4pt]
 $K\longmapsto fK$\\\figtext{Refine retained cells; rebuild edges}{残存セルを細分し、辺を再構成}};
\node[fbox,text width=5cm] (word) at (3.2,-5.2)
 {\textbf{\figtext{Word method}{語ラベル法}}\\[4pt]
 $K\ \text{\figtext{fixed}{は固定}}$\\\figtext{Contract paths; lift by the next prime}{経路を縮約し、次の素数で持ち上げ}};
\draw[fedge] (remove.south) -- ++(0,-.32) -| (ref.north);
\draw[fedge] (remove.south) -- ++(0,-.32) -| (word.north);
\draw[fedge,dashed] (ref.west) -- ++(-.5,0) |- (remove.west);
\draw[fedge,dashed] (word.east) -- ++(.5,0) |- (remove.east);
\node[fbox,text width=11.5cm] (empty) at (0,-7.2)
 {\figtext{If the retained graph becomes empty: contradiction}{残存グラフが空になれば矛盾}\\
 \figtext{A true prime-avoiding tail survives every prescribed operation}{真の素数回避軌道の尾部は各操作を生き残る}};
\draw[fedge] (ref.south)-- ++(0,-.35) -| ([xshift=-3.2cm]empty.north);
\draw[fedge] (word.south)-- ++(0,-.35) -| ([xshift=3.2cm]empty.north);
\end{tikzpicture}

\caption{Logical flow of the one-step and word methods. Arrows from an orbit to a graph express inclusion in an outer approximation. The two branches use different operations: cell refinement in the one-step method, and contraction followed by prime lifting at fixed cell count in the word method. Dashed arrows indicate iteration.}
\label{fig:certificate-flow}
\end{figure}
% END PUBLIC ADDITION figure:certificate-flow

% Markdown AST block 22: 2. Carries and nonperiodicity
\Needspace{6\baselineskip}
\section{Carries and nonperiodicity}\label{sec:2}

% Markdown AST block 23: Para
Throughout the general argument, let
\[
a,b\in\mathbb N,\qquad a>b\ge2,\qquad \gcd(a,b)=1,\qquad \xi>0,
\]
and put \(r=a/b\), with division in \(\mathbb R\). For \(n\ge0\), define
\[
x_n=\xi r^n,\qquad Q_n=\lfloor x_n\rfloor\in\mathbb Z,
\qquad \theta_n=x_n-Q_n\in[0,1),
\tag{4}\label{eq:4}
\]
and the signed integer carry
\[
c_n=bQ_{n+1}-aQ_n=a\theta_n-b\theta_{n+1}.
\tag{5}\label{eq:5}
\]
Thus
\[
1-b\le c_n\le a-1,\qquad bQ_{n+1}=aQ_n+c_n.
\tag{6}\label{eq:6}
\]
The strict real bounds \(-b<c_n<a\) imply the integer bounds in (6).
An alternative indexing writes \(e_{n+1}=c_n\); no value of \(e_0\) is
needed. Since \(x_{n+1}=rx_n>x_n\), the sequence \((Q_n)\) is
nondecreasing, and \(Q_n\ge Q_0\) for all \(n\).

% Markdown AST block 24: Para
We shall also use \(Q_n\to+\infty\) in the explicit sense
\[
\forall T\in\mathbb Z\ \exists N\in\mathbb N\ \forall n\ge N,
\qquad Q_n>T.
\tag{7}\label{eq:7}
\]
Indeed, writing \(r=1+\delta\) with \(\delta>0\), Bernoulli's inequality
gives \(r^n\ge1+n\delta\), and
\(Q_n>\xi(1+n\delta)-1\).

% Markdown AST block 25: Lemma 4 (nonperiodicity).
\begin{samepage}
\begin{lemma}[nonperiodicity]\leavevmode\label{thm:4}
The carry sequence \((c_n)\) is not
eventually periodic: there do not exist \(N\ge0\) and \(L\ge1\) such that
\[
c_{n+L}=c_n
\]
for all \(n\ge N\).
\label{thm-end:4}
\end{lemma}
\end{samepage}

% Markdown AST block 26: Proof.
\begin{proof}\leavevmode\label{proof:4}
First suppose an integer sequence \((D_k)_{k\ge0}\) satisfies
\(bD_{k+1}=aD_k\).
Induction gives
\[
b^kD_k=a^kD_0.
\]
Since \(a^k\) and \(b^k\) are coprime, \(b^k\mid D_0\) for every
\(k\). As \(b\ge2\), this forces \(D_0=0\), without any assumption on
its sign.

% Markdown AST block 27: Para
Now assume that \(c\) has period \(L\) from \(N\) onward. For a fixed
\(m\ge N\), set
\[
D_k=Q_{m+k+L}-Q_{m+k}.
\]
Subtracting the two recurrences in (6) and using periodicity gives
\(bD_{k+1}=aD_k\).
The preceding observation yields
\(Q_{m+L}=Q_m\).
This holds for every \(m\ge N\), so in particular
\(Q_{N+kL}=Q_N\)
for every \(k\), contrary to (7). \qedhere
\end{proof}

% Markdown AST block 28: Para
This proof uses only an integral recurrence and growth. The hypotheses
\(b\ge2\) and \(\gcd(a,b)=1\) are essential: integer ratios admit
identically zero carries for suitable coefficients. The earlier
nonperiodicity lemmas cited above also exclude eventual periodicity, since
a tail can be absorbed into the new positive coefficient \(\xi r^N\).

% Markdown AST block 29: Para
We also record the iterated form of (6): with \(C_0=0\) and
\[
C_{i+1}=aC_i+b^ic_{u+i},
\]
\[
b^iQ_{u+i}=a^iQ_u+C_i\qquad(i\ge0),
\tag{8}\label{eq:8}
\]
by induction, since
\[
b^{i+1}Q_{u+i+1}=b^i(aQ_{u+i}+c_{u+i})=a(a^iQ_u+C_i)+b^ic_{u+i}.
\]

% Markdown AST block 30: 3. One-step outer graphs and finite certification
\Needspace{6\baselineskip}
\section{One-step outer graphs and finite certification}\label{sec:3}

% Markdown AST block 31: 3.1 The graphs
\subsection{The graphs}\label{sec:3.1}

% Markdown AST block 32: Para
Fix integers \(M\ge1\) and \(K\ge1\). The states recording a residue
coprime to \(M\) and a fractional cell are
\[
V(M,K)=\{(q,j)\in\mathbb Z^2:0\le q<M,\ \gcd(q,M)=1,\ 0\le j<K\}.
\tag{9}\label{eq:9}
\]
For \(M=1\) the residue coordinate is \(q=0\). For
\(W\subseteq V(M,K)\), define a directed labelled graph \(G_K(W)\)
on \(W\). There is an edge \((q,j)\xrightarrow{e}(z,k)\) exactly when
the endpoints lie in \(W\) and the following signed integer conditions hold:
\[
\begin{aligned}
1-b&\le e\le a-1, &&\text{(E1)}\\
bz&\equiv aq+e\pmod M, &&\text{(E2)}\\
aj-b(k+1)&<eK<a(j+1)-bk. &&\text{(E3)}
\end{aligned}
\tag{10}\label{eq:10}
\]
Edges with the same endpoints and different labels are distinct. We
write \(G(M,K)=G_K(V(M,K))\) for the \emph{complete} one-step graph.

% Markdown AST block 33: Para
Define the state of an actual orbit by
\[
\sigma_K(n)=\bigl(Q_n\bmod M,\ \lfloor K\theta_n\rfloor\bigr),
\tag{11}\label{eq:11}
\]
where the residue is the representative in \(\{0,\ldots,M-1\}\).

% Markdown AST block 34: Proposition 5 (orbit inclusion).
\begin{samepage}
\begin{proposition}[orbit inclusion]\leavevmode\label{thm:5}
If
\(\sigma_K(n),\sigma_K(n+1)\in W\), then
\[
\sigma_K(n)\xrightarrow{c_n}\sigma_K(n+1)
\]
is an edge of \(G_K(W)\).
\label{thm-end:5}
\end{proposition}
\end{samepage}

% Markdown AST block 35: Proof.
\begin{proof}\leavevmode\label{proof:5}
Condition (E1) is (6), and reduction of (5) modulo \(M\)
gives (E2). Write \(j=\lfloor K\theta_n\rfloor\) and
\(k=\lfloor K\theta_{n+1}\rfloor\). Then
\[
j\le K\theta_n<j+1,\qquad k\le K\theta_{n+1}<k+1,
\]
so
\[
aj-b(k+1)<aK\theta_n-bK\theta_{n+1}=Kc_n<a(j+1)-bk.
\]
This is (E3). Both inequalities are strict, even though a fractional part
may equal zero. \qedhere
\end{proof}

% Markdown AST block 36: Para
The proof uses only the real identity (5) and the half-open cells; it
does not use the integrality of \(Q_n\). This observation is used in
Section 7. The graph is an outer approximation. Each orbit that
eventually avoids all prime divisors of \(M\) gives a walk in the graph,
but no converse realizability assertion about graph walks is used. For
real numbers \(u,v\) with \(j\le u<j+1\), \(k\le v<k+1\), and
\(au-bv=eK\), condition (E3) is necessary and sufficient for the
existence of such a pair; it expresses the compatibility of two adjacent
cells only.

% Markdown AST block 37: Para
The conditions allow complete integer enumeration. For a fixed \(j,e\),
(E3) and \(0\le k<K\) are equivalent to
\[
\max\!\left(0,\left\lfloor\frac{aj-eK}{b}\right\rfloor\right)
\le k\le
\min\!\left(K-1,\left\lfloor\frac{a(j+1)-eK-1}{b}\right\rfloor\right).
\tag{12}\label{eq:12}
\]
For example, if \(A=aj-eK\), then \(A<b(k+1)\) is equivalent to
\(\lfloor A/b\rfloor\le k\). The other bound follows by replacing the
strict integer inequality \(bk<B\) with \(bk\le B-1\).
All floors in (12) are mathematical floors, including for negative
numerators. Truncation toward zero would give an incorrect enumeration.
An empty integer interval contributes no edges.

% Markdown AST block 38: Para
No invertibility of \(b\) modulo \(M\) is assumed. To see explicitly how
all solutions of (E2) can be found, put \(g=\gcd(b,M)\) and \(t=aq+e\).
There are no solutions unless \(g\mid t\). If this divisibility holds,
solve the reduced congruence
\[
(b/g)z_0\equiv t/g\pmod{M/g},\qquad 0\le z_0<M/g.
\]
For \(M/g=1\), use \(z_0=0\). Otherwise the coefficient \(b/g\) is
invertible modulo \(M/g\). The full set of representatives is
\[
z=z_0+\ell(M/g),\qquad 0\le\ell<g.
\tag{13}\label{eq:13}
\]
One then retains exactly those satisfying \(\gcd(z,M)=1\) and
\((z,k)\in W\). Equivalently, one can enumerate all unit residues \(z\)
and group them by \(bz\bmod M\), as the one-step implementations do.
In the case \(M=4290\), \(\gcd(5,4290)=5\), so division by \(5\) modulo
\(4290\) is unavailable.

% BEGIN PUBLIC ADDITION figure:cell-intersection
A single compatible cell pair illustrates why both inequalities in (E3) are strict.

\begin{figure}[htbp]
\centering
\begin{tikzpicture}[paperfig]
\begin{scope}[x=2.3cm,y=2.3cm]
\draw[fedge] (.65,.65)--(2.35,.65);
\node at (1.65,.27) {$u=K\theta_n$};
\draw[fedge] (.65,.65)--(.65,2.35) node[above] {$v=K\theta_{n+1}$};
\fill[black!5] (1,1) rectangle (2,2);
\draw (1,2)--(1,1)--(2,1);
\draw[dashed] (1,2)--(2,2)--(2,1);
\draw[femph] (1.2,1)--(1.6,2);
\fill[figblue] (1.2,1) circle(1.4pt);
\draw[figblue,fill=white,line width=.8pt] (1.6,2) circle(1.7pt);
\node[flabel,rotate=63.435] at (1.4,1.52) {$5u-2v=4$};
\foreach \x in {1,2}{\draw (\x,.62)--(\x,.68);\node[below] at(\x,.62){$\x$};}
\foreach \y in {1,2}{\draw (.62,\y)--(.68,\y);\node[left] at(.62,\y){$\y$};}
\end{scope}
\node[align=left,anchor=west,text width=6.5cm] at (6,4.3)
 {$a=5,\ b=2,\ K=4$\\[4pt]
 $j=k=1,\quad e=1$\\[8pt]
 $[j,j+1)\times[k,k+1)$\\[4pt]
 \figtext{Solid: included lower boundaries}{実線：含まれる下側境界}\\
 \figtext{Dashed: excluded upper boundaries}{破線：含まれない上側境界}};
\node[fbox,align=center,text width=6.3cm] at (9.2,1.8)
 {$aj-b(k+1)<eK<a(j+1)-bk$\\[5pt]
 $1<4<8$\\[4pt]
 \figtext{Both endpoint bounds are strict}{両端の不等号は厳密}};
\end{tikzpicture}

\caption{Cell compatibility for \(a=5,b=2,K=4,j=k=1,e=1\). The line meets the half-open square along \(u=(4+2v)/5\), \(1\le v<2\). The lower endpoint is included and the upper endpoint is excluded. The cell-pair test is local; it does not assert the existence of a full orbit.}
\label{fig:cell-intersection}
\end{figure}
% END PUBLIC ADDITION figure:cell-intersection

% Markdown AST block 39: 3.2 Components with periodic output
\Needspace{6\baselineskip}
\subsection{Components with periodic output}\label{sec:3.2}

% Markdown AST block 40: Para
A strongly connected component (SCC) is an equivalence class for mutual
directed reachability, allowing paths of length zero. We call an SCC
\emph{cyclic} if it contains an internal edge. Thus a singleton is cyclic
exactly when it has a self-loop.

% Markdown AST block 41: Lemma 6.
\begin{samepage}
\begin{lemma}\leavevmode\label{thm:6}
Every infinite walk in a finite directed graph eventually
remains in one cyclic SCC.
\label{thm-end:6}
\end{lemma}
\end{samepage}

% Markdown AST block 42: Proof.
\begin{proof}\leavevmode\label{proof:6}
After leaving one SCC for a different one, a walk cannot return:
a return would make the two components mutually reachable. There can
therefore be only finitely many changes of component. The final component
contains the infinite tail and hence an internal edge. \qedhere
\end{proof}

% Markdown AST block 43: Para
Let \(C\) be a cyclic SCC of \(G_K(W)\). Choose a root \(\rho\in C\),
and let \(h(v)\) be the shortest length of a path from \(\rho\) to \(v\)
inside \(C\). Define
\[
d=\gcd\{\,|h(u)+1-h(v)|:
u\xrightarrow{e}v\text{ is an internal edge of }C\,\}.
\tag{14}\label{eq:14}
\]
There is a nonempty closed walk in \(C\). Summing the differences
\(h(u)+1-h(v)\) along a closed walk of length \(\ell>0\) gives
\(\ell\). Consequently, the differences cannot all vanish, so \(d>0\).
Every internal edge satisfies
\[
h(v)\equiv h(u)+1\pmod d.
\tag{15}\label{eq:15}
\]
The same summation shows that the length of every closed walk is divisible
by \(d\). Along any nonempty closed walk, the first \(d\) departure
phases therefore realize all residues modulo \(d\). In particular, each
phase has at least one outgoing internal edge.

% Markdown AST block 44: Para
For each phase \(s\in\{0,\ldots,d-1\}\), collect the labels on \emph{all}
internal edges whose source \(u\) satisfies \(h(u)\equiv s\pmod d\).
We certify \(C\) if each such collection is a singleton, written
\(\{\lambda_s\}\).

% Markdown AST block 45: Lemma 7 (phase criterion).
\begin{samepage}
\begin{lemma}[phase criterion]\leavevmode\label{thm:7}
Every infinite walk confined to a certified
component has a periodic label sequence, with period \(d\).
\label{thm-end:7}
\end{lemma}
\end{samepage}

% Markdown AST block 46: Proof.
\begin{proof}\leavevmode\label{proof:7}
For a walk
\(v_0\xrightarrow{\ell_0}v_1\xrightarrow{\ell_1}\cdots\), (15) gives
\[
h(v_i)\equiv h(v_0)+i\pmod d.
\]
Therefore
\[
\ell_i=\lambda_{(h(v_0)+i)\bmod d},
\qquad \ell_{i+d}=\ell_i.\qedhere
\]
\end{proof}

% Markdown AST block 47: Para
The test includes parallel edges and every internal edge, not just tree
edges or a selected cycle. Edges leaving \(C\) are irrelevant once a
walk stays in \(C\). Lemma 12 below proves soundness for any
integer-valued \(h\). For a complete test, \(h\) is chosen as the
length of a path from a common root, as above; shortest paths are
unnecessary. Theorem 21 and Corollary 22 show that the outcome is
independent of the root and of these paths.

% Markdown AST block 48: Para
The integer \(d\) in (14) is kept separate from the primitive period
\(p\) of the phase-label word. For reporting, we replace
\((\lambda_0,\ldots,\lambda_{d-1})\) by its shortest repeating word and
then by its lexicographically least cyclic rotation. This word has length
\(p\mid d\), possibly \(p<d\). Such normalization is for comparison
and display; Lemma 7 requires neither minimality of the period nor a
particular starting phase.

% Markdown AST block 49: 3.3 Tail preservation and finite certification
\Needspace{6\baselineskip}
\subsection{Tail preservation and finite certification}\label{sec:3.3}

% Markdown AST block 50: Para
Let \(R_K(W)\) be the union of all cyclic SCCs of \(G_K(W)\) that fail
the phase criterion. The other vertices are removed.

% Markdown AST block 51: Lemma 8 (retained tail).
\begin{samepage}
\begin{lemma}[retained tail]\leavevmode\label{thm:8}
If \(\sigma_K(n)\in W\) for all sufficiently
large \(n\), then \(\sigma_K(n)\in R_K(W)\) for all sufficiently large
\(n\).
\label{thm-end:8}
\end{lemma}
\end{samepage}

% Markdown AST block 52: Proof.
\begin{proof}\leavevmode\label{proof:8}
Proposition 5 turns the given tail into an infinite walk with
labels \(c_n\). By Lemma 6 it eventually stays in one cyclic SCC. If
that component were certified, Lemma 7 would make \(c_n\) eventually
periodic, contrary to Lemma 4. It must therefore be retained. \qedhere
\end{proof}

% Markdown AST block 53: Para
For an integer \(f\ge2\), define the refinement
\[
\operatorname{Refine}_f(U)
=\{(q,fj+i):(q,j)\in U,\ 0\le i<f\}.
\tag{16}\label{eq:16}
\]

% Markdown AST block 54: Lemma 9 (refinement).
\begin{samepage}
\begin{lemma}[refinement]\leavevmode\label{thm:9}
If \(\sigma_K(n)\in U\), then
\[
\sigma_{fK}(n)\in\operatorname{Refine}_f(U).
\]
\label{thm-end:9}
\end{lemma}
\end{samepage}

% Markdown AST block 55: Proof.
\begin{proof}\leavevmode\label{proof:9}
From \(j=\lfloor K\theta_n\rfloor\) one obtains
\[
fj\le fK\theta_n<f(j+1).
\]
Thus
\[
\lfloor fK\theta_n\rfloor=fj+i
\]
for a unique integer
\(0\le i<f\). The residue coordinate is unchanged. \qedhere
\end{proof}

% Markdown AST block 56: Para
Starting with \(K_0\ge1\) and \(W_0=V(M,K_0)\), compute
\[
K_t=K_0f^t,\qquad U_t=R_{K_t}(W_t),\qquad
W_{t+1}=\operatorname{Refine}_f(U_t).
\tag{17}\label{eq:17}
\]
At each new stage the edges are recomputed from (10).

% Markdown AST block 57: Theorem 10 (finite certification).
\begin{samepage}
\begin{theorem}[finite certification]\leavevmode\label{thm:10}
If \(U_s=\varnothing\) for some
finite \(s\), then for every \(\xi>0\) and every \(N\in\mathbb N\)
there is an \(n\ge\max(N,1)\) such that
\[
\gcd\!\left(\left\lfloor\xi(a/b)^n\right\rfloor,M\right)>1.
\tag{18}\label{eq:18}
\]
\label{thm-end:10}
\end{theorem}
\end{samepage}

% Markdown AST block 58: Proof.
\begin{proof}\leavevmode\label{proof:10}
Suppose instead that for some \(\xi>0\), the integers \(Q_n\)
are coprime to \(M\) for all sufficiently large \(n\). By (9) and (11),
the corresponding states lie in \(W_0\). Lemma 8 places a further tail
in \(U_0\), and Lemma 9 places the same tail at the finer scale in
\(W_1\). Iterating these two implications over the finitely many stages
puts a tail in \(U_s\), which is empty. This contradiction proves (18).
\qedhere
\end{proof}

% Markdown AST block 59: Para
The beginning of the tail may increase at each stage and may depend on
\(\xi\). Neither a uniform starting time nor preservation of every finite
prefix is needed. In particular, refinement is proved for true orbit
states; it does not presume that every coarse graph walk has a finer lift.
The graphs themselves have no \(\xi\) input. The passage to all positive
real coefficients is supplied by these lemmas, rather than by sampling
initial values or computing long numerical trajectories.

% Markdown AST block 60: Para
Theorem 10 assumes a finite successful calculation. It does not assert
termination for arbitrary choices of \(a,b,M,K_0,f\).

% Markdown AST block 61: Para
The modulus may also grow from stage to stage. For \(U\subseteq V(M,K)\)
and integers \(h,f\ge1\), put \(M'=hM\), \(K'=fK\), and
\[\begin{aligned}
&U'=\{(q+\ell M,\ fj+i):(q,j)\in U,\ 0\le\ell<h,\ 0\le i<f,\\\
&\gcd(q+\ell M,hM)=1\}\subseteq V(M',K').
\end{aligned}\]

% Markdown AST block 62: Lemma 11 (joint enlargement).
\begin{samepage}
\begin{lemma}[joint enlargement]\leavevmode\label{thm:11}
If \(\sigma_K(n)\in U\) as a state
modulo \(M\) and \(\gcd(Q_n,hM)=1\), then
\[
\sigma_{K'}(n)\in U'
\]
as a
state modulo \(M'\).
\label{thm-end:11}
\end{lemma}
\end{samepage}

% Markdown AST block 63: Proof.
\begin{proof}\leavevmode\label{proof:11}
Let \(q=Q_n\bmod M\) and \(q'=Q_n\bmod hM\). Then
\(q'\equiv q\pmod M\) and \(0\le q'<hM\), so \(q'=q+\ell M\) with
\(0\le\ell<h\), and
\[
\gcd(q',hM)=\gcd(Q_n,hM)=1.
\]
The cell coordinate
is handled by Lemma 9. \qedhere
\end{proof}

% Markdown AST block 64: Para
Since every intermediate modulus divides the final one, a tail coprime to
the final modulus is coprime to each intermediate modulus, and the
induction of Theorem 10 goes through with Lemma 11 in place of Lemma 9.

% BEGIN PUBLIC ADDITION figure:cell-refinement
Cell refinement preserves the cell containing each true orbit point, as the following example shows.

\begin{figure}[htbp]
\centering
\begin{tikzpicture}[paperfig]
\node[anchor=east] at(-.2,1.45){$K=4,\ j=1$};
\draw[line width=1pt] (0,1.45)--(9,1.45);
\foreach \x/\l in {0/{1/4},9/{1/2}} {\node[above=3pt] at(\x,1.45){$\l$};}
\fill (0,1.45) circle(2pt);\draw[fill=white] (9,1.45) circle(2pt);
\node[anchor=east] at(-.2,0){$fK=12$};
\foreach \x/\j in {0/3,3/4,6/5}{
 \draw[line width=1pt] (\x,0)--++(3,0);
 \fill (\x,0) circle(2pt);\draw[fill=white] (\x+3,0) circle(2pt);
 \node[below=3pt] at(\x+1.5,0){$\j$};
}
\node[above=3pt] at(3,0){$1/3$};\node[above=3pt] at(6,0){$5/12$};
\draw[fedge,dashed] (4.5,1.2)--(4.5,.3) node[midway,left] {$f=3$};
\draw[femph,dotted] (5.4,1.45)--(5.4,0);
\fill[figblue] (5.4,1.45) circle(2pt);\fill[figblue] (5.4,0) circle(2pt);
\node[above=3pt,align=center] at(5.4,1.5){$\theta=2/5$};
\node[fbox,text width=11.5cm] at (3.8,-1.25)
 {$\lfloor4\theta\rfloor=1,\qquad\lfloor12\theta\rfloor=4=3\cdot1+1$\\[5pt]
 $\sigma_K(n)\in U\quad\Longrightarrow\quad
 \sigma_{fK}(n)\in\operatorname{Refine}_f(U)$};
\node[align=center,text width=12cm] at(3.8,-2.5)
 {\figtext{The same true orbit point lies in exactly one finer cell.}{同じ真の軌道点は、細分後のちょうど一つのセルに入る。}\\
 \figtext{Edges are rebuilt; an arbitrary coarse walk need not lift.}{辺は再構成する。任意の粗いウォークが持ち上がるとは限らない。}};
\end{tikzpicture}

\caption{The cell \([1/4,1/2)\) at \(K=4\) splits into cells \(3,4,5\) at \(fK=12\). The point \(\theta=2/5\) belongs to cell \(4\). The residue coordinate is unchanged. This is an inclusion for true orbit states, not a lifting assertion for arbitrary graph walks.}
\label{fig:cell-refinement}
\end{figure}
% END PUBLIC ADDITION figure:cell-refinement

% Markdown AST block 65: 4. Word-labelled certificates
\Needspace{6\baselineskip}
\section{Word-labelled certificates}\label{sec:4}

% Markdown AST block 66: 4.1 Word-labelled graphs and coverings
\subsection{Word-labelled graphs and coverings}\label{sec:4.1}

% Markdown AST block 67: Para
Let \(\mathcal C\subseteq[1-b,a-1]\cap\mathbb Z\) be a finite set of
carries, the \emph{alphabet}. A \emph{word-labelled graph} \(G=(V,E)\) consists of
a finite vertex set \(V\) and a finite set \(E\) of edges
\(e=(u,w,v)\) with \(u,v\in V\) and \(w\in\mathcal C^{+}\) a nonempty
finite word; \(|w|\ge1\) is the length of the edge. Edges are identified
as triples: the same triple is one edge, and edges with the same
endpoints and different words are distinct. An infinite walk is
\[
v_0\xrightarrow{w_0}v_1\xrightarrow{w_1}\cdots
\]
with each
\((v_i,w_i,v_{i+1})\in E\); its \emph{output} is the concatenation
\[
w_0w_1w_2\cdots\in\mathcal C^{\mathbb N},
\]
an infinite word because no
edge has length zero. A one-step graph of Section 3 is the special case
in which every word has length one.

% Markdown AST block 68: Para
Let \((c_n)\) be the carry sequence of a true orbit and \(n_0\ge0\). We
say that \(G\) \emph{covers the tail from \(n_0\)}, written
\(G\models n_0\), if there are an infinite walk \((v_i,w_i)_{i\ge0}\) of
\(G\) and times \(t_0=n_0<t_1<t_2<\cdots\) with
\[
t_{i+1}=t_i+|w_i|,\qquad
w_i=(c_{t_i},c_{t_i+1},\ldots,c_{t_{i+1}-1})\qquad(i\ge0);
\]
that is, the output of the walk equals \((c_n)_{n\ge n_0}\). For each
operation \(\Phi\) below (removal, contraction, lifting) we prove \emph{tail
preservation}: if \(G\models n_0\) (together with an arithmetic
hypothesis specific to the operation), then \(\Phi(G)\models n_1\) for
some \(n_1\ge n_0\). The starting time may increase at each operation and
may depend on \(\xi\); no finite prefix is preserved.

% Markdown AST block 69: 4.2 The weighted phase criterion and removal
\Needspace{6\baselineskip}
\subsection{The weighted phase criterion and removal}\label{sec:4.2}

% Markdown AST block 70: Para
Let \(C\) be a cyclic SCC of \(G\), that is, an SCC containing at least
one internal edge \((u,w,v)\) with \(u,v\in C\). Let \(h:C\to\mathbb Z\)
be \emph{any} integer-valued function, and for an internal edge
\(e=(u,w,v)\) put
\[
\begin{aligned}
\delta(e)&=h(u)+|w|-h(v),\\
d&=d_C=\gcd\{|\delta(e)|:e\text{ an internal edge of }C\}.
\end{aligned}
\]

% Markdown AST block 71: Lemma 12 (weighted phase criterion).
\begin{samepage}
\begin{lemma}[weighted phase criterion]\leavevmode\label{thm:12}
(a) \(d\ge1\), and every
internal edge satisfies
\[
h(v)\equiv h(u)+|w|\pmod d;
\]
the total word
length of every nonempty closed walk in \(C\) is a positive multiple of \(d\).

% Markdown AST block 72: OrderedList
\begin{enumerate}
\def\labelenumi{(\alph{enumi})}
\setcounter{enumi}{1}
\tightlist
\item
  Suppose there is a map \(\lambda:\mathbb Z/d\to\mathcal C\) such that
  for every internal edge \(e=(u,w,v)\) and every letter position
  \(0\le i<|w|\),
  \[
  w_i=\lambda\bigl((h(u)+i)\bmod d\bigr).
  \tag{19}\label{eq:19}
  \]
  Then every infinite walk \((v_i,w_i)\) inside \(C\) has output
  \(y=w_0w_1\cdots\) with
  \[
  y_t=\lambda((h(v_0)+t)\bmod d)
  \]
  for all
  \(t\ge0\); in particular \(y_{t+d}=y_t\) for all \(t\).
\end{enumerate}
\label{thm-end:12}
\end{lemma}
\end{samepage}

% Markdown AST block 73: Para
We call \(C\) \emph{certified} (with respect to \(h\)) when (19) holds for
some \(\lambda\).

% Markdown AST block 74: Proof.
\begin{proof}\leavevmode\label{proof:12}
(a) Take an internal edge \(e_1=(u,w,v)\) and, by strong
connectivity, a path \(e_2\cdots e_m\) from \(v\) to \(u\) inside \(C\)
(empty if \(u=v\)). Along the closed walk \(e_1\cdots e_m\) the terms of
\(h\) telescope, so
\[
\sum_{k\le m}\delta(e_k)=\sum_{k\le m}|w_k|\ge1.
\]
Hence not all \(\delta(e)\) vanish and \(d\ge1\). The congruence is the
definition of the gcd, and the total length of a nonempty closed walk equals the
sum of its \(\delta\), a positive multiple of \(d\).

% Markdown AST block 75: OrderedList
\begin{enumerate}
\def\labelenumi{(\alph{enumi})}
\setcounter{enumi}{1}
\tightlist
\item
  Put \(\tau_i=\sum_{k<i}|w_k|\). Induction along the walk with (a)
  gives
  \[
  h(v_i)\equiv h(v_0)+\tau_i\pmod d.
  \]
  The output letter at
  position \(t=\tau_i+k\) with \(0\le k<|w_i|\) is
  \[
  (w_i)_k=\lambda((h(v_i)+k)\bmod d)=\lambda((h(v_0)+t)\bmod d).\qedhere
\]
\end{enumerate}
\end{proof}

% Markdown AST block 76: Para
Lemma 12 is a sufficient condition for any integer-valued \(h\).
For the algorithms below, choose a root \(\rho\in C\) and, for each
\(v\in C\), a path \(\pi_v\) from \(\rho\) to \(v\) inside \(C\),
and set
\[
h(v)=|\operatorname{out}(\pi_v)|.
\]
Thus \(h(v)\) is the total word length, rather than the number of
edges. The path to \(\rho\) may be empty, and the other paths need not
be shortest. The implementations use the first root paths found by
their graph traversals.

% Markdown AST block 77: Para
Unless an arbitrary \(h\) is explicitly specified, \emph{certified} and
\emph{uncertified} below refer to the full test (19) with these root path
lengths. With this convention the test is complete for eventual
periodicity of all internal infinite outputs, and its outcome is
independent of the root and paths (Theorem 21 and Corollary 22).
A test left unfinished because of a resource limit is inconclusive.

% Markdown AST block 78: Para
By Lemma 12(a), every phase \(s\in\mathbb Z/d\) carries at
least one letter of some internal edge, since a closed walk of total
length \(\ell\ge d\) occupies the phases
\(h(v_0),\ldots,h(v_0)+\ell-1\). With all words of length one, (19)
is the phase criterion of Section 3.2.

% Markdown AST block 79: Para
\textbf{Removal.} Let \(\mathcal U\) be the family of uncertified cyclic SCCs
of \(G\). Define
\[
R(G)=\Bigl(\bigcup_{C\in\mathcal U}C,\
\{(u,w,v)\in E:\exists C\in\mathcal U,\ u,v\in C\}\Bigr):
\]
only the internal edges of uncertified components are kept; edges between
components, acyclic components, and certified components are dropped.

% Markdown AST block 80: Proposition 13 (tail preservation under removal).
\begin{samepage}
\begin{proposition}[tail preservation under removal]\leavevmode\label{thm:13}
If
\(G\models n_0\), then \(R(G)\models n_1\) for some \(n_1\ge n_0\).
\label{thm-end:13}
\end{proposition}
\end{samepage}

% Markdown AST block 81: Proof.
\begin{proof}\leavevmode\label{proof:13}
The covering walk is an infinite walk in a finite graph, so by
Lemma 6 there are \(i_1\) and a cyclic SCC \(C\) with \(v_i\in C\) for
all \(i\ge i_1\); the edges \((v_i,w_i,v_{i+1})\), \(i\ge i_1\), are
internal to \(C\). If \(C\) were certified, Lemma 12(b) would give
\[
c_{n+d}=c_n
\]
for all \(n\ge t_{i_1}\), with \(d\ge1\), contrary to
Lemma 4. Hence \(C\in\mathcal U\), and \((v_i,w_i)_{i\ge i_1}\) is a
walk of \(R(G)\); take \(n_1=t_{i_1}\). \qedhere
\end{proof}

% Markdown AST block 82: Para
Keeping only internal edges is sound because Lemma 6 makes every edge of
the tail internal to one component. Dropping an internal edge of an
uncertified component would not be sound; keeping more edges would be
sound but would change the reported counts.

% BEGIN PUBLIC ADDITION figure:phase-test
The phase test can be read as a consistency check on letters assigned to cyclic positions.

\begin{figure}[htbp]
\centering
\begin{tikzpicture}[paperfig]
\foreach \x/\s in {0/L,7/R}{
 \begin{scope}[xshift=\x cm]
 \node[fnode] (\s u) at (0,0) {$u$};
 \node[fnode] (\s v) at (3.3,0) {$v$};
 \node[below=3pt of \s u] {$h(u)=0$};
 \node[below=3pt of \s v] {$h(v)=1$};
 \draw[fedge] (\s u) to[bend left=30] node[flabel,above] {$2$} (\s v);
 \draw[fedge] (\s v) to[bend left=30] node[flabel,below] {$4$} (\s u);
 \node at (1.65,-1.35) {$\delta(u\to v)=0,\quad\delta(v\to u)=2$};
 \node at (1.65,-1.85) {$d=2$};
 \end{scope}
}
\node[ftitle] at(1.65,1.3){\figtext{Certified cycle}{認証される閉路}};
\node[ftitle] at(8.65,1.3){\figtext{One extra edge}{辺を一本追加}};
\draw[fedge,dashed] (Ru) to node[flabel,above] {$4$} (Rv);
\node[fbox,text width=5.2cm] at (1.65,-2.85)
 {$\lambda(0)=2,\quad\lambda(1)=4$\\[3pt]
 \figtext{Output: $242424\ldots$ (up to shift)}{出力：$242424\ldots$（開始位相を除く）}};
\node[fbox,text width=5.2cm] at (8.65,-2.85)
 {$\lambda(0)=2\ \text{\figtext{and}{かつ}}\ \lambda(0)=4$\\[3pt]
 \figtext{Conflict at phase $0$; not certified}{位相 $0$ で不整合：認証されない}};
\node[align=center,text width=12cm] at (5.15,-4)
 {\figtext{In both panels, $h$ is the total word length of a chosen path from the root $u$.}{両図とも、$h$ は根 $u$ から選んだ経路の総語長である。}};
\end{tikzpicture}

\caption{Two abstract word graphs with one-letter labels. In both panels, \(h(u)=0,h(v)=1\) are total word lengths of paths from root \(u\), and \(d=2\). The left cycle is certified. Adding the dashed edge on the right creates two different letters at phase zero, so the full test fails. These examples illustrate the criterion, not the particular carry alphabet for \(5/2\).}
\label{fig:phase-test}
\end{figure}
% END PUBLIC ADDITION figure:phase-test

% Markdown AST block 83: 4.3 Exact contraction of deterministic paths
\Needspace{6\baselineskip}
\subsection{Exact contraction of deterministic paths}\label{sec:4.3}

% Markdown AST block 84: Para
Let \(G'=R(G)\). Identical triples are merged first; this does not change
the set of walks. Out-degrees \(\deg^+(u)\) are counted in \(G'\) after
merging.

% Markdown AST block 85: Lemma 14.
\begin{samepage}
\begin{lemma}\leavevmode\label{thm:14}
(a) Every vertex of \(G'\) has \(\deg^+\ge1\). (b) If every
vertex of a component \(C\) of \(G'\) has \(\deg^+=1\), then \(C\) is a
simple directed cycle (a single vertex with a self-loop included), and
\(C\) is certified in the sense of Lemma 12 for the root path lengths
along the cycle.
\label{thm-end:14}
\end{lemma}
\end{samepage}

% Markdown AST block 86: Proof.
\begin{proof}\leavevmode\label{proof:14}
(a) If \(|C|\ge2\), strong connectivity gives every vertex an
internal out-edge; if \(|C|=1\), cyclicity gives a self-loop.

% Markdown AST block 87: Para
\textbf{(b)} With
unique successors the internal edges form a functional graph, and a
strongly connected functional graph is one cycle
\[
v_1\xrightarrow{w_1}\cdots\xrightarrow{w_m}v_1.
\]
Put \(L=\sum|w_k|\),
\(h(v_1)=0\),
\[
h(v_{k+1})=h(v_k)+|w_k|.
\]
Then \(\delta(e_k)=0\) for
\(k<m\), \(\delta(e_m)=L\), and \(d=L\); each phase \(0\le s<L\) carries
exactly one letter, the \(s\)-th letter of \(w_1\cdots w_m\), so (19)
holds. \qedhere
\end{proof}

% Markdown AST block 88: Para
Hence, when removal uses the full test with root path lengths, no
component consisting only of out-degree-one vertices reaches the
contraction step. The implementations
nevertheless keep one \emph{anchor} vertex in such a component as a safeguard
(used when a phase test is skipped because of a resource bound); in all
computations of Section 5 this rule never fired.

% Markdown AST block 89: Definition (macro vertices and contraction).
\begin{definition}[macro vertices and contraction]
Let \(B\) consist of
all vertices of \(G'\) with \(\deg^+\ge2\), together with one anchor (the
least vertex in the vertex order) in every component that contains no such
vertex. For \(u\in B\) and each out-edge \(e=(u,w,v)\) of \(u\): while
\(v\notin B\), follow the unique out-edge \((v,w',v')\) of \(v\) and
replace \(w\leftarrow ww'\), \(v\leftarrow v'\). When first \(v\in B\),
record the \emph{macro edge} \((u,w,v)\). The contraction \(P(G')\) has vertex
set \(B\) and edge set the set of macro edges (identical triples merged).
\end{definition}

% Markdown AST block 90: Lemma 15.
\begin{samepage}
\begin{lemma}\leavevmode\label{thm:15}
The iteration in the definition terminates, and every
infinite walk of \(G'\) visits \(B\) infinitely often.
\label{thm-end:15}
\end{lemma}
\end{samepage}

% Markdown AST block 91: Proof.
\begin{proof}\leavevmode\label{proof:15}
If the iteration did not terminate, the vertices visited would
all have out-degree one and would be finitely many, so the iteration would
enter a cycle \(Z\) of out-degree-one vertices. The unique out-edges of
\(Z\) stay in \(Z\), so \(Z\) has no edge to its complement in its
component \(C\); strong connectivity gives \(C=Z\), a component of
out-degree-one vertices, whose anchor lies in \(B\cap Z\), and the
iteration stops there: a contradiction. The same argument applies to an
infinite walk that avoids \(B\) from some point on. \qedhere
\end{proof}

% Markdown AST block 92: Proposition 16 (exactness of contraction).
\begin{samepage}
\begin{proposition}[exactness of contraction]\leavevmode\label{thm:16}
Every infinite walk of
\(G'\) starting at a vertex of \(B\) determines an infinite walk of
\(P(G')\) whose vertex sequence is the sequence of visits to \(B\) of
the former, whose output is the same infinite word, and whose \(k\)-th
macro edge has length equal to the total length of the corresponding
segment; conversely, every infinite walk of \(P(G')\) arises in this way
from at least one infinite walk of \(G'\). Consequently, if \(G'\models n_1\) then
\[
P(G')\models n_2,
\]
where \(n_2\) is the first time at or after
\(n_1\) at which the covering walk visits \(B\).
\label{thm-end:16}
\end{proposition}
\end{samepage}

% Markdown AST block 93: Proof.
\begin{proof}\leavevmode\label{proof:16}
At a vertex \(u\in B\) the walk chooses an edge
\(e=(u,w,v)\); while \(v\notin B\) it is forced through the unique
out-edge, following the same vertices and words as the iteration in the
definition, until it reaches the next vertex of \(B\) after finitely many
steps (Lemma 15).

% Markdown AST block 94: Para
Thus each segment between consecutive visits to \(B\)
is exactly one macro edge, and conversely each macro edge expands to at
least one path of \(G'\) (after identical triples are merged, two
segments with the same start, word, and end give the same macro edge, so
the correspondence is surjective but need not be injective).

% Markdown AST block 95: Para
The
correspondence is segment by segment, and outputs are concatenations of
the same words; only the forward direction is used for tail
preservation. The covering statement follows,
the time \(n_2\) existing by Lemma 15. \qedhere
\end{proof}

% Markdown AST block 96: Para
Contraction is used here only between removal and lifting; the cell count
\(K\) is fixed at the initial stage of the word method, and cell
coordinates enter only through the initial edges. We do not use, and do
not claim, any refinement of cells after a contraction.

% BEGIN PUBLIC ADDITION figure:word-contraction
Contraction changes the segmentation of the output word while preserving every letter and its order.

\begin{figure}[htbp]
\centering
\begin{tikzpicture}[paperfig]
\node[ftitle,anchor=west] at(-.4,1.7){\figtext{Before contraction}{縮約前}};
\node[fnode] (a) at(0,0){$u$};\node[fnode] (x) at(4,0){$x$};\node[fnode] (b) at(8,0){$v$};
\draw[fedge,femph] (a)--node[flabel,above] {$(-1)$}(x);
\draw[fedge,femph] (x)--node[flabel,above] {$(1,3)$}(b);
\draw[fedge,dashed] (a) to[bend right=23] node[flabel,below] {$(3)$}(b);
\draw[fedge] (b) to[bend right=23] node[flabel,above] {$(1)$}(a);
\draw[fedge] (b) edge[loop right] node[right] {$(-1)$}(b);
\node[below=4pt] at(x.south){$\deg^+(x)=1$};
\node[align=center] at(10.6,-.1){$u,v\in B$\\$x\notin B$};
\draw[fedge,line width=1pt] (4,-1.45)--(4,-2.1);
\node[ftitle,anchor=west] at(-.4,-2.5){\figtext{After contraction}{縮約後}};
\node[fnode] (aa) at(0,-3.9){$u$};\node[fnode] (bb) at(8,-3.9){$v$};
\draw[fedge,femph] (aa)--node[flabel,above] {$(-1,1,3)$}(bb);
\draw[fedge,dashed] (aa) to[bend right=23] node[flabel,below] {$(3)$}(bb);
\draw[fedge] (bb) to[bend right=23] node[flabel,above] {$(1)$}(aa);
\draw[fedge] (bb) edge[loop right] node[right] {$(-1)$}(bb);
\node[align=center] at(10.6,-3.9){$\ell=3$};
\node[align=center,text width=12cm] at(5,-5.6)
 {\figtext{Concatenate every letter, in order; keep the total elapsed length.}{すべての文字を順に連結し、経過した総語長を保持する。}};
\end{tikzpicture}

\caption{An exact contraction. The vertex \(x\) has one outgoing edge; \(u\) and \(v\) are macro vertices because each has two outgoing edges. The words \((-1)\) and \((1,3)\) concatenate to \((-1,1,3)\), whose length remains three. All other edges are retained. Line styles identify corresponding edges.}
\label{fig:word-contraction}
\end{figure}
% END PUBLIC ADDITION figure:word-contraction

% Markdown AST block 97: 4.4 Lifting by a new prime
\Needspace{6\baselineskip}
\subsection{Lifting by a new prime}\label{sec:4.4}

% Markdown AST block 98: Para
Let \(p\) be a prime with \(p\nmid b\). The \emph{lift} \(\operatorname{Lift}_p(G)\)
has vertices \((u,q)\) with \(u\in V(G)\) and \(q\in\mathbb F_p^\times\).
For an edge \((u,w,v)\in E(G)\) with \(w=(c_0,\ldots,c_{\ell-1})\) and
\(q_0=q\), compute
\[
q_{i+1}=b^{-1}(aq_i+c_i)\bmod p\qquad(0\le i<\ell).
\]
There is an edge \(((u,q_0),w,(v,q_\ell))\) if and only if
\(q_0,q_1,\ldots,q_\ell\) are all nonzero. The condition on the endpoint
residue \(q_\ell\) is included, since \(q_\ell\) is the starting residue
of the next edge. Vertices incident to no edge may be deleted: a vertex
without out-edges never lies on an infinite walk, and a vertex without
in-edges lies only at the start of one, so deleting them preserves
coverings after shifting the start by one edge.

% Markdown AST block 99: Proposition 17 (tail preservation under lifting).
\begin{samepage}
\begin{proposition}[tail preservation under lifting]\leavevmode\label{thm:17}
Let \(p\nmid b\),
\(G\models n_0\), and suppose \(p\nmid Q_n\) for all \(n\ge n'\). Then
\[
\operatorname{Lift}_p(G)\models n_1
\]
for \(n_1=t_{i_0}\), where
\(i_0\) is the least index with \(t_{i_0}\ge\max(n_0,n')\).
\label{thm-end:17}
\end{proposition}
\end{samepage}

% Markdown AST block 100: Proof.
\begin{proof}\leavevmode\label{proof:17}
For \(i\ge i_0\) put
\[
q^{(i)}=Q_{t_i}\bmod p\ne0.
\]
Reducing
\(bQ_{n+1}=aQ_n+c_n\) modulo \(p\) gives
\[
Q_{n+1}\equiv b^{-1}(aQ_n+c_n).
\]

% Markdown AST block 101: Para
Hence the forward iteration along
\(w_i=(c_{t_i},\ldots,c_{t_{i+1}-1})\)
starting at \(q^{(i)}\) produces
the residues of \(Q_{t_i},Q_{t_i+1},\ldots,Q_{t_{i+1}}\), all nonzero,
ending at \(q^{(i+1)}\).

% Markdown AST block 102: Para
Thus
\[
((v_i,q^{(i)}),w_i,(v_{i+1},q^{(i+1)}))
\]
is an edge of the lift, and
\((v_i,q^{(i)})_{i\ge i_0}\) is a covering walk. \qedhere
\end{proof}

% Markdown AST block 103: Para
Enumerating all starting residues \(q\in\mathbb F_p^\times\) is what makes
the argument independent of the size of the initial integer: whatever
unit \(Q_{t_i}\bmod p\) is, the corresponding vertex exists.

% Markdown AST block 104: Lemma 18 (root formula).
\begin{samepage}
\begin{lemma}[root formula]\leavevmode\label{thm:18}
Let \(p\nmid ab\), and let
\[
C_0=0,\qquad C_{i+1}=aC_i+b^ic_i
\]
for a word
\(w=(c_0,\ldots,c_{\ell-1})\), so that (8) holds. Put
\[
R_i=-C_ia^{-i}\bmod p,
\]
so that \(R_0=0\) and
\[
R_{i+1}=R_i-c_ib^ia^{-i-1}.
\]
For a starting residue \(q\), the forward
iteration has \(q_0,\ldots,q_\ell\) all nonzero if and only if
\[
q\notin\{R_0,\ldots,R_\ell\},
\]
and then
\[
\begin{aligned}
q_\ell&=\alpha q+\beta\qquad \text{with}\qquad \alpha=(ab^{-1})^\ell,\\
\beta_0&=0,\qquad \beta_{i+1}=ab^{-1}\beta_i+b^{-1}c_i.
\end{aligned}
\]
\label{thm-end:18}
\end{lemma}
\end{samepage}

% Markdown AST block 105: Proof.
\begin{proof}\leavevmode\label{proof:18}
By (8) applied to any integer lift of \(q\),
\[
b^iq_i\equiv a^iq+C_i.
\]
Since \(p\nmid b\), \(q_i\equiv0\) if and
only if \(a^iq+C_i\equiv0\), that is, \(q\equiv-C_ia^{-i}=R_i\)
(\(p\nmid a\)).

% Markdown AST block 106: Para
The recurrence for \(R_i\) is
\[
-(aC_i+b^ic_i)a^{-i-1}=R_i-c_ib^ia^{-i-1},
\]
and the affine form of
\(q_\ell\) is the expansion of
\[
q_{i+1}=ab^{-1}q_i+b^{-1}c_i.\qedhere
\]
\end{proof}

% Markdown AST block 107: Para
Thus enumeration by forbidden starting residues and forward iteration give
the same edge set; implementation A uses the former and implementation B
the latter (Section 9). If \(p\mid b\), no inverse of \(b\) exists and
\[
0\equiv aQ_n+c_n\pmod p,
\]
so \(Q_n\bmod p\) is determined by
\(c_n\) alone and the condition \(p\nmid Q_n\) is equivalent to
\(c_n\not\equiv0\pmod p\), a condition on the alphabet. If \(p\mid a\),
then
\[
Q_{n+1}\equiv b^{-1}c_n,
\]
independent of the starting residue;
provided every letter of the alphabet is coprime to \(p\) (as for the
alphabets of Lemmas 23 and 24, whose letters are coprime to \(ab\)),
the lift multiplies vertices by \(p-1\) without removing anything, while
an edge whose word contains a letter divisible by \(p\) would be removed
for every starting residue.
This is how the primes \(2\) and \(5\) are handled in Section 5.

% BEGIN PUBLIC ADDITION figure:prime-lift
Checking an entire word under a prime lift requires every intermediate residue, not just its endpoints.

\begin{figure}[htbp]
\centering
\begin{tikzpicture}[paperfig]
\node[align=center] at(5.2,1.5)
 {$a=5,\quad b=2,\quad p=7,\quad w=(-1,1)$\\[5pt]
 $q_{i+1}=4(5q_i+c_i)\pmod7$};
\node[ftitle,anchor=west] at(-1.3,0){\figtext{Keep}{保持}};
\node[fnode] (a) at(1,0){$1$};\node[fnode] (b) at(5,0){$2$};\node[fnode] (c) at(9,0){$2$};
\draw[fedge] (a)--node[flabel,above] {$c_0=-1$}(b);
\draw[fedge] (b)--node[flabel,above] {$c_1=1$}(c);
\node[below=5pt] at(5,-.3){$q_0,q_1,q_2\ne0$};
\node[ftitle,anchor=west] at(-1.3,-2){\figtext{Discard}{除去}};
\node[fnode] (d) at(1,-2){$3$};\node[fnode,fill=black!12,double] (e) at(5,-2){$0$};\node[fnode] (f) at(9,-2){$4$};
\draw[fedge,dashed] (d)--node[flabel,above] {$c_0=-1$}(e);
\draw[fedge,dashed] (e)--node[flabel,above] {$c_1=1$}(f);
\node[below=5pt] at(5,-2.3){$q_0,q_2\ne0,\qquad q_1=0$};
\node[fbox,text width=11.5cm] at(5,-3.85)
 {\figtext{Endpoint checks alone would retain the second edge incorrectly.}{端点だけの検査では、二つ目の辺を誤って保持してしまう。}\\
 \figtext{Inspect all intermediate residues, including the last one.}{最後の剰余を含む、すべての中間剰余を検査する。}};
\end{tikzpicture}

\caption{Two lifts of the same word \((-1,1)\) for \(a=5,b=2,p=7\), where \(2^{-1}=4\pmod7\). Starting at residue \(1\) gives \(1,2,2\), so the edge survives. Starting at \(3\) gives \(3,0,4\), so it is discarded even though both endpoints are nonzero. Every displayed residue follows from integer arithmetic modulo seven.}
\label{fig:prime-lift}
\end{figure}
% END PUBLIC ADDITION figure:prime-lift

% Markdown AST block 108: 4.5 The finite success theorem
\Needspace{6\baselineskip}
\subsection{The finite success theorem}\label{sec:4.5}

% Markdown AST block 109: Para
\textbf{Hypothesis (H).} A finite set \(S_0\) of primes with
\[
M_0=\prod_{p\in S_0}p
\]
and an alphabet \(\mathcal C\) satisfy: for
every \(\xi>0\) and every \(n\),
\[
\gcd(Q_n,M_0)=\gcd(Q_{n+1},M_0)=1
\]
implies \(c_n\in\mathcal C\).

% Markdown AST block 110: Theorem 19 (word-labelled finite success).
\begin{samepage}
\begin{theorem}[word-labelled finite success]\leavevmode\label{thm:19}
Let \(a>b\ge2\),
\(\gcd(a,b)=1\), \(K\ge1\), let \(S_0,\mathcal C\) satisfy (H), and let
\(p_1,\ldots,p_s\) be distinct primes with \(p_k\nmid ab\) and
\(p_k\notin S_0\). Let \(G_0^{\rm in}\) be the graph with vertices
\(\{0,\ldots,K-1\}\) and edges \((j,(c),k)\) for all \(c\in\mathcal C\)
and \(0\le j,k<K\) with
\[
aj-b(k+1)<cK<a(j+1)-bk.
\]
Define
\[
\begin{aligned}
G_k^{\rm out}&=P(R(G_k^{\rm in}))\ \ (0\le k\le s),\\
G_k^{\rm in}&=\operatorname{Lift}_{p_k}(G_{k-1}^{\rm out})\ \ (1\le k\le s).
\end{aligned}
\]
If \(G_s^{\rm out}=\varnothing\), then for \(M=M_0\prod_{k=1}^sp_k\),
\[
\forall\xi>0\ \forall N\in\mathbb N\ \exists n\ge\max(N,1):\
\gcd(Q_n,M)>1.
\tag{20}\label{eq:20}
\]
\label{thm-end:19}
\end{theorem}
\end{samepage}

% Markdown AST block 111: Proof.
\begin{proof}\leavevmode\label{proof:19}
Fix \(\xi>0\) and suppose (20) fails: there is \(n_0\ge1\)
with
\[
\gcd(Q_n,M)=1
\]
for all \(n\ge n_0\). By (H), \(c_n\in\mathcal C\)
for all \(n\ge n_0\).

% Markdown AST block 112: Para
The cell
\(j_n=\lfloor K\theta_n\rfloor\)
lies in
\(\{0,\ldots,K-1\}\), and Proposition 5 with \(M=1\) (where (E2) is
trivial) shows that
\[
j_n\xrightarrow{c_n}j_{n+1}
\]
is an edge of
\(G_0^{\rm in}\). Hence
\[
G_0^{\rm in}\models n_0
\]
with all words of
length one.

% Markdown AST block 113: Para
Proposition 13 and Proposition 16 give
\[
G_0^{\rm out}\models n_1.
\]

% Markdown AST block 114: Para
For \(k\ge1\), since \(p_k\mid M\) we have
\(p_k\nmid Q_n\) for \(n\ge n_0\), so Proposition 17 gives
\[
G_k^{\rm in}\models n'_k,
\]
and again Propositions 13 and 16 give
\[
G_k^{\rm out}\models n_{k+1}.
\]

% Markdown AST block 115: Para
After finitely many steps
\[
G_s^{\rm out}\models n_{s+1},
\]
but the empty graph has no infinite
walk. \qedhere
\end{proof}

% Markdown AST block 116: Para
The quantification in (20) is the same as in (18); the graphs take no
\(\xi\) as input. The primes of \(S_0\) are not lifted; they enter only
through the alphabet, and (H) is an arithmetic statement to be verified
by hand for each base (Lemmas 23 and 24).

% Markdown AST block 117: Corollary 20.
\begin{samepage}
\begin{corollary}\leavevmode\label{thm:20}
Under the hypotheses of Theorem 19, with
\[
P_{\max}=\max(S_0\cup\{p_1,\ldots,p_s\}),
\]
for every \(\xi>0\) and
every \(N\) there is \(n\ge\max(N,1)\) with \(Q_n\) composite.
\label{thm-end:20}
\end{corollary}
\end{samepage}

% Markdown AST block 118: Proof.
\begin{proof}\leavevmode\label{proof:20}
By (7), \(Q_n>P_{\max}\) for all large \(n\). Apply (20)
beyond that threshold and beyond \(N\). A prime \(p\) dividing
\(\gcd(Q_n,M)\) satisfies
\[
p\le P_{\max}<Q_n,
\]
so \(Q_n=pv\) with
\(v>1\). \qedhere
\end{proof}

% Markdown AST block 119: 4.6 The all-edge test and commuting closed words
\Needspace{6\baselineskip}
\subsection{The all-edge test and commuting closed words}\label{sec:4.6}

% Markdown AST block 120: Theorem 21.
\begin{samepage}
\begin{theorem}\leavevmode\label{thm:21}
Let \(C\) be a cyclic SCC of a word-labelled graph and
\(\rho\in C\). The following are equivalent.

% Markdown AST block 121: OrderedList
\begin{enumerate}
\def\labelenumi{(\alph{enumi})}
\item
  There exist an integer-valued \(h\) on \(C\) and \(\lambda\)
  satisfying (19); moreover, if this holds for some \(h\), then it holds
  when \(h(v)\) is the total word length of any chosen path from \(\rho\)
  to \(v\) inside \(C\).
\item
  Every infinite walk inside \(C\) has eventually periodic output.
\item
  Any two nonempty closed output words at \(\rho\) (outputs of nonempty
  closed walks from \(\rho\) to \(\rho\) inside \(C\)) commute.
\end{enumerate}
\label{thm-end:21}
\end{theorem}
\end{samepage}

% Markdown AST block 122: Proof.
\begin{proof}\leavevmode\label{proof:21}
(a)\(\Rightarrow\)(b) is Lemma 12(b).

% Markdown AST block 123: Para
(b)\(\Rightarrow\)(c). Suppose closed words \(u,v\) at \(\rho\) with \(uv\ne vu\).
Let \(U,V\) be the corresponding closed walks; \(UV\) and \(VU\) are
closed walks with outputs \(x_0=uv\) and \(x_1=vu\) of the same length
\(L\), \(x_0\ne x_1\). For \(\varepsilon\in\{0,1\}^{\mathbb N}\) the walk
\(X_{\varepsilon_0}X_{\varepsilon_1}\cdots\) stays in \(C\) and has
output
\[
z_\varepsilon=x_{\varepsilon_0}x_{\varepsilon_1}\cdots.
\]

% Markdown AST block 124: Para
Reading
\(z_\varepsilon\) in blocks of length \(L\) recovers \(\varepsilon\). If
\(z_\varepsilon\) were eventually periodic with period \(P\), it would
have period \(PL\), and block decoding would make \(\varepsilon\)
eventually periodic. Choosing \(\varepsilon\) not eventually periodic
(for instance \(0\,1\,00\,1\,000\,1\cdots\)) contradicts (b).

% Markdown AST block 125: Para
(c)\(\Rightarrow\)(a). Two commuting nonempty words are powers of a common word
(induction on length, removing the shorter word from the front of the
longer), and the primitive root of a nonempty word is unique, so all
closed words at \(\rho\) are powers of one primitive word \(z\); put
\(p=|z|\). For each \(v\in C\) fix a path \(\pi_v\) from \(\rho\) to
\(v\) with
\[
h(v)=|\operatorname{out}(\pi_v)|
\]
and a return path \(\sigma_v\) from \(v\) to \(\rho\). The paths at
\(v=\rho\) are allowed to be empty.

% Markdown AST block 126: Para
For an internal edge \(e=(u,w,v)\) the words
\[
\begin{aligned}
\omega_1&=\operatorname{out}(\pi_u)\,w\,\operatorname{out}(\sigma_v)\qquad \text{and}\\
\omega_2&=\operatorname{out}(\pi_v)\operatorname{out}(\sigma_v)
\end{aligned}
\]
are closed words, with \(\omega_1\ne\varepsilon\). Both are powers
of \(z\), allowing the zeroth power when \(\omega_2\) is empty.
Therefore \(p\) divides
\[
|\omega_1|-|\omega_2|=h(u)+|w|-h(v);
\]
hence \(p\mid\delta(e)\) for every
internal edge and \(p\mid d\). Define
\[
\lambda(s)=z_{s\bmod p}
\]
for \(s\in\mathbb Z/d\), well defined since
\(p\mid d\).

% Markdown AST block 127: Para
As \(\omega_1=z^k\), its letter at position \(t\) is
\(z_{t\bmod p}\), and \(w_i\) sits at position \(h(u)+i\) of
\(\omega_1\), so
\[
w_i=z_{(h(u)+i)\bmod p}=\lambda((h(u)+i)\bmod d).
\]
This is (19) for the root path lengths \(h\). \qedhere
\end{proof}

% Markdown AST block 128: Corollary 22.
\begin{samepage}
\begin{corollary}\leavevmode\label{thm:22}
Choose a root \(\rho\in C\) and paths
\(\pi_v:\rho\to v\) inside \(C\), and put
\[
h(v)=|\operatorname{out}(\pi_v)|.
\]
Whether (19) holds is independent of the root and of these paths.
For every such choice,
\[
d=g_C:=\gcd\bigl\{|\operatorname{out}(\gamma)|:
\gamma\text{ is a nonempty closed walk in }C\bigr\}.
\]
In particular, two complete tests using root path lengths agree on
which components are certified and on the value of \(d\).
\label{thm-end:22}
\end{corollary}
\end{samepage}

% Markdown AST block 129: Proof.
\begin{proof}\leavevmode\label{proof:22}
By Theorem 21, the test holds for any such choice exactly
when every internal infinite walk has eventually periodic output.
This property does not involve the root or paths.

% Markdown AST block 130: Para
For the equality of gcds, summing \(\delta(e)\) around a closed walk
gives its total word length. Hence \(d\mid g_C\). Conversely, for
\(e=(u,w,v)\), choose a return path \(\sigma_v:v\to\rho\). Then
\[
\delta(e)
=|\operatorname{out}(\pi_u e\sigma_v)|
 -|\operatorname{out}(\pi_v\sigma_v)|.
\]
Both terms are lengths of closed walks at \(\rho\). Each is divisible
by \(g_C\), including a zero length if the second walk is empty.
Thus \(g_C\mid\delta(e)\) for every internal edge, so \(g_C\mid d\).
This proves \(d=g_C\), whether or not the component is certified. \qedhere
\end{proof}

% Markdown AST block 131: Para
The restriction on \(h\) is essential. Consider the two-edge cycle
\[
u\xrightarrow{(2)}v,\qquad v\xrightarrow{(4)}u.
\]
Its infinite outputs are the two rotations of \((2,4)^\infty\).
The root path lengths \(h(u)=0\), \(h(v)=1\) give
\[
(\delta(u\to v),\delta(v\to u))=(0,2),\qquad d=2,
\]
and (19) holds with \(\lambda(0)=2\), \(\lambda(1)=4\).
The arbitrary choice \(h(u)=h(v)=0\) instead gives
\[
(\delta(u\to v),\delta(v\to u))=(1,1),\qquad d=1.
\]
Both letters would then occupy phase zero, so (19) fails. For a
general integer-valued \(h\), the telescoping argument gives only
\(d\mid g_C\); Lemma 12 still guarantees soundness, but failure of
that test need not imply a nonperiodic output.

% Markdown AST block 132: Para
With root path lengths, Theorem 21 does give the converse: if the full
test (19) fails, the component contains an internal infinite walk whose
output is not eventually periodic. No sound test that removes a
component solely because all its infinite outputs are eventually
periodic can remove such a component from the same finite graph.
Further removal requires additional arithmetic or realizability
information.

% Markdown AST block 133: 5. The two rational bases
\Needspace{6\baselineskip}
\section{The two rational bases}\label{sec:5}

% Markdown AST block 134: 5.1 The ratio 7/5 by the one-step method
\subsection{The ratio 7/5 by the one-step method}\label{sec:5.1}

% Markdown AST block 135: Para
Use
\[
(a,b,M,K_0,f)=(7,5,4290,32,4).
\tag{21}\label{eq:21}
\]
The initial graph has \(\varphi(4290)K_0=960\cdot32=30720\) vertices.
The two one-step implementations described in Section 9.2 produce
the four stages in Table 1. The retained set is empty at \(K=2048\).

% Markdown AST block 136: Table 1.
% Caption is rendered at the start of the immediately following table.

% Markdown AST block 137: Table 1
\begingroup
\small
\setlength{\tabcolsep}{4pt}
\renewcommand{\arraystretch}{1.15}
\setlength{\LTcapwidth}{\textwidth}
\begin{longtable}{@{} r r r r r r@{}}
\caption{The complete stages for (21). Edges are counted with their
labels. The certified column counts certified cyclic SCCs, and the final
column is \(|U_t|\).}\label{tab:1}\\
\toprule
\(K_t\) & Vertices & Edges & \shortstack[r]{Cyclic\\SCCs} & \shortstack[r]{Certified\\SCCs} & \shortstack[r]{Retained\\vertices} \\
\midrule
\endfirsthead
\toprule
\(K_t\) & Vertices & Edges & \shortstack[r]{Cyclic\\SCCs} & \shortstack[r]{Certified\\SCCs} & \shortstack[r]{Retained\\vertices} \\
\midrule
\endhead
\bottomrule
\endlastfoot
32 & 30720 & 55440 & 17 & 16 & 1363 \\*
128 & 5452 & 9512 & 3 & 2 & 952 \\*
512 & 3808 & 6281 & 1 & 0 & 944 \\*
2048 & 3776 & 5951 & 15 & 15 & 0 \\
\end{longtable}
\endgroup

% Markdown AST block 138: Para
The numbers of all SCCs, including noncyclic singletons, are respectively
29312, 4440, 2865, and 3658. At the first three stages there is one
uncertified cyclic SCC, of respective size 1363, 952, and 944. Their
internal edge counts are 2332, 1550, and 1505, and all three have \(d=1\).
Their removal is not justified by the phase test; they are carried into
the next subdivision. In particular, the vertex counts of the next stages
are exactly four times these retained sizes.

% Markdown AST block 139: Para
Table 2 groups certified components by \(K,d\), and normalized primitive
word. Component sizes are listed with multiplicity. There are no certified
components at \(K=512\).

% Markdown AST block 140: Table 2.
% Caption is rendered at the start of the immediately following table.

% Markdown AST block 141: Table 2
\begingroup
\small
\setlength{\tabcolsep}{4pt}
\renewcommand{\arraystretch}{1.15}
\setlength{\LTcapwidth}{\textwidth}
\begin{longtable}{@{}>{\raggedright\arraybackslash}p{\dimexpr0.055\dimexpr\linewidth-40pt\relax} >{\raggedright\arraybackslash}p{\dimexpr0.1\dimexpr\linewidth-40pt\relax} >{\raggedright\arraybackslash}p{\dimexpr0.19\dimexpr\linewidth-40pt\relax} >{\raggedright\arraybackslash}p{\dimexpr0.035\dimexpr\linewidth-40pt\relax} >{\raggedright\arraybackslash}p{\dimexpr0.035\dimexpr\linewidth-40pt\relax} >{\raggedright\arraybackslash}p{\dimexpr0.485\dimexpr\linewidth-40pt\relax}@{}}
\caption{All nine groups of certified components. The parameter \(d\)
is (14), and \(p\) is the length of the displayed primitive word.}\label{tab:2}\\
\toprule
\(K\) & Number of SCCs & Component sizes & \(d\) & \(p\) & Primitive word, up to rotation \\
\midrule
\endfirsthead
\toprule
\(K\) & Number of SCCs & Component sizes & \(d\) & \(p\) & Primitive word, up to rotation \\
\midrule
\endhead
\bottomrule
\endlastfoot
32 & 2 & 1,\allowbreak{} 1 & 1 & 1 & \((2)\) \\
32 & 6 & 2,\allowbreak{} 2,\allowbreak{} 2,\allowbreak{} 2,\allowbreak{} 2,\allowbreak{} 2 & 2 & 2 & \((-2,\allowbreak{}4)\) \\
32 & 8 & 6,\allowbreak{} 6,\allowbreak{} 6,\allowbreak{} 6,\allowbreak{} 6,\allowbreak{} 6,\allowbreak{} 6,\allowbreak{} 6 & 6 & 2 & \((-2,\allowbreak{}4)\) \\
128 & 1 & 24 & 8 & 8 & \((-4,\allowbreak{}2,\allowbreak{}2,\allowbreak{}2,\allowbreak{}2,\allowbreak{}2,\allowbreak{}2,\allowbreak{}2)\) \\
128 & 1 & 39 & 13 & 13 & \((-4,\allowbreak{}2,\allowbreak{}2,\allowbreak{}2,\allowbreak{}2,\allowbreak{}2,\allowbreak{}-4,\allowbreak{}2,\allowbreak{}2,\allowbreak{}2,\allowbreak{}2,\allowbreak{}2,\allowbreak{}2)\) \\
2048 & 3 & 1,\allowbreak{} 1,\allowbreak{} 1 & 1 & 1 & \((2)\) \\
2048 & 3 & 4,\allowbreak{} 4,\allowbreak{} 8 & 4 & 4 & \((-2,\allowbreak{}-2,\allowbreak{}4,\allowbreak{}4)\) \\
2048 & 3 & 6,\allowbreak{} 6,\allowbreak{} 6 & 6 & 1 & \((2)\) \\
2048 & 6 & 12,\allowbreak{} 12,\allowbreak{} 12,\allowbreak{} 12,\allowbreak{} 24,\allowbreak{} 24 & 12 & 4 & \((-2,\allowbreak{}-2,\allowbreak{}4,\allowbreak{}4)\) \\
\end{longtable}
\endgroup

% Markdown AST block 142: Para
Thus five distinct primitive words occur. The rows with \(d=6,p=1\)
and \(d=12,p=4\), for example, show why \(d\) must not be reported as
the primitive word length. The words describe possible output within the
outer graph components; an infinite true orbit cannot remain in any of
these components, by Lemma 4.

% Markdown AST block 143: Para
As additional checks of the procedure, Table 3 applies the same
construction at \(K_0=32\) to the three floor results of
\hyperref[ref-dn]{DN, Theorem 1}. Each stops after its first stage. These are
reproductions of the stated divisibility conclusions, not new claims
about those three bases.

% Markdown AST block 144: Table 3.
% Caption is rendered at the start of the immediately following table.

% Markdown AST block 145: Table 3
\begingroup
\small
\setlength{\tabcolsep}{4pt}
\renewcommand{\arraystretch}{1.15}
\setlength{\LTcapwidth}{\textwidth}
\begin{longtable}{@{} r r r r r r r@{}}
\caption{Known cases reproduced by both one-step
implementations; all retained sets are empty.}\label{tab:3}\\
\toprule
\(a/b\) & \shortstack[r]{Prime\\set} & \(M\) & Vertices & Edges & \shortstack[r]{Cyclic\\SCCs} & \shortstack[r]{Certified\\SCCs} \\
\midrule
\endfirsthead
\toprule
\(a/b\) & \shortstack[r]{Prime\\set} & \(M\) & Vertices & Edges & \shortstack[r]{Cyclic\\SCCs} & \shortstack[r]{Certified\\SCCs} \\
\midrule
\endhead
\bottomrule
\endlastfoot
\(3/2\) & \(\{2,\allowbreak{}5,\allowbreak{}7,\allowbreak{}11\}\) & 770 & 7680 & 8640 & 6 & 6 \\*
\(4/3\) & \(\{2,\allowbreak{}3,\allowbreak{}5\}\) & 30 & 256 & 480 & 3 & 3 \\*
\(5/4\) & \(\{2,\allowbreak{}3,\allowbreak{}7,\allowbreak{}11,\allowbreak{}13\}\) & 6006 & 46080 & 82170 & 18 & 18 \\
\end{longtable}
\endgroup

% Markdown AST block 146: 5.2 The ratio 7/5 by the word method
\Needspace{6\baselineskip}
\subsection{The ratio 7/5 by the word method}\label{sec:5.2}

% Markdown AST block 147: Lemma 23.
\begin{samepage}
\begin{lemma}\leavevmode\label{thm:23}
Let \(a=7\), \(b=5\). If \(Q_n\) and \(Q_{n+1}\) are odd
and \(5\nmid Q_n\), then
\[
c_n\in\{-4,-2,2,4,6\}.
\]
Hence (H) holds
with
\[
S_0=\{2,5\}\qquad \text{and}\qquad \mathcal C=\{-4,-2,2,4,6\}.
\]
\label{thm-end:23}
\end{lemma}
\end{samepage}

% Markdown AST block 148: Proof.
\begin{proof}\leavevmode\label{proof:23}
By (6), \(-4\le c_n\le6\). Since \(5\equiv7\equiv1\pmod2\),
\[
c_n=5Q_{n+1}-7Q_n\equiv Q_{n+1}-Q_n\pmod2,
\]
which is even when both
terms are odd; so
\[
c_n\in\{-4,-2,0,2,4,6\}.
\]
If \(c_n=0\) then
\(5Q_{n+1}=7Q_n\) and, as \(\gcd(5,7)=1\), \(5\mid Q_n\). \qedhere
\end{proof}

% Markdown AST block 149: Para
The prime \(2\) could also be lifted by Proposition 17, but
\(\mathbb F_2^\times=\{1\}\) and the lift condition reduces to
``all letters even'', already contained in the alphabet. The prime
\(5\) divides \(b\) and cannot be lifted; by the remark after Lemma 18,
\(5\nmid Q_n\) is equivalent to \(c_n\not\equiv0\pmod5\), which within
the range \([-4,6]\) and with even letters means \(c_n\ne0\). Thus the
alphabet expresses the avoidance of \(2\) and \(5\) as a necessary
condition, which is all that Theorem 19 requires.

% Markdown AST block 150: Para
Apply Theorem 19 with \(K=2048\), the alphabet of Lemma 23, and the
primes \(p_1,p_2,p_3=3,11,13\). The value \(K=2048\) is a free
parameter, chosen to match the final cell count of Table 1; soundness
does not depend on it. Table 4 lists the four stages. Both
implementations of Section 9.3 agree on every vertex and every
word-labelled edge of the eight graphs; the output graph of the last
stage is empty. In these four stages no macro edge is produced by more
than one contracted chain, so the number of distinct edges equals the
number counted with multiplicity.

% Markdown AST block 151: Table 4.
% Caption is rendered at the start of the immediately following table.

% Markdown AST block 152: Table 4
\begingroup
\small
\setlength{\tabcolsep}{2.5pt}
\renewcommand{\arraystretch}{1.15}
\setlength{\LTcapwidth}{\textwidth}
\begin{longtable}{@{} r r r r r r r r r r@{}}
\caption{The word method for \(7/5\) with \(K=2048\),
\(\mathcal C=\{-4,-2,2,4,6\}\), primes \(3,11,13\). Input graphs are
\(G_k^{\rm in}\) and output graphs \(G_k^{\rm out}\) of Theorem 19.
``Uncertified vertices'' counts the vertices of uncertified cyclic SCCs.}\label{tab:4}\\
\toprule
\shortstack[r]{\(k\)} & \shortstack[r]{\(p_k\)} & \shortstack[r]{Input\\vertices} & \shortstack[r]{Input\\edges} & \shortstack[r]{Cyclic\\SCCs} & \shortstack[r]{Certified} & \shortstack[r]{Uncertified\\vertices} & \shortstack[r]{Output\\vertices} & \shortstack[r]{Output\\edges} & \shortstack[r]{Max.\\word\\length\\(output)} \\
\midrule
\endfirsthead
\toprule
\shortstack[r]{\(k\)} & \shortstack[r]{\(p_k\)} & \shortstack[r]{Input\\vertices} & \shortstack[r]{Input\\edges} & \shortstack[r]{Cyclic\\SCCs} & \shortstack[r]{Certified} & \shortstack[r]{Uncertified\\vertices} & \shortstack[r]{Output\\vertices} & \shortstack[r]{Output\\edges} & \shortstack[r]{Max.\\word\\length\\(output)} \\
\midrule
\endhead
\bottomrule
\endlastfoot
0 & --- & 2048 & 8366 & 1 & 0 & 2048 & 2048 & 8366 & 1 \\*
1 & 3 & 4096 & 9010 & 10 & 9 & 691 & 593 & 1292 & 13 \\*
2 & 11 & 5928 & 11276 & 4 & 3 & 89 & 62 & 131 & 8 \\*
3 & 13 & 744 & 1342 & 2 & 2 & 0 & 0 & 0 & --- \\
\end{longtable}
\endgroup

% Markdown AST block 153: Para
The total numbers of SCCs at the four stages are 1, 3387, 5840, and
704. The final prime set is \(\{2,3,5,11,13\}\), as in (21). This is
not a theorem about a new base but a second proof of (1) by the method of
Section 4, and a consistency check between the two methods.

% Markdown AST block 154: 5.3 The ratio 5/2
\Needspace{6\baselineskip}
\subsection{The ratio 5/2}\label{sec:5.3}

% Markdown AST block 155: Lemma 24.
\begin{samepage}
\begin{lemma}\leavevmode\label{thm:24}
Let \(a=5\), \(b=2\). If \(Q_n\) is odd, then
\[
c_n\in\{-1,1,3\}.
\]
Hence (H) holds with
\[
S_0=\{2\}\qquad \text{and}\qquad \mathcal C=\{-1,1,3\}.
\]
\label{thm-end:24}
\end{lemma}
\end{samepage}

% Markdown AST block 156: Proof.
\begin{proof}\leavevmode\label{proof:24}
By (6), \(-1\le c_n\le4\), and
\[
c_n=2Q_{n+1}-5Q_n\equiv Q_n\pmod2.\qedhere
\]
\end{proof}

% Markdown AST block 157: Para
Only the oddness of \(Q_n\) is used; neither the oddness of \(Q_{n+1}\)
nor avoidance of \(5\) is needed. The prime \(5\) is not included in the
modulus of (2), for three independent reasons: \(5\mid a\), so
Lemma 18 does not apply; by the remark after Lemma 18, a lift by \(5\)
would give
\[
Q_{n+1}\equiv2^{-1}c_n=3c_n\pmod5,
\]
which is nonzero for
\(c_n\in\{-1,1,3\}\) (residues \(2,3,4\)), so on a tail with odd carries
\(5\nmid Q_{n+1}\) holds automatically and the lift would remove
nothing; and the statement without \(5\) is stronger.

% Markdown AST block 158: Para
With \(K=4\), the initial graph \(G_0^{\rm in}\) has the four vertices
\(j\in\{0,1,2,3\}\) and the twelve edges \((j,(c),k)\) with
\(c\in\{-1,1,3\}\) and
\[
5j-2(k+1)<4c<5(j+1)-2k,
\]
listed in Table 5.
By Theorem 19 every tail avoiding the primes of (2) walks in this graph;
this is a covering of all cells, not a sampling of initial values, and
no converse realizability is asserted.

% BEGIN PUBLIC ADDITION figure:initial-graph
The initial four-cell graph is small enough to display completely, grouped by carry label.

\begin{figure}[htbp]
\centering
\begin{tikzpicture}[paperfig]
\begin{scope}[xshift=0.0cm]
\node[ftitle] at(1.35,2.4) {$c=-1$};
\node[fnode] (n0) at(0,1.1) {$j=0$};
\node[fnode] (n1) at(2.7,1.1) {$j=1$};
\node[fnode] (n2) at(0,-1.1) {$j=2$};
\node[fnode] (n3) at(2.7,-1.1) {$j=3$};
% exact edge: 0 -1 2
\draw[fedge] (n0) -- (n2);
% exact edge: 0 -1 3
\draw[fedge] (n0) -- (n3);
\end{scope}
\begin{scope}[xshift=4.3cm]
\node[ftitle] at(1.35,2.4) {$c=1$};
\node[fnode] (n0) at(0,1.1) {$j=0$};
\node[fnode] (n1) at(2.7,1.1) {$j=1$};
\node[fnode] (n2) at(0,-1.1) {$j=2$};
\node[fnode] (n3) at(2.7,-1.1) {$j=3$};
% exact edge: 0 1 0
\draw[fedge] (n0) edge[loop above] (n0);
% exact edge: 1 1 0
\draw[fedge] (n1) -- (n0);
% exact edge: 1 1 1
\draw[fedge] (n1) edge[loop above] (n1);
% exact edge: 1 1 2
\draw[fedge] (n1) -- (n2);
% exact edge: 2 1 3
\draw[fedge] (n2) -- (n3);
\end{scope}
\begin{scope}[xshift=8.6cm]
\node[ftitle] at(1.35,2.4) {$c=3$};
\node[fnode] (n0) at(0,1.1) {$j=0$};
\node[fnode] (n1) at(2.7,1.1) {$j=1$};
\node[fnode] (n2) at(0,-1.1) {$j=2$};
\node[fnode] (n3) at(2.7,-1.1) {$j=3$};
% exact edge: 2 3 0
\draw[fedge] (n2) -- (n0);
% exact edge: 2 3 1
\draw[fedge] (n2) -- (n1);
% exact edge: 3 3 1
\draw[fedge] (n3) -- (n1);
% exact edge: 3 3 2
\draw[fedge] (n3) -- (n2);
% exact edge: 3 3 3
\draw[fedge] (n3) edge[loop below] (n3);
\end{scope}
\node[align=center,text width=12.5cm] at(5.65,-2.1)
{$5j-2(k+1)<4c<5(j+1)-2k$\\[4pt]
\figtext{The three panels together contain all 12 labelled edges.}{三つの図を合わせると、ラベル付きの12辺すべてを含む。}};
\end{tikzpicture}

\caption{All twelve edges of the initial \(5/2\) graph at \(K=4\), separated into three panels by carry \(c\). Each panel uses the same four vertices; the panels are overlaid to recover the full graph. The edges are enumerated by the strict integer inequality shown below them and agree with Table 5.}
\label{fig:initial-graph}
\end{figure}
% END PUBLIC ADDITION figure:initial-graph

% Markdown AST block 159: Table 5.
% Caption is rendered at the start of the immediately following table.

% Markdown AST block 160: Table 5
\begingroup
\small
\setlength{\tabcolsep}{4pt}
\renewcommand{\arraystretch}{1.15}
\setlength{\LTcapwidth}{\textwidth}
\begin{longtable}{@{}>{\raggedright\arraybackslash}p{\dimexpr0.19\dimexpr\linewidth-8pt\relax} >{\raggedright\arraybackslash}p{\dimexpr0.81\dimexpr\linewidth-8pt\relax}@{}}
\caption{The twelve initial edges for \(5/2\), \(K=4\), as pairs
(carry, target cell).}\label{tab:5}\\
\toprule
Source cell \(j\) & Edges \((c,k)\) \\
\midrule
\endfirsthead
\toprule
Source cell \(j\) & Edges \((c,k)\) \\
\midrule
\endhead
\bottomrule
\endlastfoot
0 & \((-1,\allowbreak{}2),\allowbreak{}(-1,\allowbreak{}3),\allowbreak{}(1,\allowbreak{}0)\) \\*
1 & \((1,\allowbreak{}0),\allowbreak{}(1,\allowbreak{}1),\allowbreak{}(1,\allowbreak{}2)\) \\*
2 & \((1,\allowbreak{}3),\allowbreak{}(3,\allowbreak{}0),\allowbreak{}(3,\allowbreak{}1)\) \\*
3 & \((3,\allowbreak{}1),\allowbreak{}(3,\allowbreak{}2),\allowbreak{}(3,\allowbreak{}3)\) \\
\end{longtable}
\endgroup

% Markdown AST block 161: Para
Apply Theorem 19 with \(K=4\), the alphabet of Lemma 24, and the nine
primes
\[
p_1,\ldots,p_9=3,7,11,13,17,19,23,29,31.
\]
Table 6 lists the
ten stages. Both implementations of Section 9.3 agree on every vertex and
every word-labelled edge of the twenty graphs, and the output graph of
stage \(9\) is empty: all 1084 cyclic SCCs of \(G_9^{\rm in}\) are
certified. Edges are counted as sets of triples (Section 4.1); the
multiplicities with which contracted chains produce the same triple carry
no proof content and are recorded in Appendix A for comparison with the
supplied research notes.

% Markdown AST block 162: Table 6.
% Caption is rendered at the start of the immediately following table.

% Markdown AST block 163: Table 6
\begingroup
\small
\setlength{\tabcolsep}{2.5pt}
\renewcommand{\arraystretch}{1.15}
\setlength{\LTcapwidth}{\textwidth}
\begin{longtable}{@{} r r r r r r r r r r@{}}
\caption{The word method for \(5/2\) with \(K=4\),
\(\mathcal C=\{-1,1,3\}\). Columns as in Table 4.}\label{tab:6}\\
\toprule
\shortstack[r]{\(k\)} & \shortstack[r]{\(p_k\)} & \shortstack[r]{Input\\vertices} & \shortstack[r]{Input\\edges} & \shortstack[r]{Cyclic\\SCCs} & \shortstack[r]{Certified} & \shortstack[r]{Uncertified\\vertices} & \shortstack[r]{Output\\vertices} & \shortstack[r]{Output\\edges} & \shortstack[r]{Max.\\word\\length\\(output)} \\
\midrule
\endfirsthead
\toprule
\shortstack[r]{\(k\)} & \shortstack[r]{\(p_k\)} & \shortstack[r]{Input\\vertices} & \shortstack[r]{Input\\edges} & \shortstack[r]{Cyclic\\SCCs} & \shortstack[r]{Certified} & \shortstack[r]{Uncertified\\vertices} & \shortstack[r]{Output\\vertices} & \shortstack[r]{Output\\edges} & \shortstack[r]{Max.\\word\\length\\(output)} \\
\midrule
\endhead
\bottomrule
\endlastfoot
0 & --- & 4 & 12 & 1 & 0 & 4 & 4 & 12 & 1 \\
1 & 3 & 8 & 17 & 1 & 0 & 7 & 5 & 12 & 2 \\
2 & 7 & 30 & 58 & 4 & 2 & 16 & 9 & 19 & 4 \\
3 & 11 & 90 & 159 & 3 & 0 & 68 & 39 & 84 & 8 \\
4 & 13 & 462 & 812 & 4 & 1 & 342 & 250 & 524 & 15 \\
5 & 17 & 3985 & 7020 & 11 & 8 & 2913 & 2059 & 4058 & 38 \\
6 & 19 & 36843 & 57660 & 73 & 70 & 15933 & 8890 & 16696 & 56 \\
7 & 23 & 194242 & 293181 & 23 & 21 & 84824 & 40051 & 73286 & 98 \\
8 & 29 & 1101968 & 1542910 & 311 & 310 & 269245 & 100401 & 171256 & 210 \\
9 & 31 & 2856938 & 3256019 & 1084 & 1084 & 0 & 0 & 0 & --- \\
\end{longtable}
\endgroup

% Markdown AST block 164: Para
The total numbers of SCCs at the ten stages are 1, 2, 16, 25, 123, 1060,
20800, 109409, 832175, and 2850774. The maximal word length of an input
graph can be smaller than that of the preceding output graph (207 at
stage 9 against 210 at stage 8), because a lift discards every edge whose
word forces a zero residue for all starting residues. The anchor rule of
Section 4.3 never fired. A direct product of all residues modulo the
modulus of (2) with four cells would have
\[
4\varphi(40112098026)=30656102400
\]
vertices; the staged construction
never forms it.

% BEGIN PUBLIC ADDITION figure:stage-counts
The vertex counts show how growth under lifting alternates with reduction inside each stage.

\begin{figure}[htbp]
\centering
\begin{tikzpicture}[paperfig,x=1.16cm,y=.61cm]
\draw[fedge] (-.3,0)--(9.5,0) node[right] {$k$};
\draw[fedge] (0,-.2)--(0,7.1);
\node[anchor=west] at(.05,7.3) {$\log_{10}(N+1)$};
\draw[black!18] (0,1)--(9,1);
\node[anchor=east] at(-.12,1) {$ 1 $};
\draw[black!18] (0,2)--(9,2);
\node[anchor=east] at(-.12,2) {$ 2 $};
\draw[black!18] (0,3)--(9,3);
\node[anchor=east] at(-.12,3) {$ 3 $};
\draw[black!18] (0,4)--(9,4);
\node[anchor=east] at(-.12,4) {$ 4 $};
\draw[black!18] (0,5)--(9,5);
\node[anchor=east] at(-.12,5) {$ 5 $};
\draw[black!18] (0,6)--(9,6);
\node[anchor=east] at(-.12,6) {$ 6 $};
\draw (0,0)--(0,-.12);
\node[below] at(0,-.12) {$ 0 $};
\node[below] at(0,-.62) {$ -- $};
\draw (1,0)--(1,-.12);
\node[below] at(1,-.12) {$ 1 $};
\node[below] at(1,-.62) {$ 3 $};
\draw (2,0)--(2,-.12);
\node[below] at(2,-.12) {$ 2 $};
\node[below] at(2,-.62) {$ 7 $};
\draw (3,0)--(3,-.12);
\node[below] at(3,-.12) {$ 3 $};
\node[below] at(3,-.62) {$ 11 $};
\draw (4,0)--(4,-.12);
\node[below] at(4,-.12) {$ 4 $};
\node[below] at(4,-.62) {$ 13 $};
\draw (5,0)--(5,-.12);
\node[below] at(5,-.12) {$ 5 $};
\node[below] at(5,-.62) {$ 17 $};
\draw (6,0)--(6,-.12);
\node[below] at(6,-.12) {$ 6 $};
\node[below] at(6,-.62) {$ 19 $};
\draw (7,0)--(7,-.12);
\node[below] at(7,-.12) {$ 7 $};
\node[below] at(7,-.62) {$ 23 $};
\draw (8,0)--(8,-.12);
\node[below] at(8,-.12) {$ 8 $};
\node[below] at(8,-.62) {$ 29 $};
\draw (9,0)--(9,-.12);
\node[below] at(9,-.12) {$ 9 $};
\node[below] at(9,-.62) {$ 31 $};
\node[anchor=east] at(-.15,-.4) {$k$};
\node[anchor=east] at(-.15,-.9) {$p_k$};
\draw[solid,line width=.85pt] plot coordinates {(0,0.698970004) (1,0.954242509) (2,1.491361694) (3,1.959041392) (4,2.665580991) (5,3.600537294) (6,4.566366774) (7,5.288345377) (8,6.042169377) (9,6.455900968)};
\node[draw,circle,fill=white,inner sep=1.2pt] at(0,0.698970004){};
\node[draw,circle,fill=white,inner sep=1.2pt] at(1,0.954242509){};
\node[draw,circle,fill=white,inner sep=1.2pt] at(2,1.491361694){};
\node[draw,circle,fill=white,inner sep=1.2pt] at(3,1.959041392){};
\node[draw,circle,fill=white,inner sep=1.2pt] at(4,2.665580991){};
\node[draw,circle,fill=white,inner sep=1.2pt] at(5,3.600537294){};
\node[draw,circle,fill=white,inner sep=1.2pt] at(6,4.566366774){};
\node[draw,circle,fill=white,inner sep=1.2pt] at(7,5.288345377){};
\node[draw,circle,fill=white,inner sep=1.2pt] at(8,6.042169377){};
\node[draw,circle,fill=white,inner sep=1.2pt] at(9,6.455900968){};
\draw[solid,line width=.85pt] (.6,6.2)--(1.3,6.2);
\node[draw,circle,fill=white,inner sep=1.2pt] at(.95,6.2){};
\node[anchor=west] at(1.45,6.2) {\figtext{Input vertices}{入力頂点数}};
\draw[dashed,line width=.85pt] plot coordinates {(0,0.698970004) (1,0.903089987) (2,1.230448921) (3,1.838849091) (4,2.535294120) (5,3.464489547) (6,4.202324813) (7,4.928523868) (8,5.430149260) (9,0.000000000)};
\node[draw,rectangle,fill=white,inner sep=1.2pt] at(0,0.698970004){};
\node[draw,rectangle,fill=white,inner sep=1.2pt] at(1,0.903089987){};
\node[draw,rectangle,fill=white,inner sep=1.2pt] at(2,1.230448921){};
\node[draw,rectangle,fill=white,inner sep=1.2pt] at(3,1.838849091){};
\node[draw,rectangle,fill=white,inner sep=1.2pt] at(4,2.535294120){};
\node[draw,rectangle,fill=white,inner sep=1.2pt] at(5,3.464489547){};
\node[draw,rectangle,fill=white,inner sep=1.2pt] at(6,4.202324813){};
\node[draw,rectangle,fill=white,inner sep=1.2pt] at(7,4.928523868){};
\node[draw,rectangle,fill=white,inner sep=1.2pt] at(8,5.430149260){};
\node[draw,rectangle,fill=white,inner sep=1.2pt] at(9,0.000000000){};
\draw[dashed,line width=.85pt] (.6,5.65)--(1.3,5.65);
\node[draw,rectangle,fill=white,inner sep=1.2pt] at(.95,5.65){};
\node[anchor=west] at(1.45,5.65) {\figtext{After removal}{除去後の頂点数}};
\draw[densely dotted,line width=.85pt] plot coordinates {(0,0.698970004) (1,0.778151250) (2,1.000000000) (3,1.602059991) (4,2.399673721) (5,3.313867220) (6,3.948950610) (7,4.602624207) (8,5.001742364) (9,0.000000000)};
\node[draw,diamond,fill=white,inner sep=1.2pt] at(0,0.698970004){};
\node[draw,diamond,fill=white,inner sep=1.2pt] at(1,0.778151250){};
\node[draw,diamond,fill=white,inner sep=1.2pt] at(2,1.000000000){};
\node[draw,diamond,fill=white,inner sep=1.2pt] at(3,1.602059991){};
\node[draw,diamond,fill=white,inner sep=1.2pt] at(4,2.399673721){};
\node[draw,diamond,fill=white,inner sep=1.2pt] at(5,3.313867220){};
\node[draw,diamond,fill=white,inner sep=1.2pt] at(6,3.948950610){};
\node[draw,diamond,fill=white,inner sep=1.2pt] at(7,4.602624207){};
\node[draw,diamond,fill=white,inner sep=1.2pt] at(8,5.001742364){};
\node[draw,diamond,fill=white,inner sep=1.2pt] at(9,0.000000000){};
\draw[densely dotted,line width=.85pt] (.6,5.1)--(1.3,5.1);
\node[draw,diamond,fill=white,inner sep=1.2pt] at(.95,5.1){};
\node[anchor=west] at(1.45,5.1) {\figtext{After contraction}{縮約後の頂点数}};
\node[flabel,anchor=east] at(8.6,.75) {$N=0$};
\draw[fedge] (8.65,.6)--(8.98,.04);
\node[align=center,text width=12cm] at(4.5,-1.65)
{\figtext{A lift can enlarge the graph; removal and contraction reduce each stage.}{持ち上げでグラフは拡大し得る。各段階で除去と縮約を行う。}};
\end{tikzpicture}

\caption{Vertex counts from Table 6. The ordinate is \(\log_{10}(N+1)\), so the final empty output \(N=0\) is visible at height zero. Solid circles show input counts, dashed squares counts after removal, and dotted diamonds counts after contraction. Connecting segments guide the eye; they do not assert monotone decay across stages. All input cyclic SCCs are certified at stage nine.}
\label{fig:stage-counts}
\end{figure}
% END PUBLIC ADDITION figure:stage-counts

% Markdown AST block 165: 5.4 Proofs of Theorem 1 and Corollary 2
\Needspace{6\baselineskip}
\subsection{Proofs of Theorem 1 and Corollary 2}\label{sec:5.4}

% Markdown AST block 166: Proof of Theorem 1.
\begin{proof}[Proof of Theorem 1]\leavevmode\label{proof:1}
(i) The finite calculation verifies the graphs,
retained sets, and refinements in (17) for the parameters (21), with
\(s=3\) and \(U_3=\varnothing\). Verification includes the full edge
sets and component conditions, as detailed in Section 9; the summary
counts alone are not the certificate. Theorem 10 gives (1).

% Markdown AST block 167: Para
Alternatively, Table 4 verifies the hypotheses of Theorem 19 with
\(S_0=\{2,5\}\), \(\mathcal C=\{-4,-2,2,4,6\}\) (Lemma 23),
\(K=2048\), primes \(3,11,13\), and \(G_3^{\rm out}=\varnothing\);
then (20) with
\[
M=10\cdot3\cdot11\cdot13=4290
\]
is (1).

% Markdown AST block 168: OrderedList
\begin{enumerate}
\def\labelenumi{(\roman{enumi})}
\setcounter{enumi}{1}
\tightlist
\item
  Table 6 verifies the hypotheses of Theorem 19 with
  \(S_0=\{2\}\), \(\mathcal C=\{-1,1,3\}\) (Lemma 24), \(K=4\), the nine
  primes \(3,7,11,13,17,19,23,29,31\), each distinct, coprime to
  \(10\), and not in \(S_0\), and \(G_9^{\rm out}=\varnothing\). Then
% Markdown AST block 169: OrderedList
(20) with
\[
M=2\cdot3\cdot7\cdot11\cdot13\cdot17\cdot19\cdot23\cdot29\cdot31
=40112098026
\]
is (2). \qedhere
\end{enumerate}
\end{proof}

% Markdown AST block 170: Proof of Corollary 2.
\begin{proof}[Proof of Corollary 2]\leavevmode\label{proof:2}
By (7), eventually \(Q_n>13\) for the
sequence with ratio \(7/5\) and \(Q_n>31\) for the sequence with ratio
\(5/2\). Given any lower bound on the index, apply Theorem 1 beyond both
that bound and the growth threshold. A prime divisor of the gcd in (1)
lies in \(\{2,3,5,11,13\}\), and one of the gcd in (2) lies in
\(\{2,3,7,11,13,17,19,23,29,31\}\); in either case it is a proper divisor
of \(Q_n\). Hence \(Q_n\) is composite, at arbitrarily large indices.
This is Corollary 20 in both cases. \qedhere
\end{proof}

% Markdown AST block 171: 6. Quantitative waiting time for 7/5
\Needspace{6\baselineskip}
\section{Quantitative waiting time for 7/5}\label{sec:6}

% Markdown AST block 172: 6.1 A finite version of nonperiodicity
\subsection{A finite version of nonperiodicity}\label{sec:6.1}

% Markdown AST block 173: Lemma 25.
\begin{samepage}
\begin{lemma}\leavevmode\label{thm:25}
Let \(u\in\mathbb N\), \(1\le d\le\ell\), and suppose
\[
c_{u+i+d}=c_{u+i}
\]
for \(0\le i<\ell-d\). Put
\[
D_i=Q_{u+i+d}-Q_{u+i}.
\]
Then
\[
b^{\ell-d}\mid D_0.
\tag{22}\label{eq:22}
\]
If \(D_0>0\), then with
\(\nu_b(D_0)=\max\{k\ge0:b^k\mid D_0\}\),
\[
\begin{aligned}
\ell&\le d+\nu_b(D_0),\\
b^{\ell-d}&\le D_0<r^d(Q_u+1)-Q_u.
\end{aligned}
\tag{23}\label{eq:23}
\]
\label{thm-end:25}
\end{lemma}
\end{samepage}

% Markdown AST block 174: Proof.
\begin{proof}\leavevmode\label{proof:25}
For \(0\le i<\ell-d\), subtracting the recurrences (6) gives
\[
bD_{i+1}=(aQ_{u+i+d}+c_{u+i+d})-(aQ_{u+i}+c_{u+i})=aD_i.
\]
Induction
gives
\[
b^kD_k=a^kD_0
\]
for \(0\le k\le\ell-d\), and
\(\gcd(a^k,b^k)=1\) gives (22).

% Markdown AST block 175: Para
If \(D_0>0\), the positive divisor
\(b^{\ell-d}\) is at most \(D_0\), so \(\ell-d\le\nu_b(D_0)\).

% Markdown AST block 176: Para
Finally
\[
Q_{u+d}\le x_{u+d}=r^dx_u<r^d(Q_u+1).
\]
The integer \(b\) need not be
prime. \qedhere
\end{proof}

% Markdown AST block 177: Lemma 26.
\begin{samepage}
\begin{lemma}\leavevmode\label{thm:26}
If \(x_n\ge b/(a-b)\), then \(Q_{n+1}\ge Q_n+1\). In
particular, if \(Q_u\ge b/(a-b)\), then \(Q_n\) is strictly increasing
for \(n\ge u\), and
\[
D_0=Q_{u+d}-Q_u\ge d\ge1
\]
in Lemma 25.
\label{thm-end:26}
\end{lemma}
\end{samepage}

% Markdown AST block 178: Proof.
\begin{proof}\leavevmode\label{proof:26}
\[
rx_n-x_n=x_n(a-b)/b\ge1,
\]
so
\[
\lfloor rx_n\rfloor\ge\lfloor x_n+1\rfloor=Q_n+1.\qedhere
\]
\end{proof}

% Markdown AST block 179: Para
For \(7/5\), \(b/(a-b)=5/2\), so \(Q_u\ge3\) suffices. A positive
difference must not be assumed when the integer parts stay at \(0\) or
\(1\).

% Markdown AST block 180: 6.2 From a rank certificate to a waiting time
\Needspace{6\baselineskip}
\subsection{From a rank certificate to a waiting time}\label{sec:6.2}

% Markdown AST block 181: Para
Let \(H\) be a finite labelled graph with edges of length one, together
with a partition of \(V(H)\) into \emph{blocks} and the following data.

% Markdown AST block 182: BulletList
\begin{itemize}
\tightlist
\item
  (Q-rank) An integer \(\rho(B)\in[1,\rho_{\max}]\) for every block
  \(B\) such that every edge \(u\to v\) between distinct blocks satisfies
  \[
  \rho(B(v))\ge\rho(B(u))+1.
  \]
  Put \(R=\rho_{\max}-1\).
\item
  (Q-block) Every block either has no internal edge, or is \emph{periodic}:
  there are \(d_B\in[1,D]\), \(h_B:B\to\mathbb Z\), and
  \(\lambda_B:\mathbb Z/d_B\to\mathcal C\) such that every internal edge
  \(u\xrightarrow{c}v\) satisfies
  \[
  \begin{aligned}
  &h_B(v)\equiv h_B(u)+1\pmod{d_B}\qquad\text{and}\\
  &c=\lambda_B(h_B(u)\bmod d_B).
  \end{aligned}
  \]
\item
  (Q-cyc) An integer \(\eta(B)\in[0,m]\) for every block with
  \[
  \begin{aligned}
  &\eta(B)\ge\mathbf 1[B\text{ periodic}],\qquad\text{and}\\
  &\eta(B(v))\ge\eta(B(u))+\mathbf 1[B(v)\text{ periodic}]
  \end{aligned}
  \]
  for every
  edge between distinct blocks.
\end{itemize}

% Markdown AST block 183: Para
No block with internal edges may remain uncertified; blocks need not be
SCCs. An SCC decomposition supplies such data through its acyclic
quotient, with \(\rho\) the length of a longest block sequence ending at
the block and \(\eta\) the largest number of periodic blocks on such a
sequence.

% Markdown AST block 184: Para
By (Q-rank), every finite walk in \(H\) has at most \(R\) edges between
distinct blocks, visits each block during one contiguous interval, and
by (Q-cyc) visits at most \(m\) periodic blocks.

% Markdown AST block 185: Para
An \emph{orbit segment of length \(T\) lies in \(H\)} if the states
\(\sigma(0),\ldots,\sigma(T)\) of the orbit (in the coordinates of
\(H\)) are vertices of \(H\) and \(\sigma(n)\xrightarrow{c_n}\sigma(n+1)\)
is an edge of \(H\) for \(0\le n<T\).

% Markdown AST block 186: Theorem 27 (waiting time from a rank certificate).
\begin{samepage}
\begin{theorem}[waiting time from a rank certificate]\leavevmode\label{thm:27}
Let
\[
\begin{aligned}
&\kappa=\log_br>0,\qquad \gamma\ge1+\kappa,\\
&P=Q_0\ge b/(a-b),\qquad\text{and}\qquad A\ge\log_b(P+2).
\end{aligned}
\]
If an orbit segment of length \(T\) lies in \(H\),
then
\[
T\le R\gamma^m+(A+\gamma D)\sum_{j=0}^{m-1}\gamma^j .
\tag{24}\label{eq:24}
\]
For \(m=0\) the sum is empty and \(T\le R\).
\label{thm-end:27}
\end{theorem}
\end{samepage}

% Markdown AST block 187: Proof.
\begin{proof}\leavevmode\label{proof:27}
\emph{Step 1: stay in one periodic block.} Suppose the segment
enters a periodic block \(B\) at time \(u\) and the states at times
\(u,\ldots,u+\ell\) lie in \(B\), so the edges at times
\(u,\ldots,u+\ell-1\) are internal. Induction with (Q-block) gives
\[
c_{u+i}=\lambda_B((h_B(\sigma(u))+i)\bmod d_B),
\]
hence
\[
c_{u+i+d_B}=c_{u+i}
\]
for \(0\le i<\ell-d_B\).

% Markdown AST block 188: Para
If \(\ell\ge d_B\),
Lemma 25 applies, and Lemma 26 (with \(Q_u\ge P\ge b/(a-b)\)) gives
\(D_0>0\); so
\[
\begin{aligned}
&b^{\ell-d_B}\le D_0<r^{d_B}(Q_u+1)\qquad\text{and}\\
&\ell-d_B<\log_b(Q_u+1)+d_B\kappa.
\end{aligned}
\]

% Markdown AST block 189: Para
Since
\[
Q_u\le x_u=r^u\xi<r^u(P+1),
\]
we have
\[
Q_u+1<r^u(P+1)+1\le r^u(P+2),
\]
so
\[
\log_b(Q_u+1)<u\kappa+A.
\]
Together,
\[
\ell<A+\kappa u+(1+\kappa)d_B\le A+(\gamma-1)u+\gamma D.
\]

% Markdown AST block 190: Para
If \(\ell<d_B\) the same bound holds trivially, since
\[
\ell<d_B\le D\le\gamma D.
\]
Thus the exit time (or the end of the
segment) satisfies
\[
u+\ell<\gamma u+A+\gamma D.
\]

% Markdown AST block 191: Para
\emph{Step 2: bookkeeping.} Let \(B_0,\ldots,B_s\) (\(s\le R\)) be the
blocks visited, with entry times \(u_0=0<u_1<\cdots<u_s\), let
\(\ell_k\) be the number of internal edges traversed in \(B_k\), so that
\[
u_{k+1}=u_k+\ell_k+1\qquad \text{and}\qquad T=u_s+\ell_s;
\]
edgeless blocks have
\(\ell_k=0\).

% Markdown AST block 192: Para
Let \(k_1<\cdots<k_{m'}\) (\(m'\le m\)) be the indices of
the periodic blocks, \(k_0=0\), \(x_0=0\),
\[
x_j=u_{k_j}+\ell_{k_j}
\]
the exit time from the \(j\)-th periodic block,
\[
\begin{aligned}
&\delta_j=k_j-k_{j-1}\qquad \text{for}\qquad 1\le j\le m',\qquad\text{and}\\
&\delta_{m'+1}=s-k_{m'}.
\end{aligned}
\]
Between periodic blocks only edgeless blocks
occur, so time advances by one per inter-block edge:
\[
\begin{aligned}
&u_{k_j}=x_{j-1}+\delta_j,\qquad T=x_{m'}+\delta_{m'+1},\qquad\text{and}\\
&\sum_{j=1}^{m'+1}\delta_j=s\le R.
\end{aligned}
\]

% Markdown AST block 193: Para
\emph{Step 3: expansion.} By Step 1,
\[
x_j<\gamma u_{k_j}+A+\gamma D=\gamma(x_{j-1}+\delta_j)+A+\gamma D.
\]
Induction gives
\[
x_{m'}<\sum_{j=1}^{m'}\gamma^{m'-j+1}\delta_j+(A+\gamma D)\sum_{j=0}^{m'-1}\gamma^j .
\]
Since \(\gamma\ge1\), \(\gamma^{m'-j+1}\le\gamma^{m'}\) and
\(1\le\gamma^{m'}\), so
\[
\begin{aligned}
T&=x_{m'}+\delta_{m'+1}\\
&<\gamma^{m'}\sum_{j=1}^{m'+1}\delta_j+(A+\gamma D)\sum_{j<m'}\gamma^j\\
&\le R\gamma^{m}+(A+\gamma D)\sum_{j<m}\gamma^j .
\end{aligned}
\]
If \(m'=0\) then
\[
T=s\le R.\qedhere
\]
\end{proof}

% Markdown AST block 194: Para
Each inter-block edge is amplified by a factor of at most \(\gamma^m\),
once for every periodic block visited afterwards. The proof handles
finite walks only and uses neither infinite walks nor samples of
\(\xi\). If any cyclic block is uncertified the theorem does not apply,
since a walk may stay in it forever.

% BEGIN PUBLIC ADDITION figure:rank-certificate
The ranks count blocks and periodic blocks along paths, whereas the period records the output inside a block.

\begin{figure}[htbp]
\centering
\begin{tikzpicture}[paperfig]
\node[ftitle] at(6,1.85){\figtext{Illustrative condensation DAG (not the full computed graph)}{SCC凝縮DAGの概念例（実計算グラフ全体ではない）}};
\node[fbox] (a) at(0,0){$A$\\$\rho=1,\ \eta=0$};
\node[fbox,double] (b) at(4,0.6){$B$\\$\rho=2,\ \eta=1$\\$d_B=2$};
\node[fbox] (c) at(4,-1.3){$C$\\$\rho=2,\ \eta=0$};
\node[fbox,double] (d) at(8,0.6){$D$\\$\rho=3,\ \eta=2$\\$d_D=3$};
\node[fbox] (e) at(12,0){$E$\\$\rho=4,\ \eta=2$};
\draw[fedge] (a)--(b);\draw[fedge] (a)--(c);\draw[fedge] (b)--(d);\draw[fedge] (c)--(d);\draw[fedge] (d)--(e);
\draw[fedge,dashed] (c) to[bend right=12](e);
\node[align=center,text width=12.7cm] at(6,-2.55)
 {\figtext{Double border: periodic block. Single border: no internal edge.}{二重枠：周期ブロック。一重枠：内部辺なし。}\\[4pt]
 $\rho$: \figtext{maximum number of blocks}{最大ブロック数}\qquad
 $\eta$: \figtext{maximum number of periodic blocks}{最大周期ブロック数}\\[4pt]
 $d_B$: \figtext{internal output period}{内部出力の周期}};
\node[fbox,text width=12cm] at(6,-4.3)
 {\figtext{Computed $7/5$ certificate (Table 7)}{実計算の $7/5$ 証明書（表7）}\\[5pt]
 $491\,275$ \figtext{blocks}{ブロック},\quad $35$ \figtext{cyclic blocks}{巡回ブロック}\\[4pt]
 $\max\rho=65,\quad\max\eta=6,\quad\max d_B=13$};
\end{tikzpicture}

\caption{An illustrative condensation DAG distinguishes three quantities: \(\rho\) is the maximum number of blocks on a path ending at a block, \(\eta\) is the maximum number of periodic blocks on such a path, and \(d_B\) is an internal output period. The five-block diagram is schematic. The separate box records the actual \(7/5\) certificate values from Table 7; those values do not describe the five-block example.}
\label{fig:rank-certificate}
\end{figure}
% END PUBLIC ADDITION figure:rank-certificate

% Markdown AST block 195: 6.3 The certificate for 7/5
\Needspace{6\baselineskip}
\subsection{The certificate for 7/5}\label{sec:6.3}

% Markdown AST block 196: Lemma 28 (reduction of the modulus).
\begin{samepage}
\begin{lemma}[reduction of the modulus]\leavevmode\label{thm:28}
Suppose \(\gcd(Q_n,4290)=1\)
for \(0\le n\le T\). Then \(c_n\in\mathcal C=\{-4,-2,2,4,6\}\) for
\(0\le n<T\), and with \(M=429=3\cdot11\cdot13\) and \(K=2048\), the
orbit segment of length \(T\) lies in the graph
\[
H=G(429,2048)\ \text{restricted to labels in }\mathcal C .
\]
\label{thm-end:28}
\end{lemma}
\end{samepage}

% Markdown AST block 197: Proof.
\begin{proof}\leavevmode\label{proof:28}
Lemma 23 gives the alphabet. Since \(\gcd(Q_n,429)=1\), the
states \((Q_n\bmod429,\lfloor2048\theta_n\rfloor)\) lie in
\(V(429,2048)\), and Proposition 5 makes consecutive states adjacent with
label \(c_n\). \qedhere
\end{proof}

% Markdown AST block 198: Para
Because \(\gcd(5,429)=1\), the congruence (E2) has the unique solution
\[
z\equiv5^{-1}(7q+c)\pmod{429};
\]
this is a consequence of reducing the
modulus from \(4290\) to \(429\), where \(5\) is a non-unit. The
avoidance of \(2\) and \(5\) is carried by the alphabet, and the prime
\(7\) is not assumed.

% Markdown AST block 199: Para
The graph \(H\) has
\[
\varphi(429)\cdot2048=240\cdot2048=491520
\]
vertices. Its SCC decomposition and the ranks \(\rho,\eta\) computed on
the acyclic quotient give the data of Table 7. Both implementations of
Section 9.3 enumerate all vertices and edges of \(H\), agree on the
block partition, on the attributes of every block, and on the phase
labels of every periodic block, and verify the inequalities of
(Q-rank), (Q-block), (Q-cyc) on every edge.

% Markdown AST block 200: Table 7.
% Caption is rendered at the start of the immediately following table.

% Markdown AST block 201: Table 7
\begingroup
\small
\setlength{\tabcolsep}{4pt}
\renewcommand{\arraystretch}{1.15}
\setlength{\LTcapwidth}{\textwidth}
\begin{longtable}{@{}>{\raggedright\arraybackslash}p{\dimexpr0.65\dimexpr\linewidth-8pt\relax} >{\raggedright\arraybackslash}p{\dimexpr0.35\dimexpr\linewidth-8pt\relax}@{}}
\caption{The rank certificate for \(7/5\) at \(M=429\),
\(K=2048\).}\label{tab:7}\\
\toprule
Quantity & Value \\
\midrule
\endfirsthead
\toprule
Quantity & Value \\
\midrule
\endhead
\bottomrule
\endlastfoot
Legal vertices & 491520 \\
Labelled edges & 891990 \\
Blocks (SCCs) & 491275 \\
Cyclic blocks & 35 \\
Certified cyclic blocks & 35 \\
Uncertified cyclic blocks & 0 \\
\(\max\rho\) & 65 \\
\(\max\eta\) & 6 \\
\(\max d_B\) & 13 \\
\end{longtable}
\endgroup

% Markdown AST block 202: Para
Thus (Q-rank) holds with \(\rho_{\max}=65\), \(R=64\); (Q-cyc) with
\(m=6\); (Q-block) with \(D=13\). Since \(7^4=2401<3125=5^5\), we have
\((7/5)^4<5\), \(\kappa=\log_5(7/5)<1/4\), and \(\gamma=5/4\) is
admissible. With
\[
A=L(P)=\min\{h\ge0:5^h\ge P+2\}\ge\log_5(P+2),
\]
the
constants in (24) are, exactly in rational arithmetic,
\[
\begin{aligned}
\sum_{j=0}^{5}\Bigl(\tfrac54\Bigr)^j&=\frac{11529}{1024},\\
64\Bigl(\tfrac54\Bigr)^6+\tfrac54\cdot13\cdot\frac{11529}{1024}
&=\frac{15625}{64}+\frac{749385}{4096}\\
&=\frac{1749385}{4096},
\end{aligned}
\]
\[
\text{and}\quad
1749385<428\cdot4096=1753088,\qquad 11529<12\cdot1024=12288.
\]
Hence, for an orbit segment of length \(T\) lying in \(H\) with
\(P\ge3\),
\[
T\le\frac{1749385}{4096}+\frac{11529}{1024}\,L(P)<428+12L(P).
\tag{25}\label{eq:25}
\]

% Markdown AST block 203: Proof of Theorem 3.
\begin{proof}[Proof of Theorem 3]\leavevmode\label{proof:3}
Let \(P=\lfloor\xi\rfloor>13\) and
\(T_0=428+12L(P)\). If \(\gcd(Q_n,4290)=1\) for all
\(0\le n\le T_0\), then by Lemma 28 the orbit segment of length \(T_0\)
lies in \(H\), and Theorem 27 with the certificate of Table 7 (Lemma 26
applies since \(P>13\ge3\)) gives
\[
T_0<428+12L(P)=T_0,
\]
a
contradiction.

% Markdown AST block 204: Para
Hence some \(0\le n\le T_0\) has
\(\gcd(Q_n,4290)>1\).

% Markdown AST block 205: Para
For such \(n\), a prime \(p\in\{2,3,5,11,13\}\)
divides
\[
Q_n\ge Q_0=P>13\ge p,
\]
so \(Q_n=pv\) with \(v>1\) is
composite.

% Markdown AST block 206: Para
For the last assertion, apply the result to
\(\xi'=\xi(7/5)^N\), for which
\[
\lfloor\xi'(7/5)^n\rfloor=Q_{N+n}
\]
and
\(\lfloor\xi'\rfloor=Q_N\); the graph \(H\) does not depend on \(\xi\),
and \(P\) enters only through the bound \(Q_u<r^u(P+1)\). \qedhere
\end{proof}

% Markdown AST block 207: Para
The bound (25) itself holds for \(P\ge3\); the hypothesis \(P>13\) is
used only to conclude compositeness. The integer definition of \(L(P)\)
avoids rounding questions; it equals \(\lceil\log_5(P+2)\rceil\) for
every \(P\ge0\).

% Markdown AST block 208: 7. Scope of the certificate method
\Needspace{6\baselineskip}
\section{Scope of the certificate method}\label{sec:7}

% Markdown AST block 209: Para
This section and the next study the reach of the procedure of
Sections 3--4 when its operations are fixed precisely. Throughout,
\[
\Sigma=\{1-b,\ldots,a-1\},
\]
and \(\operatorname{rad}M\) is the product
of the distinct primes dividing \(M\).

% Markdown AST block 210: 7.1 The general carry alphabet
\Needspace{6\baselineskip}
\subsection{The general carry alphabet}\label{sec:7.1}

% Markdown AST block 211: Para
For a reduced base \(a/b\) define
\[
\mathcal C(a,b)=\{c\in\mathbb Z:1-b\le c\le a-1,\ c\equiv b-a\pmod2,\
\gcd(c,ab)=1\}.
\tag{26}\label{eq:26}
\]

% Markdown AST block 212: Lemma 29.
\begin{samepage}
\begin{lemma}\leavevmode\label{thm:29}
Let \(\operatorname{rad}(2ab)\mid M\). Every edge
\((q,j)\xrightarrow{e}(z,k)\) of \(G(M,K)\) has \(e\in\mathcal C(a,b)\).
In particular \(c_n\in\mathcal C(a,b)\) whenever
\(\gcd(Q_nQ_{n+1},2ab)=1\). Conversely, for such \(M\), the label
conditions imposed by the primes dividing \(2ab\) are exactly
\(e\in\mathcal C(a,b)\).
\label{thm-end:29}
\end{lemma}
\end{samepage}

% Markdown AST block 213: Proof.
\begin{proof}\leavevmode\label{proof:29}
Let \(p\mid a\); then \(p\nmid b\) and (E2) gives
\[
bz\equiv e\pmod p,
\]
so \(p\mid e\) if and only if \(p\mid z\), which
is excluded as \(z\) is a unit.

% Markdown AST block 214: Para
Let \(p\mid b\); then
\[
0\equiv aq+e\pmod p,
\]
so \(p\mid e\) if and only if \(p\mid q\),
excluded.

% Markdown AST block 215: Para
If \(a,b\) are odd and \(2\mid M\), then
\[
z\equiv q+e\pmod2
\]
with \(q,z\) odd, so \(e\) is even, that is, \(e\equiv b-a\pmod2\); if
one of \(a,b\) is even the parity condition \(e\equiv b-a\pmod2\) is
the same as \(\gcd(e,ab)\) being odd, which was shown. The converse is
the same computation read backwards. \qedhere
\end{proof}

% Markdown AST block 216: Para
For \(7/5\),
\[
\mathcal C(7,5)=\{-4,-2,2,4,6\},
\]
and for \(5/2\),
\[
\mathcal C(5,2)=\{-1,1,3\}:
\]
the alphabets of Lemmas 23 and 24 are
exactly the general ones. When \(\mathcal C(a,b)\) is used to express
avoidance of all primes dividing \(2ab\), the conclusion of a certificate
is a divisibility statement whose prime set contains all primes dividing
\(2ab\); the individual reductions of Section 5 avoid some of these
primes (the prime \(7\) for \(7/5\), the prime \(5\) for \(5/2\)).

% Markdown AST block 217: 7.2 The procedure and its three soundness notions
\Needspace{6\baselineskip}
\subsection{The procedure and its three soundness notions}\label{sec:7.2}

% Markdown AST block 218: Para
\textbf{Word-labelled graphs with expansions.} In this section a word-labelled
graph has vertex set contained in some \(V(M,K)\), and each edge
\(u\xrightarrow{w}v\) records an \emph{expansion path}
\[
u=u_0\xrightarrow{w_0}u_1\cdots\xrightarrow{w_{|w|-1}}u_{|w|}=v
\]
of
edges of \(G(M,K)\). One-step graphs are the case of length one.
For \(M\mid M'\) and \(K'=fK\), the \emph{projection} is
\[
\pi(q,J)=(q\bmod M,\lfloor J/f\rfloor).
\]

% Markdown AST block 219: Para
\textbf{Pseudo-orbits.} A \emph{right-directed \(M\)-pseudo-orbit} is a sequence
\((q_n,\theta_n,c_n)_{n\ge n_0}\) with \(q_n\in(\mathbb Z/M)^\times\),
\(\theta_n\in[0,1)\), \(c_n\in\mathbb Z\), and
\[
a\theta_n-b\theta_{n+1}=c_n\ \text{(exactly, in }\mathbb R),\qquad
bq_{n+1}\equiv aq_n+c_n\pmod M .
\]
A \emph{left-directed} pseudo-orbit satisfies the same conditions for
\(n\le n_0\). A true orbit whose tail avoids the primes of \(M\) gives a
right-directed pseudo-orbit with \(q_n=Q_n\bmod M\); a pseudo-orbit gives,
for every \(K\), a walk \((q_n,\lfloor K\theta_n\rfloor)\) of
\(G(M,K)\), by the proof of Proposition 5, which uses only the real
identity and the half-open cells. A pseudo-orbit modulo \(M\) reduces to
one modulo any divisor of \(M\). What a true orbit has and a pseudo-orbit
lacks is a single integer sequence \(Q_n\) linking residues and fractional
parts.

% Markdown AST block 220: Para
\textbf{Operations.} An \emph{execution} is a finite sequence
\(\Gamma_0,\ldots,\Gamma_s\) of word-labelled graphs with parameters
\((M_t,K_t)\), \(M_t\mid M_{t+1}\), \(K_t\mid K_{t+1}\), built by the
following operations.

% Markdown AST block 221: BulletList
\begin{itemize}
\tightlist
\item
  \begin{enumerate}
  \def\labelenumi{(\Roman{enumi})}
  \tightlist
  \item
    \emph{Initial graph.} Either \(\Gamma_0=G(M_0,K_0)\) on all legal
    states, with all labels in \(\Sigma\); or, when the final modulus is to
    contain \(\operatorname{rad}(2ab)\), \(\Gamma_0=G'(R_0,K_0)\), the
    graph with vertices \((q,j)\), \(q\in(\mathbb Z/R_0)^\times\)
    (\(R_0=1\) allowed, \(\gcd(R_0,2ab)=1\)), edges labelled by
    \(e\in\mathcal C(a,b)\) with
    \[
    bz\equiv aq+e\pmod{R_0}
    \]
    and (E3); in
    this form we set \(M_0=\operatorname{rad}(2ab)R_0\), and the residues
    modulo the primes of \(2ab\) are represented by the alphabet
    (Lemma 29 and Proposition 36).
  \end{enumerate}
\item
  \begin{enumerate}
  \def\labelenumi{(\Alph{enumi})}
  \setcounter{enumi}{15}
  \tightlist
  \item
    \emph{Phase removal.} Replace \(\Gamma_t\) by the union of its
    uncertified cyclic SCCs with their internal edges. In each component,
    use the total word lengths of paths from a common root and apply (19)
    to every internal edge and letter (Section 4.2; Lemma 7 for length
    one). The outcome is independent of these choices by Corollary 22.
  \end{enumerate}
\item
  \begin{enumerate}
  \def\labelenumi{(\Alph{enumi})}
  \setcounter{enumi}{2}
  \tightlist
  \item
    \emph{Contraction.} The operation \(P\) of Section 4.3; expansion
    paths are concatenated.
  \end{enumerate}
\item
  \begin{enumerate}
  \def\labelenumi{(\Alph{enumi})}
  \setcounter{enumi}{17}
  \tightlist
  \item
    \emph{Cell refinement.} \(K_{t+1}=fK_t\); a vertex \((q,j)\) is
    replaced by \((q,fj+i)\), \(0\le i<f\), and each edge by all
    refinements of its expansion path that satisfy the edge conditions of
    \(G(M_t,K_{t+1})\) at every step.
  \end{enumerate}
\item
  \begin{enumerate}
  \def\labelenumi{(\Alph{enumi})}
  \setcounter{enumi}{11}
  \tightlist
  \item
    \emph{Modulus enlargement.} \(M_{t+1}=hM_t\); each edge is replaced by
    all sequences of residues modulo \(M_{t+1}\) along its expansion path
    that are units and satisfy (E2) at every step (all solutions, with no
    inverse of \(b\) assumed). Lifting by a prime \(p\nmid ab\)
    (Section 4.4) is the case \(h=p\), and in the second form of (I) the
    new primes do not divide \(2ab\).
  \end{enumerate}
\item
  \begin{enumerate}
  \def\labelenumi{(\Alph{enumi})}
  \tightlist
  \item
    \emph{Carry restriction.} If \(\operatorname{rad}(2ab)\mid M_t\),
    restrict labels to \(\mathcal C(a,b)\).
  \end{enumerate}
\end{itemize}

% Markdown AST block 222: Para
Let \(\mathfrak M_0\) be the procedure allowing exactly these operations.
An execution \emph{succeeds} if \(\Gamma_s=\varnothing\). The executions of
Section 5 are instances: Table 1 uses (I), (P), (R); Tables 4 and 6 use
the second form of (I) with \(R_0=1\), then (P), (C), (L).

% Markdown AST block 223: Para
\textbf{Soundness.} Let \(\omega\) be an infinite walk (right- or
left-directed) of \(G(M_*,K_*)\) with \(M_t\mid M_*\) and
\(K_t\mid K_*\), whose output is not eventually periodic (in the
direction of the walk). We say that \emph{a tail of \(\omega\) becomes a walk
of \(\Gamma_t\)} if there is an infinite walk of \(\Gamma_t\) whose
successive expansion paths coincide with the projection of the state
sequence of \(\omega\) to \((M_t,K_t)\), and whose output is a tail of
the carry sequence of \(\omega\). An operation is

% Markdown AST block 224: BulletList
\begin{itemize}
\tightlist
\item
  \emph{walk-sound} if it preserves this property for every such \(\omega\);
\item
  \emph{pseudo-orbit-sound} if it preserves it for every \(M\)-pseudo-orbit
  with non-eventually-periodic carries, \(M_t\mid M\), in place of
  \(\omega\);
\item
  \emph{orbit-sound} if it preserves it for every true orbit tail avoiding the
  primes of the final modulus.
\end{itemize}

% Markdown AST block 225: Para
Walk-soundness implies pseudo-orbit-soundness, which implies
orbit-soundness, by the inclusions above. Orbit-soundness of the
operations is the content of Theorems 10 and 19.

% Markdown AST block 226: Lemma 30 (preservation).
\begin{samepage}
\begin{lemma}[preservation]\leavevmode\label{thm:30}
The operations (P), (C), (R), (L), (A) are
walk-sound, for right- and for left-directed walks.
\label{thm-end:30}
\end{lemma}
\end{samepage}

% Markdown AST block 227: Proof.
\begin{proof}\leavevmode\label{proof:30}
(R) By Lemma 31 below, the projection of \(\omega\) to
\((M_t,K_{t+1})\) is a walk of \(G(M_t,K_{t+1})\) whose further
projection to \((M_t,K_t)\) is the current expansion path; so it is one of
the refined edges.

% Markdown AST block 228: Para
\textbf{(L)} Likewise, the residues of \(\omega\) modulo
\(M_{t+1}\) are units satisfying (E2), hence among all solutions.

% Markdown AST block 229: Para
\textbf{(A)} By
Lemma 29, every label of \(\omega\) lies in \(\mathcal C(a,b)\) once
\(\operatorname{rad}(2ab)\mid M_t\mid M_*\).

% Markdown AST block 230: Para
\textbf{(P)} A right-directed infinite
walk of \(\Gamma_t\) eventually stays in one cyclic SCC \(C\)
(Lemma 6). For a left-directed walk, the set of times spent in each SCC
is an interval, since an SCC cannot be re-entered; finitely many
intervals cover \(\mathbb Z_{\le n_0}\), so the leftmost is unbounded
below, and that SCC contains an internal edge. If \(C\) were certified,
the output of the tail would be periodic: for right-directed walks by
Lemma 12(b); for left-directed walks, the congruence of Lemma 12(a)
propagated backwards shows that the letter at time \(t\le-1\) (relative
to the last vertex \(v_0\)) equals \(\lambda((h(v_0)+t)\bmod d)\). This
contradicts the non-periodicity of the tail of \(\omega\). Hence \(C\)
is retained and the tail is a walk of the retained subgraph.

% Markdown AST block 231: Para
\textbf{(C)} For
right-directed walks this is Proposition 16. For a left-directed walk, if
it avoided the macro vertex set \(B\) from some point to the left, all
its vertices there would have out-degree one; by finiteness some vertex
repeats, \(v_m=v_{m'}\) with \(m<m'\), and \(Z=\{v_m,\ldots,v_{m'-1}\}\)
is a cycle of out-degree-one vertices closed under successors. Its SCC
is then \(Z\) itself, all of whose vertices have out-degree one, so its
anchor lies in \(B\cap Z\), contradicting avoidance of \(B\). Hence the
walk visits \(B\) infinitely often to the left, and its segments between
visits are macro edges, as in Proposition 16. \qedhere
\end{proof}

% Markdown AST block 232: Para
Any further operation that is pseudo-orbit-sound in both directions may
be adjoined to \(\mathfrak M_0\) without affecting the obstruction
theorems of Section 8. An example, the whole-word interval test, is treated in Appendix E.4
(Proposition 73); it is pseudo-orbit-sound but not walk-sound, and it is
excluded from Theorem 33.

% Markdown AST block 233: 7.3 Projection and inheritance of certification
\Needspace{6\baselineskip}
\subsection{Projection and inheritance of certification}\label{sec:7.3}

% Markdown AST block 234: Lemma 31 (projection).
\begin{samepage}
\begin{lemma}[projection]\leavevmode\label{thm:31}
Let \(M\mid M'\) and \(K'=fK\). The
projection \(\pi\) maps every edge \((q,J)\xrightarrow{e}(z,H)\) of
\(G(M',K')\) to an edge of \(G(M,K)\) with the same label.
\label{thm-end:31}
\end{lemma}
\end{samepage}

% Markdown AST block 235: Proof.
\begin{proof}\leavevmode\label{proof:31}
(E1) is unchanged and (E2) reduces modulo \(M\). For (E3),
write \(J=fj+i\), \(H=fk+\ell\) with \(0\le i,\ell<f\). Then
\(eK'=efK\)
and
\[
\begin{aligned}
f\bigl(aj-b(k+1)\bigr)&\le afj+ai-bfk-b(\ell+1)\\
&<efK<afj+a(i+1)-bfk-b\ell\\
&\le f\bigl(a(j+1)-bk\bigr),
\end{aligned}
\]
the outer inequalities using \(i\ge0\), \(\ell+1\le f\) and
\(i+1\le f\), \(\ell\ge0\). Dividing by \(f\) gives (E3) for
\((j,k)\). \qedhere
\end{proof}

% Markdown AST block 236: Lemma 32 (inheritance).
\begin{samepage}
\begin{lemma}[inheritance]\leavevmode\label{thm:32}
In the setting of Lemma 31, let
\(W\subseteq V(M,K)\) and \(W'\subseteq V(M',K')\) with
\(\pi(W')\subseteq W\), and let \(C'\) be a cyclic SCC of \(G_{K'}(W')\)
with root \(\rho'\) and root path lengths \(h'\). Then \(\pi(C')\) lies
in one cyclic SCC \(C\) of \(G_K(W)\), with root path lengths \(h\) and
gcd \(d\); \(d\mid d'\); and if \(C\) is certified then so is \(C'\).
\label{thm-end:32}
\end{lemma}
\end{samepage}

% Markdown AST block 237: Proof.
\begin{proof}\leavevmode\label{proof:32}
By Lemma 31 and \(\pi(W')\subseteq W\), the projection of
every internal edge of \(C'\) is an edge of the induced subgraph
\(G_K(W)\); hence the image of \(C'\) is strongly connected in
\(G_K(W)\) and contains an internal edge, so it lies in a cyclic SCC
\(C\).

% Markdown AST block 238: Para
Projecting the chosen
path from \(\rho'\) to \(x\) and using (15) along it gives
\[
h(\pi x)\equiv h(\pi\rho')+h'(x)\pmod d.
\]

% Markdown AST block 239: Para
For an internal edge
\(u'\to v'\) of \(C'\),
\[
h'(u')+1-h'(v')\equiv h(\pi u')+1-h(\pi v')
\equiv0\pmod d,
\]
so \(d\mid d'\).

% Markdown AST block 240: Para
If \(C\) is certified, the label of
\(u'\to v'\) is
\[
\lambda_{h(\pi u')\bmod d}
=\lambda_{(h(\pi\rho')+h'(u'))\bmod d},
\]
which depends only on
\(h'(u')\bmod d'\); hence \(C'\) is certified. \qedhere
\end{proof}

% Markdown AST block 241: 7.4 Normal form of finite certificates
\Needspace{6\baselineskip}
\subsection{Normal form of finite certificates}\label{sec:7.4}

% Markdown AST block 242: Theorem 33 (contraction does not enlarge the class).
\begin{samepage}
\begin{theorem}[contraction does not enlarge the class]\leavevmode\label{thm:33}
Let
\(\Gamma_0,\ldots,\Gamma_s\) be a successful execution of
\(\mathfrak M_0\), and let \(M_*=M_s\) (with
\(\operatorname{rad}(2ab)\mid M_*\) in the second form of (I)) and
\(K_*=K_s\). Then every cyclic SCC of the complete one-step graph
\(G(M_*,K_*)\) passes the phase test of Section 3.2. Conversely, a
one-step certificate as in Theorem 10 (or Lemma 11) is an execution of
\(\mathfrak M_0\) without (C).
\label{thm-end:33}
\end{theorem}
\end{samepage}

% Markdown AST block 243: Proof.
\begin{proof}\leavevmode\label{proof:33}
Suppose \(G(M_*,K_*)\) has an infinite walk \(\omega\) with
non-eventually-periodic output. Its projection to \((M_0,K_0)\) is a walk
of \(\Gamma_0\): in the first form of (I) by Lemma 31; in the second
form, its labels lie in \(\mathcal C(a,b)\) by Lemma 29 and its residues
modulo \(R_0\) satisfy (E2).

% Markdown AST block 244: Para
By Lemma 30, inductively, a tail of
\(\omega\) becomes a walk of every \(\Gamma_t\), including
\(\Gamma_s=\varnothing\): a contradiction.

% Markdown AST block 245: Para
Hence every infinite walk of
\(G(M_*,K_*)\) has eventually periodic output, in particular every
infinite walk inside a cyclic SCC, and Theorem 21 ((b)\(\Rightarrow\)(a)) shows that
every cyclic SCC is certified. The converse is clear. \qedhere
\end{proof}

% Markdown AST block 246: Para
Consequently, for a fixed \(K\) and a fixed set of primes, the order in
which primes are added does not affect whether an execution can succeed;
it affects only the intermediate sizes. The conclusion ``no walk of
\(G(M_*,K_*)\) has non-eventually-periodic output'' does not use
Theorem 21; only its restatement as a phase test does.

% Markdown AST block 247: Theorem 34 (projection monotonicity).
\begin{samepage}
\begin{theorem}[projection monotonicity]\leavevmode\label{thm:34}
Let a one-step execution with
operations (P), (R), (L) start from \(W_0=V(M_0,K_0)\), with
\(U_t=R_{K_t}(W_t)\) and \(W_{t+1}\) the full enlargement of \(U_t\)
(Lemmas 9 and 11), and suppose \(U_s=\varnothing\). Then
\[
R_{K_s}(V(M_s,K_s))=\varnothing:
\]
the complete graph at the final
parameters certifies in one stage. Consequently, if \(M_s\mid M'\) and
\(K_s\mid K'\), then
\[
R_{K'}(V(M',K'))=\varnothing
\]
as well; in
particular the pairs \((s!,s!)\), \(s=1,2,\ldots\), miss no one-step
certificate. The certificate at \((M',K')\) proves only
\(\gcd(Q_n,M')>1\) infinitely often; it does not recover the conclusion
for the original prime set.
\label{thm-end:34}
\end{theorem}
\end{samepage}

% Markdown AST block 248: Proof.
\begin{proof}\leavevmode\label{proof:34}
We show
\[
\pi_t(R_{K_s}(V(M_s,K_s)))\subseteq U_t
\]
by
induction on \(t\), where \(\pi_t\) projects to \((M_t,K_t)\). For
\(t=0\), Lemma 32 applies since \(W_0\) is the full vertex set.

% Markdown AST block 249: Para
For the
step, let \(v\) lie in an uncertified cyclic SCC \(C'\) of the final
graph. By induction
\[
\pi_t(C')\subseteq U_t,
\]
hence
\[
\pi_{t+1}(C')\subseteq W_{t+1}.
\]
The set \(\pi_{t+1}(C')\) is
strongly connected inside \(G_{K_{t+1}}(W_{t+1})\) (an induced subgraph
keeps the edges) and lies in one cyclic SCC \(C\); if \(C\) were
certified, Lemma 32 would certify \(C'\). So \(C\subseteq U_{t+1}\) and
\(\pi_{t+1}(v)\in U_{t+1}\).

% Markdown AST block 250: Para
At \(t=s\), \(\pi_s\) is the identity.

% Markdown AST block 251: Para
The last assertion follows from
Lemma 32 applied to the projection from \((M',K')\) to \((M_s,K_s)\) with
\(W=V(M_s,K_s)\), \(W'=V(M',K')\). \qedhere
\end{proof}

% Markdown AST block 252: Para
This is relative completeness only: a search over the pairs \((s!,s!)\)
finds a certificate whenever one exists at some parameters, but with the
weaker conclusion \(\gcd(Q_n,s!)>1\), and no termination is asserted.

% Markdown AST block 253: Proposition 35 (prime powers).
\begin{samepage}
\begin{proposition}[prime powers]\leavevmode\label{thm:35}
For every \(M\) and \(K\), the sets
of infinite output words (together with their cell sequences) of
\(G(M,K)\) and of \(G(\operatorname{rad}M,K)\) coincide.
\label{thm-end:35}
\end{proposition}
\end{samepage}

% Markdown AST block 254: Proof.
\begin{proof}\leavevmode\label{proof:35}
One inclusion is Lemma 31. Conversely, let
\((\bar q_n,j_n)_{n\ge0}\) be an infinite walk of
\(G(\operatorname{rad}M,K)\) with output \((c_n)\). For each
\(p^e\parallel M\) we construct units \(q_n\) modulo \(p^e\) with
\[
q_n\equiv\bar q_n\pmod p\qquad \text{and}\qquad bq_{n+1}\equiv aq_n+c_n\pmod{p^e}.
\]

% Markdown AST block 255: Para
If \(p\nmid b\): take any lift \(q_0\) of \(\bar q_0\) and
\[
q_{n+1}=b^{-1}(aq_n+c_n)\bmod p^e;
\]
modulo \(p\) this is the same
recurrence with the same initial value, so \(q_n\equiv\bar q_n\).

% Markdown AST block 256: Para
If
\(p\mid b\) (so \(p\nmid a\)): put
\[
q_n=-\sum_{j=0}^{e-1}b^ja^{-j-1}c_{n+j}\bmod p^e.
\]
Then
\[
\begin{aligned}
bq_{n+1}-aq_n&=-\sum_{j=1}^{e}b^ja^{-j}c_{n+j}+\sum_{j=0}^{e-1}b^ja^{-j}c_{n+j}\\
&=c_n-b^ea^{-e}c_{n+e}\\
&\equiv c_n\pmod{p^e},
\end{aligned}
\]
and modulo \(p\),
\[
q_n\equiv-a^{-1}c_n\equiv\bar q_n
\]
because the original walk satisfies
\(0\equiv a\bar q_n+c_n\pmod p\).

% Markdown AST block 257: Para
Combine by the Chinese remainder
theorem; (E3) involves only cells and labels. \qedhere
\end{proof}

% Markdown AST block 258: Proposition 36 (primes dividing \(2ab\)).
\begin{samepage}
\begin{proposition}[primes dividing \(2ab\)]\leavevmode\label{thm:36}
Let
\(M_0=\operatorname{rad}(2ab)\) and \(\gcd(R,M_0)=1\). The sets of
infinite output words (with cell sequences) of \(G(M_0R,K)\) and of
\(G'(R,K)\) (defined in (I)) coincide.
\label{thm-end:36}
\end{proposition}
\end{samepage}

% Markdown AST block 259: Proof.
\begin{proof}\leavevmode\label{proof:36}
A walk of \(G(M_0R,K)\) has labels in \(\mathcal C(a,b)\) by
Lemma 29 and residues modulo \(R\) satisfying (E2); this gives a walk of
\(G'(R,K)\).

% Markdown AST block 260: Para
Conversely, given a walk \((\bar q_n,j_n)\) of \(G'(R,K)\)
with output \(c_n\in\mathcal C(a,b)\), define residues for each prime
\(p\mid M_0\): if \(p\mid a\), let \(q_0\) be any unit and
\[
q_{n+1}\equiv b^{-1}c_n
\]
(then \(bq_{n+1}\equiv c_n\equiv aq_n+c_n\));
if \(p\mid b\), let
\[
q_n\equiv-a^{-1}c_n
\]
(then \(bq_{n+1}-aq_n\equiv c_n\)); if \(p=2\) and \(a,b\) are odd, let
\(q_n=1\) (then \(b\equiv a+c_n\pmod2\) since \(c_n\) is even). All are
units since \(p\nmid c_n\). Combine with \(\bar q_n\) by the Chinese
remainder theorem. \qedhere
\end{proof}

% Markdown AST block 261: Para
Define the class of \emph{finitely certifiable bases}
\[
\begin{aligned}
\mathcal F_{\rm cert}&=\Bigl\{\tfrac ab:\ a>b\ge2,\ \gcd(a,b)=1,\\\
\exists K\ \exists S\text{ finite},\ p\nmid2ab\ (p&\in S):\\\
&\text{every cyclic SCC of }G'\bigl(\textstyle\prod_{p\in S}p,K\bigr)\\
&\text{ satisfies Theorem 21(c)}\Bigr\}.
\end{aligned}
\]

% Markdown AST block 262: Theorem 37 (normal form).
\begin{samepage}
\begin{theorem}[normal form]\leavevmode\label{thm:37}
For a reduced base \(a/b\) the following
are equivalent.

% Markdown AST block 263: OrderedList
\begin{enumerate}
\def\labelenumi{(\arabic{enumi})}
\item
  Some execution of \(\mathfrak M_0\) succeeds.
\item
  Some one-step execution (operations (I), (P), (R), (L), first form)
  succeeds for some modulus \(M\) and some sequence of cell counts.
\item
  There are \(K\) and a finite set \(S\) of primes not dividing
  \(2ab\) such that every cyclic SCC of \(G'(\prod_{p\in S}p,K)\) passes
  the phase test.
\item
  \(a/b\in\mathcal F_{\rm cert}\).
\end{enumerate}
\label{thm-end:37}
\end{theorem}
\end{samepage}

% Markdown AST block 264: Proof.
\begin{proof}\leavevmode\label{proof:37}
(3)\(\Leftrightarrow\)(4) is Theorem 21.

% Markdown AST block 265: Para
(3)\(\Rightarrow\)(2): by Proposition 36 and
Theorem 21, every cyclic SCC of \(G(M_0R,K)\) is certified, which is a
one-step execution with \(s=0\).

% Markdown AST block 266: Para
(2)\(\Rightarrow\)(1) is clear.

% Markdown AST block 267: Para
(1)\(\Rightarrow\)(3): by
Theorem 33 every cyclic SCC of \(G(M_*,K_*)\) is certified. By Lemma 31
and Theorem 21, the same holds for \(G(M_0M_*,K_*)\) with
\(M_0=\operatorname{rad}(2ab)\): every infinite walk of the finer graph
projects to one with eventually periodic output.

% Markdown AST block 268: Para
By Proposition 35 it
holds for
\[
G(\operatorname{rad}(M_0M_*),K_*)=G(M_0R,K_*),
\]
where \(R\)
is the product of the primes of \(M_*\) not dividing \(2ab\); by
Proposition 36 and Theorem 21 it holds for \(G'(R,K_*)\). \qedhere
\end{proof}

% Markdown AST block 269: Para
The equivalence concerns the \emph{existence} of some certificate, not the
prime set of the conclusion. A certificate in (2) proves
\(\gcd(Q_n,M)>1\) infinitely often; a certificate in (3) proves
\(\gcd(Q_n,2abR)>1\) infinitely often, a conclusion whose prime set
contains every prime dividing \(2ab\). For example, (1) uses
\(M=4290\), which does not contain \(7\), while the normal form for
\(7/5\) uses \(\{2,5,7\}\cup\{3,11,13\}\). Membership of
\(3/2,4/3,5/4,7/5,5/2\) in \(\mathcal F_{\rm cert}\) follows from the
certificates of Section 5 (Tables 1, 3, 4, 6); no other base is claimed
to belong to it, and no simple inequality or congruence in \(a,b\)
characterizes the class.

% Markdown AST block 270: 8. Obstructions and the external-prime budget
\Needspace{6\baselineskip}
\section{Obstructions and the external-prime budget}\label{sec:8}

% Markdown AST block 271: Para
The results of this section have the following common form: for a given
modulus \(M\), \emph{no execution of \(\mathfrak M_0\) all of whose moduli
\(M_t\) divide \(M\) succeeds}; the same holds when \(\mathfrak M_0\) is
extended by operations that are pseudo-orbit-sound in both directions.
They are statements about the non-termination of a procedure. They do
\emph{not} imply
\[
\exists\,\xi>0\ \exists\,n_0\ \forall n\ge n_0:\ \gcd(\lfloor\xi r^n\rfloor,M)=1,
\]
and they are not counterexamples to any statement about composite terms.
The witnesses are pseudo-orbits, which lack an integer sequence linking
residues and fractional parts; Lemma 41 shows that every finite window
of such a witness is realized by a true orbit, in the form
\(\forall N\ \exists\xi_N\), which cannot be exchanged for
\(\exists\xi\ \forall N\). Procedures that use information available only
for true orbits (integrality of \(Q_n\), forward realizability, or the
recurrence hypothesis of Section 10) are not covered by these theorems.

% Markdown AST block 272: Lemma 38 (type of the obstruction).
\begin{samepage}
\begin{lemma}[type of the obstruction]\leavevmode\label{thm:38}
Let \(M\ge2\). If there is a
right- or left-directed \(M\)-pseudo-orbit whose carry sequence is not
eventually periodic (in its direction), and whose carries lie in
\(\mathcal C(a,b)\) whenever the second form of (I) is used, then no
execution of \(\mathfrak M_0\) (or of an extension by pseudo-orbit-sound
operations) with all \(M_t\mid M\) succeeds.
\label{thm-end:38}
\end{lemma}
\end{samepage}

% Markdown AST block 273: Proof.
\begin{proof}\leavevmode\label{proof:38}
The pseudo-orbit yields, for every \(K\), a walk of
\(G(M,K)\) with the same carries; its projection to \((M_0,K_0)\) is a
walk of \(\Gamma_0\) (in the second form of (I), the labels lie in
\(\mathcal C(a,b)\) by hypothesis and the residues modulo \(R_0\) satisfy
(E2)). By Lemma 30, applied inductively, a tail of it becomes a walk of
every \(\Gamma_t\); hence \(\Gamma_s\ne\varnothing\). \qedhere
\end{proof}

% Markdown AST block 274: 8.1 Large numerators
\Needspace{6\baselineskip}
\subsection{Large numerators}\label{sec:8.1}

% Markdown AST block 275: Theorem 39.
\begin{samepage}
\begin{theorem}\leavevmode\label{thm:39}
Let \(M\ge2\), \(R=\operatorname{rad}M\), and
\[
a\ge2R.
\]
Then no execution of \(\mathfrak M_0\) (or of an extension by
pseudo-orbit-sound operations) with all \(M_t\mid M\) and \(\Gamma_0\) as
in (I) succeeds. Hence a successful execution requires
\[
a<2\operatorname{rad}(M).
\tag{27}\label{eq:27}
\]
\label{thm-end:39}
\end{theorem}
\end{samepage}

% Markdown AST block 276: Proof.
\begin{proof}\leavevmode\label{proof:39}
\emph{Step 1: a left-directed nonperiodic carry sequence.} Put
\(\theta_0=0\). Given \(\theta_{n+1}\in[0,1)\), choose
\[
\begin{aligned}
c_n&\in[-b\theta_{n+1},\,a-b\theta_{n+1})\cap(b-a+R\mathbb Z),\\
\theta_n&=\frac{b\theta_{n+1}+c_n}{a}.
\end{aligned}
\]
The half-open interval has length \(a\ge2R\), so it contains at least
\(\lfloor a/R\rfloor\ge2\) terms of the progression. Then
\[
0\le b\theta_{n+1}+c_n<a
\]
gives \(\theta_n\in[0,1)\), the identity
\[
a\theta_n-b\theta_{n+1}=c_n
\]
is exact, and \(-b<c_n<a\) gives
\(1-b\le c_n\le a-1\). Using only two ordered choices at each step, the
map \(\{0,1\}^{\mathbb N}\to\Sigma^{\mathbb Z_{\le0}}\) is injective
(at the first index where two choice sequences differ, the values
\(\theta_{n+1}\) agree and the \(c_n\) differ), so uncountably many
sequences arise. A sequence \((c_n)_{n\le0}\) that is \emph{left-eventually
periodic} (there are \(N\le0\) and \(L\ge1\) with \(c_{n-L}=c_n\) for
all \(n\le N\)) is determined by \((N,L,c_{N-L+1},\ldots,c_0)\), so
there are only countably many. Choose a sequence that is not
left-eventually periodic. Modulo \(R\), \(q_n=1\) satisfies
\(bq_{n+1}\equiv aq_n+c_n\) because \(c_n\equiv b-a\pmod R\).

% Markdown AST block 277: Para
\emph{Step 2: units modulo prime powers.} For each \(p^k\parallel M\) we
define \(q_n\equiv1\pmod p\) with
\[
bq_{n+1}\equiv aq_n+c_n\pmod{p^k}.
\]
If \(p\nmid a\): \(q_0=1\) and
\[
q_n=a^{-1}(bq_{n+1}-c_n)\bmod p^k
\]
for \(n<0\); since
\(c_n\equiv b-a\pmod p\), \(q_{n+1}\equiv1\) gives
\[
q_n\equiv a^{-1}(b-(b-a))=1.
\]

% Markdown AST block 278: Para
If \(p\mid a\) (so \(b\) is invertible
modulo \(p^k\)): choose \(L\) with \(p^k\mid a^L\) and put
\[
q_n=\sum_{j=0}^{L-1}a^jb^{-j-1}c_{n-1-j}\bmod p^k.
\]
Then
\[
\begin{aligned}
bq_{n+1}-aq_n&=\sum_{j=0}^{L-1}a^jb^{-j}c_{n-j}-\sum_{j=1}^{L}a^jb^{-j}c_{n-j}\\
&=c_n-a^Lb^{-L}c_{n-L}\\
&\equiv c_n\pmod{p^k},
\end{aligned}
\]
and modulo \(p\),
\[
q_n\equiv b^{-1}c_{n-1}\equiv b^{-1}(b-a)\equiv1.
\]
All indices used are \(\le-1\).

% Markdown AST block 279: Para
The Chinese remainder theorem gives
\(\gcd(q_n,M)=1\) and
\[
bq_{n+1}\equiv aq_n+c_n\pmod M.
\]

% Markdown AST block 280: Para
\emph{Step 3.} \((q_n,\theta_n,c_n)_{n\le0}\) is a left-directed
\(M\)-pseudo-orbit with non-left-eventually-periodic carries. If
\(\operatorname{rad}(2ab)\mid M\), then \(c_n\equiv b-a\pmod p\) for every
\(p\mid2ab\) gives \(c_n\in\mathcal C(a,b)\) (for \(p\mid a\),
\(c_n\equiv b\); for \(p\mid b\), \(c_n\equiv-a\); for \(p=2\) with
\(a,b\) odd, \(c_n\) is even). Lemma 38 applies. \qedhere
\end{proof}

% Markdown AST block 281: Para
The hypothesis on the initial graph is essential: a procedure that starts
from a different sound initial restriction derived from properties of
true orbits does not belong to \(\mathfrak M_0\), and the witness need
not pass it.

% Markdown AST block 282: Corollary 40.
\begin{samepage}
\begin{corollary}\leavevmode\label{thm:40}
(a) For a fixed finite set \(\mathcal P\) of primes,
the reduced bases \(a/b\) for which some execution of \(\mathfrak M_0\)
with moduli composed of primes in \(\mathcal P\) (with any exponents)
succeeds are finitely many:
\[
a<2\prod_{p\in\mathcal P}p
\]
and
\(b<a\).

% Markdown AST block 283: OrderedList
\begin{enumerate}
\def\labelenumi{(\alph{enumi})}
\setcounter{enumi}{1}
\tightlist
\item
  The bases \(r_t=(2t+7)/(2t+5)\), \(t\ge0\), are reduced. A
  successful execution for \(r_t\) with final modulus \(M_t\) requires
  \[
  \operatorname{rad}(M_t)>(2t+7)/2,
  \]
  that is,
  \[
  \operatorname{rad}(M_t)\ge t+4.
  \]
  With the primes \(\{2,3,5,11,13\}\)
  of (1) alone, no execution succeeds for \(t\ge4287\). The same holds
  for \((14k+1)/(10k+1)\), which tend to \(7/5\), for \(k\ge613\).
\end{enumerate}
\label{thm-end:40}
\end{corollary}
\end{samepage}

% Markdown AST block 284: Proof.
\begin{proof}\leavevmode\label{proof:40}
(a) is (27). (b) The pairs are coprime: their differences are
\(2\) and their members are odd for the first family; for the second, a
common divisor divides
\[
5(14k+1)-7(10k+1)=-2
\]
and both members are
odd. The thresholds are
\[
2t+7\ge8580\iff t\ge4287
\]
(as
\(2\cdot4286+7=8579\)) and
\[
14k+1\ge8580\iff k\ge613.\qedhere
\]
\end{proof}

% Markdown AST block 285: Para
Part (b) says that the certificate for \(7/5\) does not transfer to
nearby bases by the same prime set; it says nothing about composite
terms for those bases.

% Markdown AST block 286: 8.2 Finite windows are realized by true orbits
\Needspace{6\baselineskip}
\subsection{Finite windows are realized by true orbits}\label{sec:8.2}

% Markdown AST block 287: Lemma 41.
\begin{samepage}
\begin{lemma}\leavevmode\label{thm:41}
Let integers \(q_i\) with \(\gcd(q_i,M)=1\), reals
\(\theta_i\in[0,1)\) (\(0\le i\le L\)), and integers \(c_i\)
(\(0\le i<L\)) satisfy
\[
\begin{aligned}
&a\theta_i-b\theta_{i+1}=c_i\qquad\text{and}\\
&bq_{i+1}\equiv aq_i+c_i\pmod M.
\end{aligned}
\]
Then for every \(T_0\) there is
\(\xi>T_0\) with \(\lfloor\xi r^i\rfloor\equiv q_i\pmod M\) and
\(\{\xi r^i\}=\theta_i\) for \(0\le i\le L\).
\label{thm-end:41}
\end{lemma}
\end{samepage}

% Markdown AST block 288: Proof.
\begin{proof}\leavevmode\label{proof:41}
Write \(Q_i=q_i+Mt_i\) with integers \(t_i\) to be found, and
\[
d_i=(aq_i+c_i-bq_{i+1})/M\in\mathbb Z.
\]

% Markdown AST block 289: Para
The requirement
\(bQ_{i+1}=aQ_i+c_i\) is
\[
bt_{i+1}=at_i+d_i.
\]
Put \(S_0=0\),
\(S_{j+1}=aS_j+b^jd_j\), so that
\[
b^jt_j=a^jt_0+S_j
\]
by induction.

% Markdown AST block 290: Para
Since \(\gcd(a,b^L)=1\), choose
\[
t_0\equiv-a^{-L}S_L\pmod{b^L}.
\]
For
\(j\le L\),
\[
S_L\equiv a^{L-j}S_j\pmod{b^j},
\]
hence
\[
a^jt_0+S_j\equiv-a^ja^{-L}a^{L-j}S_j+S_j\equiv0\pmod{b^j},
\]
so
\[
t_j=(a^jt_0+S_j)/b^j
\]
is an integer, and \(bt_{j+1}=at_j+d_j\) holds
as an identity of rationals, hence of integers.

% Markdown AST block 291: Para
Adding \(b^LT\) to
\(t_0\) adds \(a^jb^{L-j}T\) to \(t_j\); taking \(T\) large makes all
\(Q_i\) positive and as large as desired.

% Markdown AST block 292: Para
Finally
\(x_i=Q_i+\theta_i\) satisfies
\[
bx_{i+1}=aQ_i+c_i+b\theta_{i+1}=ax_i,
\]
so with \(\xi=x_0\) we get \(x_i=\xi r^i\) and
\(\lfloor x_i\rfloor=Q_i\). \qedhere
\end{proof}

% Markdown AST block 293: 8.3 The two-carry obstruction
\Needspace{6\baselineskip}
\subsection{The two-carry obstruction}\label{sec:8.3}

% Markdown AST block 294: Para
Let \(\Delta=a-b\) (not to be confused with the phase gcd \(d\)).

% Markdown AST block 295: Theorem 42.
\begin{samepage}
\begin{theorem}\leavevmode\label{thm:42}
Let \(M\ge2\) and let \(R\) be the product of the primes
dividing \(M\) that do not divide \(2ab\) (\(R=1\) allowed). Suppose
there is an integer \(m\ge1\) with
\[
h=2Rm\le a-1,\qquad c=\Delta-h,\qquad \gcd(c,ab)=1 .
\tag{28}\label{eq:28}
\]
Then no execution of \(\mathfrak M_0\) (or of an extension by
pseudo-orbit-sound operations) with all \(M_t\mid M\) succeeds.
\label{thm-end:42}
\end{theorem}
\end{samepage}

% Markdown AST block 296: Proof.
\begin{proof}\leavevmode\label{proof:42}
\emph{Carry conditions.} We have
\[
\begin{aligned}
&1-b\le c<\Delta\le a-1,\qquad c\equiv\Delta\pmod2\quad(h\text{ is even}),\\
&c\equiv\Delta\pmod R,\\
&\text{and}\quad\gcd(\Delta,ab)=\gcd(a-b,a)\gcd(a-b,b)=1;
\end{aligned}
\]
so both \(c\) and
\(\Delta\) lie in \(\mathcal C(a,b)\). (The value \(c=0\) is excluded by
\(\gcd(c,ab)=1\).)

% Markdown AST block 297: Para
\emph{Fixed points and an invariant interval.} Put
\[
\begin{aligned}
\psi_\Delta(x)&=(bx+\Delta)/a=1+\tfrac ba(x-1)\qquad \text{and}\\
\psi_c(x)&=(bx+c)/a=\psi_\Delta(x)-h/a;
\end{aligned}
\]
these are the backward maps
\(\theta_{n+1}\mapsto\theta_n\) for the carries \(\Delta\) and \(c\).

% Markdown AST block 298: Para
For the word \(W_L=\Delta^Lc\) (in forward time order), the backward
composition is
\[
\begin{aligned}
\Phi_L&=\psi_\Delta^L\circ\psi_c,\\
\Phi_L(x)&=1+\Bigl(\tfrac ba\Bigr)^{L+1}(x-1)-\tfrac ha\Bigl(\tfrac ba\Bigr)^L,
\end{aligned}
\]
with fixed point
\[
x_L=1-hb^L/(a^{L+1}-b^{L+1}).
\]
Since \(a>b\),
\(x_L<x_{L+1}\), and \(x_L\uparrow1\) with \(x_L<1\). As
\(c\ge1-b\), \(\max(0,-c/b)<1\), so there is a least \(L\) with
\[
\max(0,-c/b)<x_L;
\]
put \(I=[x_L,x_{L+1}]\). The maps
\(\Phi_L,\Phi_{L+1}\) are increasing affine maps with slopes in
\((0,1)\) and fixed points at the left and right endpoints of \(I\)
respectively, so both map \(I\) into \(I\).

% Markdown AST block 299: Para
For \(x\in I\),
\(\psi_c(x)>0\) because \(x>-c/b\), and
\[
\psi_c(x)<\psi_\Delta(x)<1
\]
because \(c<\Delta\) and \(x<1\);
\(\psi_\Delta\) maps \((0,1)\) into \((\Delta/a,1)\). Hence all
intermediate values of a backward application of \(W_L\) or
\(W_{L+1}\) starting in \(I\) lie in \((0,1)\).

% Markdown AST block 300: Para
\emph{Infinite concatenation.} Let \(U=W_L\), \(V=W_{L+1}\), and let
\(B_1B_2\cdots\) be any sequence of blocks in \(\{U,V\}\), block \(i\)
starting at time \(n_i\) (\(n_1=0\)). For \(k\ge i\) the sets
\[
I_i^{(k)}=\Phi_{B_i}\circ\cdots\circ\Phi_{B_k}(I)
\]
are closed
subintervals of \(I\), decreasing in \(k\), of length at most
\[
(b/a)^{(L+1)(k-i+1)}|I|\to0;
\]
their intersection is a single point
\(\theta_{n_i}\in I\), and
\[
\Phi_{B_i}(\theta_{n_{i+1}})=\theta_{n_i}
\]
by construction.

% Markdown AST block 301: Para
Inside each block the values \(\theta_n\) are defined
by the backward recursion from \(\theta_{n_{i+1}}\), whose composition is
\(\Phi_{B_i}\). Thus \(0<\theta_n<1\) and
\[
a\theta_n-b\theta_{n+1}=c_n
\]
exactly for all \(n\ge0\).

% Markdown AST block 302: Para
\emph{Nonperiodic coding.} The words \(UV\) and \(VU\) have the same length
\(2L+3\) and differ at position \(L\) (letters \(c\) and \(\Delta\)).
Encode a binary sequence that is not eventually periodic (for instance
\(1\,0\,1\,00\,1\,000\cdots\)) by \(0\mapsto UV\), \(1\mapsto VU\). If
the carry sequence were eventually periodic with period \(P\), it would
have period \(P(2L+3)\) and block decoding would make the binary
sequence eventually periodic.

% Markdown AST block 303: Para
\emph{Unit residues modulo \(M\).} Every prime of \(M\) divides \(2abR\). For
\(p^e\parallel M\): if \(p\mid a\), take \(q_0=1\) and
\[
q_{n+1}=b^{-1}(aq_n+c_n)\bmod p^e;
\]
modulo \(p\),
\[
q_{n+1}\equiv b^{-1}c_n\ne0
\]
since \(c_n\in\{c,\Delta\}\) is coprime
to \(a\).

% Markdown AST block 304: Para
If \(p\mid b\), take
\[
q_n=-\sum_{j=0}^{e-1}b^ja^{-j-1}c_{n+j}\bmod p^e;
\]
as in the proof of
Proposition 35,
\[
\begin{aligned}
&bq_{n+1}-aq_n\equiv c_n\pmod{p^e},\qquad\text{and}\\
&q_n\equiv-a^{-1}c_n\ne0\pmod p.
\end{aligned}
\]

% Markdown AST block 305: Para
If \(p\mid R\), take
\(q_0\equiv-1\pmod{p^e}\) and
\[
q_{n+1}=b^{-1}(aq_n+c_n)\bmod p^e
\]
(\(b\) is a unit modulo \(p^e\) since \(p\nmid2ab\)); reducing modulo
\(p\), \(c_n\equiv\Delta\) gives
\[
q_{n+1}\equiv b^{-1}(-a+a-b)=-1,
\]
so every \(q_n\) is a unit.

% Markdown AST block 306: Para
If \(p=2\) and \(a,b\) are odd,
then \(c_n\) is even; take \(q_0\) odd, and
\[
q_{n+1}\equiv q_n+c_n\equiv q_n\pmod2.
\]
Combine by the Chinese
remainder theorem.

% Markdown AST block 307: Para
The result is a right-directed \(M\)-pseudo-orbit with carries in
\(\mathcal C(a,b)\) that are not eventually periodic. Lemma 38
applies. \qedhere
\end{proof}

% Markdown AST block 308: Corollary 43 (external primes are necessary).
\begin{samepage}
\begin{corollary}[external primes are necessary]\leavevmode\label{thm:43}
For every reduced base
\(a>b\ge2\), no execution of \(\mathfrak M_0\) whose moduli involve only
primes dividing \(2ab\) succeeds.
\label{thm-end:43}
\end{corollary}
\end{samepage}

% Markdown AST block 309: Proof.
\begin{proof}\leavevmode\label{proof:43}
Take \(R=1\) in Theorem 42. If \(\Delta\) is odd and
\(\ge3\), take
\[
c=1,\qquad h=\Delta-1\ge2.
\]

% Markdown AST block 310: Para
If \(\Delta\) is even and
\(\ge4\), then \(a,b\) are odd, \(c=2\) is coprime to \(ab\), and
\[
h=\Delta-2\ge2.
\]

% Markdown AST block 311: Para
If \(\Delta=1\), take
\[
c=-1\ge1-b,\qquad h=2\le a-1
\]
(as \(a\ge3\)).

% Markdown AST block 312: Para
If \(\Delta=2\), then \(a,b\) are odd, \(b\ge3\),
\(a\ge5\), and
\[
c=-2\ge1-b,\qquad h=4\le a-1.
\]

% Markdown AST block 313: Para
In every case
\(h\le a-1\) is even and positive, so (28) holds with \(m=h/2\). \qedhere
\end{proof}

% Markdown AST block 314: Para
For \(7/5\) with \(M=4290\) we have \(R=3\cdot11\cdot13=429\) and
\(h=858m\le6\) is impossible, consistent with Theorem 1(i).

% Markdown AST block 315: 8.4 The external-prime budget
\Needspace{6\baselineskip}
\subsection{The external-prime budget}\label{sec:8.4}

% Markdown AST block 316: Para
Let \(J(N)\) denote the Jacobsthal function: the least positive integer
\(J\) such that every \(J\) consecutive integers contain one coprime to
\(N\); \(J(1)=1\). Equivalently, \(J(N)\) is the largest gap between
consecutive integers coprime to \(N\). Let
\[
N_0=\prod\{p\mid ab:p\text{ odd}\}
\]
and \(k=\omega(N_0)\) its number
of prime factors; since \(a>b\ge2\) are coprime, \(N_0\ge3\).

% Markdown AST block 317: Theorem 44.
\begin{samepage}
\begin{theorem}\leavevmode\label{thm:44}
Let \(R\) be as in Theorem 42. If some execution of
\(\mathfrak M_0\) with all \(M_t\mid M\) succeeds, then
\[
a\le2R\,J(N_0).
\tag{29}\label{eq:29}
\]
Moreover, for every odd squarefree \(N\) with \(\omega(N)=k\),
\[
J(N)\le3^k;
\]
hence
\[
R\ge\frac{a}{2\cdot3^{k}},\qquad 2^kk!\le N_0<a^2,
\tag{30}\label{eq:30}
\]
so that \(3^k=a^{o(1)}\) and
\[
\liminf\log R/\log a\ge1
\]
along any
sequence of successful certificates with \(a\to\infty\). If the primes
available for \(R\) are restricted to a fixed finite set with product
\(R_0\), then
\[
a\le2R_0\,3^{k_0-1},
\]
where \(k_0\ge4\) is any integer
with
\[
2^{k_0}k_0!>4R_0^2\,9^{k_0};
\]
in particular only finitely many
reduced bases admit a successful execution with that set.
\label{thm-end:44}
\end{theorem}
\end{samepage}

% Markdown AST block 318: Proof.
\begin{proof}\leavevmode\label{proof:44}
By Theorem 42, success implies
\[
\gcd(\Delta-2Rm,ab)>1
\]
for
every
\[
1\le m\le H:=\lfloor(a-1)/(2R)\rfloor.
\]

% Markdown AST block 319: Para
Since \(2R\) is
invertible modulo \(N_0\), let
\[
u\equiv\Delta(2R)^{-1}\pmod{N_0};
\]
for
an odd prime \(p\mid ab\), \(p\mid\Delta-2Rm\) if and only if
\(m\equiv u\pmod p\). The prime \(2\) never divides \(\Delta-2Rm\) when
\(2\mid ab\) (then \(\Delta\) is odd), and is not a factor of \(ab\)
otherwise.

% Markdown AST block 320: Para
Hence
\[
\gcd(\Delta-2Rm,ab)=1\iff\gcd(m-u,N_0)=1.
\]

% Markdown AST block 321: Para
If
\(H\ge J(N_0)\), the \(H\) consecutive integers \(1-u,\ldots,H-u\)
contain one coprime to \(N_0\), giving a witness \(m\). So success
forces \(H\le J(N_0)-1\), that is,
\[
(a-1)/(2R)<J(N_0),
\]
which is
(29).

% Markdown AST block 322: Para
For the bound on \(J\): in a window of \(H\) consecutive integers, the
number coprime to \(N\) is
\[
\sum_{t\mid N}\mu(t)\,\#\{x:t\mid x\},
\qquad\text{and}\qquad
\#\{x:t\mid x\}=H/t+\epsilon_t
\]
with \(|\epsilon_t|<1\),
\(\epsilon_1=0\); so the count is at least
\[
H\varphi(N)/N-(2^k-1).
\]

% Markdown AST block 323: Para
As \(N\) is odd,
\[
\varphi(N)/N=\prod(1-1/p)\ge(2/3)^k,
\]
and \(H=3^k\) gives a count
\(\ge1\).

% Markdown AST block 324: Para
The first inequality in (30) follows from (29); the second from
\[
N_0\ge\prod_{i\le k}(2i+1)>2^kk!\qquad \text{and}\qquad N_0\le ab<a^2.
\]

% Markdown AST block 325: Para
From
\[
k!<a^2\qquad \text{and}\qquad k!\ge(k/e)^k
\]
one gets
\[
k\log k\ll\log a,
\]
so
\(k=O(\log a/\log\log a)\) and
\[
3^k=a^{O(1/\log\log a)}.
\]

% Markdown AST block 326: Para
For the effective bound, the ratio \(2^kk!/(4R_0^2\,9^k)\) is multiplied
by \(2(k+1)/9>1\) when \(k\ge4\), so it exceeds \(1\) for all
\(k\ge k_0\).

% Markdown AST block 327: Para
If \(k\ge k_0\), then
\[
a^2>N_0\ge2^kk!>4R_0^2\,9^k\ge(2R_0\,3^k)^2\ge a^2
\]
by (29) and
\(J(N_0)\le3^k\), a contradiction; so \(k\le k_0-1\) and
\[
a\le2R_0\,3^{k_0-1}.\qedhere
\]
\end{proof}

% Markdown AST block 328: Para
The condition (29) is necessary, not sufficient: closing every \(m\le H\)
does not guarantee a certificate.

% Markdown AST block 329: Corollary 45.
\begin{samepage}
\begin{corollary}\leavevmode\label{thm:45}
(a) If \(a=p\) is an odd prime and \(b=2^s<p\), then
\(N_0=p\), \(J(p)=2\), and a successful execution requires
\[
R\ge p/4.
\]
(b) If \(a,b\) are distinct odd primes, then \(J(ab)=3\)
and a successful execution requires
\[
R\ge a/6.
\]
\label{thm-end:45}
\end{corollary}
\end{samepage}

% Markdown AST block 330: Proof.
\begin{proof}\leavevmode\label{proof:45}
Among two consecutive integers at most one is divisible by
\(p\), so \(J(p)=2\). Among three consecutive integers at most one is
divisible by \(a\) and at most one by \(b\), so \(J(ab)\le3\); by the
Chinese remainder theorem there are consecutive \(x\equiv0\pmod a\),
\(x+1\equiv0\pmod b\), so \(J(ab)=3\). Apply (29). \qedhere
\end{proof}

% Markdown AST block 331: Para
These are constraints on any certificate for the bases named; they are
not statements about composite terms of those bases.

% Markdown AST block 332: Remark (Iwaniec’s bound).
\begin{remark}[Iwaniec’s bound]
Let \(C(k)\) be the maximal number of
consecutive integers each divisible by one of \(k\) arbitrarily chosen
primes. Iwaniec proved
\[
C(k)\ll k^2\log^2k
\]
\hyperref[ref-iw78]{IW78, Corollary, p.226}; see also
\hyperref[ref-cw13]{CW13, §1}, where the same statement is written as
\(g(n)\le X(k\log k)^2\) for \(k=\omega(n)\). Since
\(J(N)\le C(\omega(N))+1\) for \(\omega(N)\ge2\) and \(J(N)=2\) for
\(\omega(N)=1\), this gives
\[
J(N_0)\ll(k\log2k)^2
\]
with an absolute
implied constant, and with \(k\log2k\ll\log a\) from (30), the
condition (29) yields
\[
R\gg\frac{a}{(\log a)^2}.
\]
This sharpens (30) quantitatively. It is an external input, not a result
of this paper; the constant is absolute but not specified here, and no
theorem above depends on it. For explicit finite inputs, (29) and the
gcd test of Theorem 42 are used directly.
\end{remark}

% Markdown AST block 333: Para
The explicit certificates for the bases \(101/2\) and \(101/100\), the
selection of primes that eliminate persistent return words, and the
whole-word interval test are given in Appendix E.4.

% Markdown AST block 334: 9. Verification and reproducibility
\Needspace{6\baselineskip}
\section{Verification and reproducibility}\label{sec:9}

% BEGIN PUBLIC ADDITION availability
The paper sources, computation code, Lean formal proofs, reference data, and reproduction instructions are publicly available at:

\url{https://github.com/ixixi/rational-floor-certificates}

The repository README describes the contents; REPRODUCE.md and ENVIRONMENT.md explain how to reproduce the results and set up the required environment. Appendix A gives the main commands.
% END PUBLIC ADDITION availability

% Markdown AST block 335: 9.1 What the finite computations establish
\subsection{What the finite computations establish}\label{sec:9.1}

% Markdown AST block 336: Para
The table entries of Section 5 and Table 7 are summaries, not the
certificates. The mathematical obligations discharged by the finite
computations, and checked element by element by two implementations,
are the following.

% Markdown AST block 337: OrderedList
\begin{enumerate}
\def\labelenumi{\arabic{enumi}.}
\tightlist
\item
  \emph{Initial coverage.} The initial vertex set is the full set of legal
  states (all of \(V(M,K_0)\) for Table 1 and for Table 7; all \(K\)
  cells for Tables 4 and 6).
\item
  \emph{Complete edge sets.} Every edge satisfying (E1)--(E3), or the initial
  inequalities of Theorem 19, is present; the lift of every edge is
  computed for every starting residue.
\item
  \emph{Residues at every intermediate time.} In a lift, all
  \(\ell+1\) residues along a word of length \(\ell\), including the
  endpoint, are nonzero.
\item
  \emph{All-letter phase conditions.} For every certified component, (19)
  holds at every letter position of every internal edge, with a single
  label per phase.
\item
  \emph{Rank monotonicity.} For Table 7, the inequalities of (Q-rank) and
  (Q-cyc) hold on every edge between distinct blocks, and the bounds
  \(\rho\le65\), \(\eta\le6\), \(d_B\le13\) hold.
\item
  \emph{Retention.} Every vertex and every internal edge of every uncertified
  cyclic component is kept, and nothing else.
\item
  \emph{Output-preserving contraction.} Macro vertices are exactly the
  vertices of out-degree at least two (plus anchors), each chain of
  out-degree-one vertices ends at a macro vertex, and macro words are the
  concatenations of the chain words.
\item
  \emph{Final emptiness.} The last retained set (Table 1) or the last output
  graph (Tables 4 and 6) is empty; for Table 7, no cyclic block is
  uncertified.
\end{enumerate}

% Markdown AST block 338: Para
The strongly connected components themselves are verified by sufficient
conditions: the proposed blocks are nonempty, disjoint, cover all
vertices, are internally reachable forwards and backwards, and admit a
strictly monotone integer rank on all edges between distinct blocks
(acyclic quotient). These conditions determine the SCC partition.
Agreement of counts or of hashes is recorded but is not the criterion;
the criterion is the element-wise comparison of the complete data.

% Markdown AST block 339: 9.2 The one-step computation for 7/5
\Needspace{6\baselineskip}
\subsection{The one-step computation for 7/5}\label{sec:9.2}

% Markdown AST block 340: Para
The one-step calculation of Tables 1--3 was performed by two
implementations in Python's arbitrary-precision integer arithmetic,
using only the standard library and no floating point in the proof
computation. Implementation A enumerates target cells with (12) and
computes SCCs by iterative Kosaraju traversal. Implementation B tests
(E3) directly for every \(k\in\{0,\ldots,K-1\}\), caching the resulting
lists by \((j,e)\), and computes SCCs by an iterative form of Tarjan's
algorithm \hyperref[ref-t]{T, §4}. Both classify all unit residues by
\(bz\bmod M\) to solve (E2). B was developed from the mathematical
specification and a common output schema, without reading or importing
A or the supplied baseline program; its design, source, and outputs were
fixed before the comparison with A.

% Markdown AST block 341: Para
The validation checks every vertex and labelled edge. It recomputes
shortest distances, all edge differences, phase label collections,
certification flags, and primitive words, checks that the retained set is
exactly the union of uncertified cyclic components, and checks that every
child in (16) is present in the next stage. For a rejected phase test,
two internal edges with the same departure phase and distinct labels
provide a checkable witness. Across the four main stages and the three
first stages of Table 3, the comparison covers 97772 vertices, 168474
labelled edges, and 94006 SCCs; no mismatch in the scientific data was
found. The recorded executions used CPython 3.11.13 for A and
CPython 3.10.12 for B on Linux, each case with a 600-second wall-time
limit and a 4 GiB address-space limit. Reaching a limit is recorded as
\texttt{inconclusive} and supplies neither a counterexample nor a minimality
claim; all adopted cases completed. The programs are included under \nolinkurl{supplement/computations/implementation_a/} and \nolinkurl{supplement/computations/implementation_b/}.
The public verification procedure and the all-element comparison are described in the \href{supplement/COMPUTATION.md}{computation guide} and the \href{REPRODUCE.md}{reproduction guide}.

% Markdown AST block 342: 9.3 The word-labelled computations
\Needspace{6\baselineskip}
\subsection{The word-labelled computations}\label{sec:9.3}

% Markdown AST block 343: Para
The computations of Tables 4, 6, and 7 were performed by two
implementations with disjoint scientific cores.

% Markdown AST block 344: Para
\emph{Implementation A} (\href{supplement/COMPUTATION.md}{\nolinkurl{word_a}}) is a
provenance-tracked working copy of the programs supplied with the
research notes on which this paper is based (C++ constructions driven by
Python, with a Python third check for the first five stages of
Table 6), together with a driver that fixes resource limits, records
inputs and outputs, and writes the graphs in a normalized text format.
The supplied programs contain two constructions labelled ``A'' and ``B'';
these share a common origin, file formats, and conventions, and their
mutual agreement is a consistency check within one lineage, not the
independent verification described next. The scientific cores are retained in the public sources. The provenance and
packaging changes are described in the \href{supplement/COMPUTATION.md}{computation guide}; the \href{MANIFEST.json}{source manifest}
identifies every distributed source by SHA-256.

% Markdown AST block 345: Para
\emph{Implementation B} (\href{supplement/COMPUTATION.md}{\nolinkurl{word_b}}) was
written by a different programmer from a mathematical
specification equivalent to Section 4 and the
comparison contract
\href{supplement/computations/word-schema.md}{\nolinkurl{word-schema.md}},
which fixes the definitions, the text format, and the normalization
only, without reading the supplied programs, implementation A, or the
research notes' tables. It consists of a C++17 production program and a
separately written literal Python reference program (direct inequality
tests for all \((j,c,k)\), forward iteration for all starting residues,
Kosaraju's algorithm, and Kahn's algorithm for acyclicity), whose
normalized outputs coincide byte for byte on all graphs. B computes
\(h\) as the word length of the first root path found, uses forward
iteration rather than Lemma 18, and re-verifies its SCC partitions by the
sufficient conditions of Section 9.1. Its outputs were fixed, with sizes
and hashes recorded, before A's outputs were opened.

% Markdown AST block 346: Para
\emph{Comparison.} B's comparator compares, for each of the 29 file pairs
(ten input and ten output graphs for Table 6, four and four for Table 4,
one graph for Table 7), the header parameters, the vertex sets, and the
sets of word-labelled triples \((u,w,v)\), and checks the sort order,
absence of duplicates, and consistency of each file. For Table 7 it
compares the block partition as a family of member sets (independent of
block numbering), the attributes of each block (size, internal edges,
\(d_B\), \(\rho\), \(\eta\)), the phase-label sequences, and the absence
of uncertified cyclic blocks. All 29 pairs agree element by element. The comparison procedure and
reference data are given in the \href{supplement/COMPUTATION.md}{computation guide} and \href{checks/reference.json}{checks/reference.json};
Appendix A gives the public reproduction entry point.
Both implementations compute \(h\) from total word lengths of paths
inside each component from a common root. Corollary 22 therefore
implies agreement of the reported \(d_B\) values, even when their
roots or paths differ.

% Markdown AST block 347: Para
\emph{Environment and limits.} A ran under Python 3.10.12 with g++ 11.4.0
on Linux, with wall-time and address-space limits of 600 s and 4 GiB per
supplied construction for Table 6 (1500 s overall) and 300 s and 4 GiB
otherwise. B ran its C++ program under limits of 3600 s and 8 GiB
(Table 6), 300 s and 4 GiB (Table 4), and 600 s and 4 GiB (Table 7),
its Python reference under 1800 s and 8 GiB, and the comparison under
1800 s and 8 GiB. No limit was reached; a limit reached would have been
recorded as \texttt{inconclusive}. Wall times are recorded in the run records
and are not performance claims.

% Markdown AST block 348: Para
\emph{Meaning of independence.} The independence concerns the scientific
cores of the two implementations: edge enumeration, lifting, component
decomposition, certification, and contraction were coded separately from
the specification. It does not extend to the mathematical lemmas, which
are common, to the shared text format, or to human refereeing, and it
does not mean clean-room development in a legal sense. Agreement of the
two implementations is evidence that they computed the same finite
objects; the soundness of the operations (Section 4) and the passage to
all \(\xi>0\) (Theorem 19) are mathematical statements proved above, not
outputs of the programs.

% Markdown AST block 349: 9.4 The Lean formalization and its boundary
\Needspace{6\baselineskip}
\subsection{The Lean formalization and its boundary}\label{sec:9.4}

% Markdown AST block 350: Para
A Lean formalization covers both parts of Theorem 1 and Corollary 2,
with different trust bases for the finite evaluations. It uses
Lean 4.29.1 and Mathlib at commit
\nolinkurl{5e932f97dd25535344f80f9dd8da3aab83df0fe6}, with transitive dependencies
pinned in \nolinkurl{supplement/lean/lake-manifest.json}.

% Markdown AST block 351: Para
\textbf{The one-step certificate.} The one-step formalization follows the
integer-difference proof of Lemma 4, orbit inclusion (Proposition 5), a
rank-and-phase version of tail preservation, and the finite iteration of
Theorem 10.

% Markdown AST block 352: Para
The formal certificate avoids a formal proof of the SCC algorithms. At
each stage it assigns each vertex to a block, and each block a natural
number rank. On every edge between distinct blocks the rank must strictly
decrease (the opposite orientation to the strictly increasing \(\rho\) of
(Q-rank) in Section 6.2; either convention yields a well-founded order).
Each block is designated as one of the following:

% Markdown AST block 353: OrderedList
\begin{enumerate}
\def\labelenumi{\arabic{enumi}.}
\tightlist
\item
  \textbf{Retained:} every vertex of the block is included in the retained set.
\item
  \textbf{Edgeless:} the block contains no internal edge.
\item
  \textbf{Periodic:} a positive integer \(d\), a phase in
  \(\{0,\ldots,d-1\}\) for every vertex, and a label for every phase
  satisfy the phase-increment and departure-label conditions on every
  internal edge.
\end{enumerate}

% Markdown AST block 354: Para
These are sufficient conditions for Lemma 8. Along an infinite walk the
rank cannot strictly decrease infinitely often, so the walk eventually
remains in one block. That block cannot be edgeless. If it is periodic,
its label sequence contradicts Lemma 4. It must therefore be retained.
An SCC decomposition supplies such blocks and ranks through its acyclic
quotient, but the formal implication does not require the blocks to be
identified as SCCs.

% Markdown AST block 355: Para
The abstract structure \texttt{FiniteSuccess} requires initial coverage of all
legal states, the rank/phase condition on \textbf{every} edge satisfying
(E1)--(E3) between current states, coverage of all retained-block
vertices, inclusion of all refined states at the next stage, and an empty
final retained set. The following declarations, all with namespace prefix
\texttt{MathPaper.}, identify the principal formal dependencies.

% Markdown AST block 356: Table 8.
% Caption is rendered at the start of the immediately following table.

% Markdown AST block 357: Table 8
\begingroup
\small
\setlength{\tabcolsep}{4pt}
\renewcommand{\arraystretch}{1.15}
\setlength{\LTcapwidth}{\textwidth}
\begin{longtable}{@{}>{\raggedright\arraybackslash}p{\dimexpr0.44\dimexpr\linewidth-8pt\relax} >{\raggedright\arraybackslash}p{\dimexpr0.56\dimexpr\linewidth-8pt\relax}@{}}
\caption{Principal Lean declarations. All names have prefix
\texttt{MathPaper.}; \texttt{Pipe.} denotes the executable word-labelled pipeline.}\label{tab:8}\\
\toprule
Mathematical content & Lean declaration \\
\midrule
\endfirsthead
\toprule
Mathematical content & Lean declaration \\
\midrule
\endhead
\bottomrule
\endlastfoot
Growth (7) & \nolinkurl{floor_orbit_eventually_gt} \\
Lemma 4 & \nolinkurl{floor_carry_not_eventually_periodic} \\
Proposition 5 and legal orbit states & \nolinkurl{orbit_edge},\allowbreak{} \nolinkurl{orbit_state_legal} \\
Phase-periodic labels & \nolinkurl{phase_labels_periodic} \\
Tail preservation from rank/phase data & \nolinkurl{rank_certificate_retained_tail},\allowbreak{} \nolinkurl{certified_orbit_retained_tail} \\
Lemma 9 & \nolinkurl{floor_refinement_bounds},\allowbreak{} \nolinkurl{refined_state_mem} \\
Theorem 10 in certificate form & \nolinkurl{finite_success_visits} \\
Conditional compositeness conclusion & \nolinkurl{finite_success_composite} \\
Signed interval enumeration (12) & \nolinkurl{edge_interval_iff} \\
Concrete four-stage certificate & \texttt{mainFiniteSuccess} \\
Theorem 1(i) & \nolinkurl{floor_seven_fifths_visits} \\
Corollary 2 for \(7/5\) & \nolinkurl{floor_seven_fifths_composite} \\
Word-labelled covering and phase output,\allowbreak{} Lemma 12(b) & \texttt{Covers},\allowbreak{} \nolinkurl{phase_output_periodic} \\
Word-labelled tail preservation,\allowbreak{} Proposition 13 & \nolinkurl{WordCert.stage_tail} \\
Reaching macro vertices and preserving the output,\allowbreak{} Lemma 15 and the forward direction of Proposition 16 & \nolinkurl{macro_visited},\allowbreak{} \nolinkurl{chain_realized},\allowbreak{} \nolinkurl{contract_tail} \\
Prime lifting,\allowbreak{} Proposition 17 & \nolinkurl{runWord_orbit},\allowbreak{} \nolinkurl{lift_tail} \\
Theorem 19 and Corollary 20 in certificate form & \nolinkurl{word_finite_success_visits},\allowbreak{} \nolinkurl{word_finite_success_composite} \\
Lemmas 23--24 & \nolinkurl{carry_alphabet_75},\allowbreak{} \nolinkurl{carry_alphabet_52} \\
Initial coverage,\allowbreak{} stage checking,\allowbreak{} lifting,\allowbreak{} and pipeline soundness & \nolinkurl{Pipe.initialStage_sound},\allowbreak{} \nolinkurl{Pipe.checkStage_sound},\allowbreak{} \nolinkurl{Pipe.liftArr_sound},\allowbreak{} \nolinkurl{Pipe.wordPipeline_sound} \\
Finite word-labelled evaluations & \nolinkurl{five_halves_ok},\allowbreak{} \nolinkurl{seven_fifths_ok} \\
Theorem 1(ii) and Corollary 2 for \(5/2\) & \nolinkurl{floor_five_halves_visits},\allowbreak{} \nolinkurl{floor_five_halves_composite} \\
Theorem 1(i),\allowbreak{} alternative word-labelled proof & \nolinkurl{floor_seven_fifths_visits_word} \\
\end{longtable}
\endgroup

% Markdown AST block 358: Para
The closed term \texttt{mainFiniteSuccess\ :\ FiniteSuccess\ 7\ 5\ 4290} discharges
the certificate argument of the general reduction. The two one-step final theorems
take only a real coefficient \(\xi\), a proof of \(\xi>0\), and an
arbitrary \(N\in\mathbb N\), and return an \(n\ge\max(N,1)\) satisfying
the respective conclusion, with no unproved finite-computation
hypothesis. The finite checker does not assume that the edge list
exported by Python is complete: proved enumeration lemmas place every
edge satisfying (E1)--(E3) among three carry candidates, five residue
lifts, and three target-cell candidates, and the checker tests the
rank/phase condition on an outer approximation of these edges. The
Boolean checks are proved by \texttt{decide\ +kernel}, split into smaller
statements joined by proved composition lemmas. The compositeness proof
uses the eventually valid bound \(Q_n>M\) and defines compositeness as
\(Q_n=uv\) with integers \(u,v>1\), without dependence on prime
factorization.

% Markdown AST block 359: Para
The complete one-step certificate, main module, and aggregate build succeeded.
The axiom audit of 30 designated declarations, including
\texttt{mainFiniteSuccess} and both final theorems, shows dependencies confined
to Lean's standard \texttt{propext}, \texttt{Classical.choice}, and \texttt{Quot.sound}; no
\texttt{sorryAx}, custom axiom, or external-computation axiom occurs
(the \href{supplement/FORMALIZATION.md}{formalization guide}). An independent audit
rebuilt all 79 modules of the one-step project from a fresh project cache and checked
the types and axiom dependencies of 60 declarations (58 from the project
and two independently written restatements of Theorem 1(i) and its
corollary), with the same three standard axioms only. The declared scope and
commands for repeating the type and axiom checks are given in the \href{supplement/FORMALIZATION.md}{formalization guide}.
The rebuild used compiled library caches for the pinned Mathlib and its
dependencies, whose revisions were checked against the lock file; it did
not rebuild every dependency from source. For this one-step proof the
trust base comprises the fixed Lean kernel, its standard axioms, and
those fixed libraries. The declaration-level coverage is summarized in the \href{supplement/FORMALIZATION.md}{formalization guide}.
The public build procedure is given in the \href{REPRODUCE.md}{reproduction guide}.

% Markdown AST block 360: Para
\textbf{The word-labelled certificates.} The added formalization in
\nolinkurl{MathPaper/Word/} extends the block/rank method to nonempty edge words.
Its covering relation records a walk and strictly increasing boundary
times, with each edge word equal to the corresponding segment of the
carry sequence. In a periodic block the phase advances by the full word
length, and every letter must agree with its phase label. In a retained
block the checker verifies a macro-vertex set and a remaining word and
destination for each non-macro vertex. The consistency conditions make
the remaining word length decrease along a non-macro path and preserve
the concatenated output on passage to macro edges. Thus the formal proof
uses the forward direction of Proposition 16; it need not show that
every walk of the output graph expands to an input walk.

% Markdown AST block 361: Para
The kernel-checked theorem \nolinkurl{Pipe.wordPipeline_sound} proves that
\texttt{Pipe.wordPipeline\ a\ b\ K\ alphabet\ primes\ =\ true} implies
\texttt{WordFiniteSuccess\ a\ b\ K\ alphabet\ primes}. The pipeline enumerates the
initial cell edges using the strict inequalities (E3). At each stage,
auxiliary algorithms compute proposed blocks, ranks, phases, contraction
data, and output edges. Their correctness is not assumed: the proved
checker verifies the required conditions on every input edge. The
proved lifting function considers every nonzero starting residue and
every letter, including the residue at the final endpoint; its inverse
of \(b\) is checked modulo the newly added prime. Vertex renumbering
uses preservation of covering under an arbitrary vertex map, so no
unproved injectivity assumption is needed. The pipeline accepts only
after the final output is empty. Together with the two proved carry
alphabet lemmas, its soundness theorem discharges the finite-certificate
arguments of the general visitation and compositeness theorems.

% Markdown AST block 362: Para
The two accepted finite evaluations are

% Markdown AST block 363: reproduction command
\begin{Verbatim}[fontsize=\small,breaklines=true,breakanywhere=true]
Pipe.wordPipeline 5 2 4 [-1, 1, 3] [3, 7, 11, 13, 17, 19, 23, 29, 31] = true
Pipe.wordPipeline 7 5 2048 [-4, -2, 2, 4, 6] [3, 11, 13] = true
\end{Verbatim}

% Markdown AST block 364: Para
They are proved by \nolinkurl{native_decide}. In the fixed Lean version this
mechanism asserts the native evaluation through a generated axiom for
each theorem; it does not use \texttt{Lean.ofReduceBool}. The actual additional
axioms, with namespace prefix \texttt{MathPaper.}, are
\nolinkurl{five_halves_ok._native.native_decide.ax_1_1} and
\nolinkurl{seven_fifths_ok._native.native_decide.ax_1_1}, respectively. Accordingly,
the \(5/2\) final theorems depend on the first generated axiom in
addition to the three standard axioms, and the alternative \(7/5\)
theorem depends on the second. The word-labelled theory and checker
soundness depend only on the standard axioms. No \texttt{sorryAx} occurs in
the recorded dependencies. The added trust is in the compiler and native
evaluation of the Lean pipeline, not in a Python or C++ success flag,
exported graph, or hash. These final theorems take only \(\xi>0\) and
\(N\), with no certificate argument, and give exactly the quantifiers
and moduli of Theorem 1 and the \(5/2\) part of Corollary 2.

% Markdown AST block 365: Para
The added modules, aggregate build, and axiom display were accepted.
The \href{supplement/FORMALIZATION.md}{formalization guide} records their coverage and axiom boundary. A separate
independent audit of the word-labelled proof rebuilt all 14 added Word modules, reran both native finite evaluations,
and checked the types and axiom dependencies of 46 project declarations
and five independent restatements or checks. It confirmed the trust
boundary described above. That audit reused compiled one-step
project dependencies and the fixed library caches; it was not a fresh
recompilation of all one-step modules or of the libraries.
An attempted \texttt{decide\ +kernel}
evaluation of a smaller pipeline instance stopped at opaque hashing
operations, so neither word-labelled evaluation is claimed as a
kernel-only finite computation.

% Markdown AST block 366: Para
\textbf{Boundary.} The formalization establishes the arithmetic conclusions
of Theorem 1 and Corollary 2 with the trust bases just stated. It does
not formally identify the reported blocks as SCCs, prove that \(h\)
is a shortest distance or that \(d\) is the specified gcd, or establish
the reported table counts and words. Those data are checked by the
A/B computations; the Lean checker uses sufficient rank, phase, and
contraction conditions. Within Section 4, the gcd positivity argument
of Lemma 12(a), Lemma 14, the reverse correspondence in Proposition 16,
the root formula of Lemma 18, and Theorem 21 with Corollary 22 are not
formalized. Positive periods are required by the checker, and lifting
uses the forward recurrence. Theorem 3, Section 6 (including the
waiting-time certificate of Table 7), Sections 7--8, and the further
results of Section 10 and Appendices B--E also remain outside the
formalization. Their evidence consists of the proofs in this paper and,
where finite computation is required, the recorded A/B comparisons.

% Markdown AST block 367: 10. Further results deferred to the appendices
\Needspace{6\baselineskip}
\section{Further results deferred to the appendices}\label{sec:10}

% Markdown AST block 368: Para
Two groups of further results are stated here without proof; the proofs
are given in Appendices B and D, and the supporting computations were
performed by the two implementations of Section 9.3 and compared element
by element (Sections D.5 and E.2). They are not used in Sections 1--9.

% Markdown AST block 369: Para
\textbf{Recurrent carries.} An infinite word \(w=w_0w_1\cdots\) over a finite
alphabet is \emph{recurrent} if for every \(L\ge1\) and every \(B\ge1\) there
is \(n\ge B\) with
\[
w_{n+i}=w_i
\]
for \(0\le i<L\); every tail of a
recurrent word is recurrent.

% Markdown AST block 370: Theorem 46 (recurrent carries force composites; proof in Appendix B.3).
\begin{samepage}
\begin{theorem}[recurrent carries force composites; proof in Appendix B.3]\leavevmode\label{thm:46}
Let \(a>b\ge2\), \(\gcd(a,b)=1\), \(\xi>0\), and
\emph{suppose that the carry sequence \(c_0c_1\cdots\) of the orbit has a
recurrent tail} \(c_{j_0}c_{j_0+1}\cdots\). Then \((Q_n)\) contains
infinitely many composite terms. Moreover, if \(j\ge j_0\) satisfies
\(Q_j>1\) and \(\gcd(Q_j,ab)=1\), then \(Q_j\) divides infinitely many
later terms.
\label{thm-end:46}
\end{theorem}
\end{samepage}

% Markdown AST block 371: Para
The recurrence assumption concerns the carry word itself. Confinement
of a walk to a strongly connected component does not imply recurrence
of its output.

% Markdown AST block 371 (continued): Para
Theorem 46 is a return argument in the spirit of a theorem of Dubickas
\hyperref[ref-dub09]{Dub09, Theorem 4}: if integers \(x_n\) satisfy
\(x_{n+d}=kx_n+F(x_{n+1},\ldots,x_{n+d-1})\) with an integer \(k\ne0\)
and \(F\in\mathbb Z[z_1,\ldots,z_{d-1}]\), then \((x_n)\) is purely
periodic modulo every \(q\) with \(\gcd(k,q)=1\), and it contains
infinitely many composite terms if \(|x_n|\to\infty\). Stephan
\hyperref[ref-s26]{S26, Proposition 3.1} applies the case \(d=1\) to
ultimately periodic digit words in an integer base. For a nonintegral
base the carry word is never eventually periodic (Lemma 4). In
Theorem 46 the letters act on \(\mathbb Z/M\) by the affine permutations
\(x\mapsto b^{-1}(ax+c)\), the word is only assumed to be recurrent,
and the finite-prefix copying lemma (Lemma 49) replaces periodicity in
producing a return to the initial residue (Theorem 50).

% Markdown AST block 372: Para
\textbf{Finite-height results.} For the bases \(6/5\) and \(9/7\), a sieve
over carry words combined with exact interval and divisor tests covers
all initial integers up to an explicit height.

% Markdown AST block 373: Theorem 47 (finite height; proof in Appendix D.4).
\begin{samepage}
\begin{theorem}[finite height; proof in Appendix D.4]\leavevmode\label{thm:47}
(a) If \(6<\lfloor\xi\rfloor\le5^{40}\), then
\(\lfloor\xi(6/5)^j\rfloor\) is composite for some \(0\le j\le40\).
(b) If \(9<\lfloor\xi\rfloor\le7^{18}\), then
\(\lfloor\xi(9/7)^j\rfloor\) is composite for some \(0\le j\le18\).
\label{thm-end:47}
\end{theorem}
\end{samepage}

% Markdown AST block 374: Para
For \(6/5\) the shifted sequence
\(\lfloor\xi(6/5)^n\rfloor-1\) is treated in
\hyperref[ref-dn]{DN, Theorem 3, p.637}; Theorem 47(a) concerns the unshifted
sequence on a finite range and is a different statement.

% Markdown AST block 375: Para
Appendix C treats the family \((6t+3)/(6t+1)\), for which the
all-coefficient statement reduces to an exact integer halting problem
(Theorem 58), and Appendix E treats arbitrarily long prime prefixes
along that family (Theorem 68, using known results on narrow arithmetic
progressions of primes as external inputs), explicit return-word
obstructions for \(9/7\) and \(6/5\) (Corollary 70), and the items
deferred from Section 8.

% Markdown AST block 376: 11. Open problems
\Needspace{6\baselineskip}
\section{Open problems}\label{sec:11}

% Markdown AST block 377: OrderedList
\begin{enumerate}
\def\labelenumi{\arabic{enumi}.}
\tightlist
\item
  \emph{Further bases and infinite families.} It remains to determine which
  further rational bases admit finite certificates and to construct a
  uniform certificate for an infinite family. Theorem 44 gives a
  necessary budget of external primes, \(R\ge a^{1-o(1)}\), but no construction meeting it.
  Eliminating specific return words (Appendix E.3) and covering all
  recurrent components are different tasks.
\item
  \emph{Finite height versus all coefficients.} Theorem 47 leaves the
  initial integers above the stated heights untreated. For an initial
  integer \(s+kb^N\) with the same first \(N\) carries as \(s\), the
  terms are \(Q_j(s)+ka^jb^{N-j}\) by (8), and a divisor found for
  \(k=0\) need not divide the terms for other \(k\).
\item
  \emph{Recurrence.} Theorem 46 excludes recurrent carry tails on an
  orbit all of whose terms are prime, but it is open whether a
  non-recurrent all-prime tail can exist (Corollary 51); confinement
  to a strongly connected component of an outer graph does not imply
  recurrence (Appendix B.1).
\item
  \emph{Additional operations.} The whole-word interval test (Appendix E.4)
  is pseudo-orbit-sound but not walk-sound; whether adjoining it changes
  the class \(\mathcal F_{\rm cert}\) is open. Whether the halting
  program of Appendix C halts for every \(t\) (Theorem 58) is open.
\item
  \emph{Minimality and subsequences.} The certificates give finite prime
  sets without establishing their minimality. Whether the conclusions
  hold along prescribed arithmetic progressions or sparse subsequences
  of indices, and what bounds hold for the density of prime terms,
  remain separate questions.
\end{enumerate}

% Markdown AST block 378: Appendix A. Reproducing the computations and the formal proof
\appendix
\Needspace{6\baselineskip}\section{Reproducing the computations and the formal proof}\label{app:A}

% Public reproduction instructions.
The paper sources, scientific programs, and Lean project are distributed in
the public repository \href{https://github.com/ixixi/rational-floor-certificates}{rational-floor-certificates}.
Run commands from the root of the downloaded distribution. The \href{REPRODUCE.md}{reproduction guide}
gives the dependencies, commands, resource limits, and output locations.
The common entry point is:

\begin{Verbatim}[fontsize=\small,breaklines=true,breakanywhere=true]
python3 -B reproduce.py verify
python3 -B reproduce.py check --output ../certificate-preflight
python3 -B reproduce.py recompute --output ../certificate-replay
\end{Verbatim}

The distribution identifies its sources by the \nolinkurl{source_sha256}
field of the \href{MANIFEST.json}{source manifest}, computed from the SHA-256 values of the
distributed source and check files. \nolinkurl{SHA256SUMS} lists the
individual file hashes. \href{checks/reference.json}{checks/reference.json} records the finite reference
data used for comparison. Generated graphs and logs are written to the
chosen output location. Reaching a resource limit is recorded as
\nolinkurl{inconclusive}; an incomplete run is not an accepted certificate.

\textbf{Word-labelled computations, implementation A.} The public driver
\nolinkurl{supplement/computations/word_a/run_word.py}
applies resource limits, runs the supplied scientific constructions, writes
normalized graphs, and runs the same-lineage post-checks. The cases
\nolinkurl{five-halves}, \nolinkurl{cross-75}, and \nolinkurl{q75} regenerate the data of Tables 6, 4,
and 7, respectively. The \href{REPRODUCE.md}{reproduction guide} gives the complete commands;
the \href{supplement/COMPUTATION.md}{computation guide} describes the source provenance and comparison scope.
The C++ constructions and the Python normalization code are included.
Regeneration therefore uses the distributed sources.

\textbf{Edge multiplicities.} In Tables 4 and 6 edges are counted as distinct
triples. Counting contracted chains with multiplicity, and propagating
the multiplicity through lifting as the supplied research notes do,
gives for the ten stages of Table 6 the input edge counts 12, 17, 58,
167, 820, 7083, 60242, 315483, 1685458, 3754612 and the output edge
counts 12, 12, 20, 85, 529, 4256, 18069, 80704, 201199, 0; for Table 4
all multiplicities are one. Both implementations report these numbers
and agree on them; they carry no proof content.

\textbf{Word-labelled computations, implementation B.} The separate C++
program, Python reference, and comparator are included under
\nolinkurl{supplement/computations/word_b/}. The following commands show
the mathematical inputs and use an initially absent output directory
outside the distribution. The public entry point applies the limits and
performs the complete comparison described in Section 9.3:

\begin{Verbatim}[fontsize=\small,breaklines=true,breakanywhere=true]
mkdir -p ../certificate-manual/build \
  ../certificate-manual/OUT52 ../certificate-manual/OUT75 \
  ../certificate-manual/OUTQ75 ../certificate-manual/REF52
g++ -O2 -std=c++17 -o ../certificate-manual/build/wgb supplement/computations/word_b/wgb.cpp
../certificate-manual/build/wgb pipeline --a 5 --b 2 --K 4    --carries -1,1,3      --primes 3,7,11,13,17,19,23,29,31 --out ../certificate-manual/OUT52
../certificate-manual/build/wgb pipeline --a 7 --b 5 --K 2048 --carries -4,-2,2,4,6 --primes 3,11,13 --out ../certificate-manual/OUT75
../certificate-manual/build/wgb q75      --a 7 --b 5 --M 429 --K 2048 --carries -4,-2,2,4,6 --out ../certificate-manual/OUTQ75
python3 supplement/computations/word_b/wgb_ref.py pipeline --a 5 --b 2 --K 4 --carries -1,1,3 --primes 3,7,11,13,17,19,23,29,31 --out ../certificate-manual/REF52
python3 supplement/computations/word_b/compare_wordgraph.py batch --out ../certificate-manual/CMPDIR LABEL=FILE_A:FILE_B ...
\end{Verbatim}

The shared normalization contracts are
\href{supplement/computations/comparison-schema.md}{comparison-schema.md} and
\href{supplement/computations/word-schema.md}{word-schema.md}.
The reference data and artifact hashes are recorded in \href{checks/reference.json}{checks/reference.json}
and the \href{MANIFEST.json}{source manifest}. Hashes identify finite artifacts; they do not
replace element-wise comparison or the mathematical completeness arguments.

\textbf{One-step computation.} The public reproduction procedure also
regenerates the four cases of Tables 1–3 with both one-step implementations
and compares their complete outputs. Run without Python optimization,
since the independent verifier uses assertions for checks. The source
programs, reference checks, and generated Lean certificate sources are included;
the \href{supplement/COMPUTATION.md}{computation guide} specifies the compared data and the \href{REPRODUCE.md}{reproduction guide} gives the commands.

\textbf{Formal proof.} The generated Lean data and proof sources are included,
so checking the one-step proof does not require rerunning their generator.
The word-labelled modules define and execute their finite pipeline inside
Lean; they do not import the A/B graph dumps. Follow the \href{REPRODUCE.md}{reproduction guide}
to obtain the pinned library dependencies and build the project modules
sequentially with a fresh project cache. After the one-step finite modules
and before the final \nolinkurl{MathPaper} target and aggregate build, check
the word-labelled modules sequentially:

\begin{Verbatim}[fontsize=\small,breaklines=true,breakanywhere=true]
python3 supplement/lean/tools/check_finite.py \
  MathPaper.Word.Graph MathPaper.Word.Cert MathPaper.Word.Lift \
  MathPaper.Word.Chain MathPaper.Word.Alphabet \
  MathPaper.Word.Pipeline.Data MathPaper.Word.Pipeline.Check \
  MathPaper.Word.Pipeline.Compute MathPaper.Word.Pipeline.LiftArr \
  MathPaper.Word.Pipeline.Initial MathPaper.Word.Pipeline.Main \
  MathPaper.Word.FiveHalves MathPaper.Word.SevenFifths MathPaper.Word
\end{Verbatim}

The \nolinkurl{FiveHalves} and \nolinkurl{SevenFifths} targets rerun the two finite
\nolinkurl{native_decide} evaluations when built from a fresh project cache.
Then run the final targets, aggregate build, and \nolinkurl{Audit.lean} axiom
display as instructed in the \href{REPRODUCE.md}{reproduction guide}. Compare the dependencies with
the declared axiom boundary in the \href{supplement/FORMALIZATION.md}{formalization guide}. The additional evaluation
axioms are the trust boundary described in Section 9.4; these evaluations
are not kernel-only computations. Section 9.4 also specifies the scope of
the completed independent rebuilds and type and axiom checks.

% Markdown AST block 397: Appendix B. Recurrent carries
\Needspace{6\baselineskip}\section{Recurrent carries}\label{app:B}

% Markdown AST block 398: Para
This appendix proves Theorem 46 and the integer copying condition for
the family \((6t+3)/(6t+1)\). Its results are general lemmas about
integer recurrences and finite permutations; they do not depend on
Sections 3--9 and are not used there.

% Markdown AST block 399: B.1 Recurrent words
\Needspace{6\baselineskip}
\subsection{Recurrent words}\label{sec:B.1}

% Markdown AST block 400: Para
Recall from Section 10 that an infinite word \(w=w_0w_1\cdots\) over a
finite alphabet is \emph{recurrent} if for every \(L\ge1\) and \(B\ge1\)
there is \(n\ge B\) with
\[
w[n:n+L]=w[0:L],
\]
where
\[
w[i:j]=w_iw_{i+1}\cdots w_{j-1}.
\]
Equivalently, every finite prefix
occurs infinitely often, or again every factor that occurs at all occurs
infinitely often (a factor lies in some prefix).

% Markdown AST block 401: Lemma 48 (tails).
\begin{samepage}
\begin{lemma}[tails]\leavevmode\label{thm:48}
If \(w\) is recurrent, then so is \(w[k:]\) for
every \(k\ge0\).
\label{thm-end:48}
\end{lemma}
\end{samepage}

% Markdown AST block 402: Proof.
\begin{proof}\leavevmode\label{proof:48}
Given \(L,B\ge1\), apply recurrence of \(w\) to the prefix
length \(k+L\) and the bound \(B\): there is \(n\ge B\) with
\[
w[n:n+k+L]=w[0:k+L].
\]
Comparing positions \(k\) onward gives
\[
w[n+k:n+k+L]=w[k:k+L],
\]
that is,
\[
w'[n:n+L]=w'[0:L]
\]
for
\(w'=w[k:]\). \qedhere
\end{proof}

% Markdown AST block 403: Para
Uniform recurrence is not required. For example, let \(w_0=0\) and
\[
w_{k+1}=w_k\,1\,0^{k+1}\,1\,w_k;
\]
each \(w_k\) is a prefix of
\(w_{k+1}\), so a limit word \(w_\infty\) exists, in which every \(w_k\)
occurs infinitely often (it occurs in \(w_{j+1}\) at position
\(|w_j|+j+3\) for every \(j\ge k\)), while the runs of zeros are
unbounded and \(1\) occurs infinitely often, so \(w_\infty\) is not
eventually periodic. Conversely, the word \(1\,0\,1\,00\,1\,000\cdots\),
in which the runs of zeros have lengths \(1,2,3,\ldots\), has no
recurrent tail: a prefix of any tail that contains two letters \(1\)
contains a factor \(1\,0^k\,1\) which, the run lengths being strictly
increasing, occurs only once in the whole word. This word is generated by
a one-vertex graph with two self-loops, so confinement of a walk to a
strongly connected component does not imply recurrence of its output.

% Markdown AST block 404: Para
Neither example is claimed to be the carry sequence of a true orbit.

% Markdown AST block 405: B.2 The finite-prefix copying lemma
\Needspace{6\baselineskip}
\subsection{The finite-prefix copying lemma}\label{sec:B.2}

% Markdown AST block 406: Lemma 49 (finite-prefix copying).
\begin{samepage}
\begin{lemma}[finite-prefix copying]\leavevmode\label{thm:49}
Let \(X\) be a set with \(m\)
elements and let each letter \(c\) act by a permutation \(f_c\) of
\(X\). Let \(w\) be a finite word, \(t_1,\ldots,t_m\) positive integers,
\(S_0=0\), \(S_k=t_1+\cdots+t_k\), and \(L\ge0\), such that
\[
w[t_k:\,t_k+S_{k-1}+L]=w[0:\,S_{k-1}+L]\qquad(k=1,\ldots,m)
\tag{31}\label{eq:31}
\]
and \(|w|\ge S_m+L\). Put
\[
G_n=f_{w_{n-1}}\circ\cdots\circ f_{w_0}
\]
(\(G_0=\mathrm{id}\); maps act from the right, so \(G_n\) applies the
first letter first). Then for every \(x\in X\) there are
\(0\le i<j\le m\) such that
\[
n=t_{i+1}+\cdots+t_j>0
\]
satisfies
\[
G_n(x)=x\qquad \text{and}\qquad w[n:n+L]=w[0:L].
\]
\label{thm-end:49}
\end{lemma}
\end{samepage}

% Markdown AST block 407: Proof.
\begin{proof}\leavevmode\label{proof:49}
\emph{(1) Composition rule.} If \(0\le s\le S_{k-1}\), then
\[
G_{t_k+s}=G_s\circ G_{t_k}:
\]
indeed
\[
G_{t_k+s}=f_{w_{t_k+s-1}}\circ\cdots\circ f_{w_{t_k}}\circ G_{t_k},
\]
and by (31), \(w_{t_k+i}=w_i\) for \(0\le i<s\), so the left factor is
\(G_s\).

% Markdown AST block 408: Para
\emph{(2) Copies compose.} For indices \(i_1<\cdots<i_h\) and
\(n=t_{i_1}+\cdots+t_{i_h}\),
\[
\begin{aligned}
G_n&=G_{t_{i_1}}\circ G_{t_{i_2}}\circ\cdots\circ G_{t_{i_h}},\\
w[n:n+L]&=w[0:L].
\end{aligned}
\]
By induction on \(h\); \(h=0\) is \(G_0=\mathrm{id}\). For \(h\ge1\)
put
\[
n'=t_{i_1}+\cdots+t_{i_{h-1}}
\]
and \(k=i_h\); since all
\(i_j<k\), \(n'\le S_{k-1}\). By (1),
\[
G_n=G_{t_k+n'}=G_{n'}\circ
G_{t_k},
\]
and by induction
\[
G_{n'}=G_{t_{i_1}}\circ\cdots\circ
G_{t_{i_{h-1}}}.
\]

% Markdown AST block 409: Para
For the word, \(n'+L\le S_{k-1}+L\), so (31) gives
\[
w[t_k+n':t_k+n'+L]=w[n':n'+L]=w[0:L].
\]

% Markdown AST block 410: Para
\emph{(3) Pigeonhole.} Put \(H_0=\mathrm{id}\) and
\[
H_k=G_{t_1}\circ\cdots\circ G_{t_k};
\]
by (2) applied to
\(\{1,\ldots,k\}\), \(H_k=G_{S_k}\), the actual state at time \(S_k\).

% Markdown AST block 411: Para
Among the \(m+1\) values \(H_0(x),\ldots,H_m(x)\) in the \(m\)-element
set \(X\) two coincide: \(H_i(x)=H_j(x)\) with \(i<j\).

% Markdown AST block 412: Para
With
\[
B=G_{t_{i+1}}\circ\cdots\circ G_{t_j}
\]
we have \(H_j=H_i\circ B\),
so
\[
H_i(x)=H_i(B(x)),
\]
and injectivity of \(H_i\) (a composition of
permutations) gives \(B(x)=x\).

% Markdown AST block 413: Para
By (2) applied to \(\{i+1,\ldots,j\}\),
\[
\begin{aligned}
&B=G_n\qquad \text{with}\qquad n=t_{i+1}+\cdots+t_j>0\qquad \text{and}\\
&w[n:n+L]=w[0:L].
\end{aligned}\qedhere
\]
\end{proof}

% Markdown AST block 414: Para
Injectivity is used exactly once, in step (3); the conclusion fails for
non-injective actions. Applying the pigeonhole principle directly to the
states along the orbit would only give a return to \emph{some} state, not to
the initial state; the copies align the compositions so that the common
left factor can be cancelled.

% Markdown AST block 415: Theorem 50 (residue recurrence).
\begin{samepage}
\begin{theorem}[residue recurrence]\leavevmode\label{thm:50}
Let integer sequences
\((Q_n)_{n\ge0}\), \((c_n)_{n\ge0}\) satisfy
\[
bQ_{n+1}=aQ_n+c_n,
\]
let
the word \(c=c_0c_1\cdots\) be recurrent, and let \(M\ge2\) with
\(\gcd(M,ab)=1\). Then for every \(B\ge1\) and \(L\ge0\) there is
\(n\ge B\) with
\[
Q_n\equiv Q_0\pmod M\qquad \text{and}\qquad c[n:n+L]=c[0:L].
\]
\label{thm-end:50}
\end{theorem}
\end{samepage}

% Markdown AST block 416: Proof.
\begin{proof}\leavevmode\label{proof:50}
Let \(X=\mathbb Z/M\) and
\[
f_c(x)=b^{-1}(ax+c);
\]
since
\(a,b\) are units modulo \(M\), each \(f_c\) is an affine permutation,
and
\[
Q_n\equiv G_n(Q_0).
\]

% Markdown AST block 417: Para
Choose copy times successively: for
\(k=1,\ldots,M\), the required prefix length \(S_{k-1}+L\) is finite, so
by recurrence (when this length is positive; otherwise any \(t_k\ge B\))
there is \(t_k\ge B\) satisfying (31), \(S_{k-1}\) being determined by
\(t_1,\ldots,t_{k-1}\).

% Markdown AST block 418: Para
Lemma 49 with \(m=M\) and \(x=Q_0\bmod M\)
gives
\[
n=t_{i+1}+\cdots+t_j\ge t_j\ge B,
\]
\(G_n(Q_0)\equiv Q_0\), and
the return of the prefix. \qedhere
\end{proof}

% Markdown AST block 419: Para
The hypothesis \(\gcd(M,ab)=1\) cannot be dropped. If a prime
\(p\mid M\) divides \(a\), then \(f_c(x)\equiv b^{-1}c\pmod p\) is
constant, not injective, and \(Q_n\bmod p\) is determined by \(c_{n-1}\)
alone; if \(Q_0\bmod p\) is not of the form \(b^{-1}c\) for a letter
\(c\) of the word (for instance \(p\mid Q_0\) and no letter is divisible
by \(p\)), then \(Q_n\equiv Q_0\pmod p\) never holds for \(n\ge1\),
however recurrent the word is.

% Markdown AST block 420: B.3 Proof of Theorem 46
\Needspace{6\baselineskip}
\subsection{Proof of Theorem 46}\label{sec:B.3}

% Markdown AST block 421: Proof of Theorem 46.
\begin{proof}[Proof of Theorem 46]\leavevmode\label{proof:46}
Let \(c[j_0:]\) be recurrent.

% Markdown AST block 422: Para
\emph{Case 1: some \(j\ge j_0\) has \(Q_j>1\) and \(\gcd(Q_j,ab)=1\).} By
Lemma 48, \(c[j:]\) is recurrent. Apply Theorem 50 to the sequences
\(Q'_n=Q_{j+n}\), \(c'_n=c_{j+n}\) with \(M=Q_j\ge2\) and \(L=0\): for
every \(B\) there is \(n\ge B\) with
\[
Q_{j+n}\equiv Q_j\equiv0\pmod{Q_j}.
\]
By (7),
\[
Q_{j+n}>Q_j>1
\]
for
all large \(n\), so \(Q_j\) is a proper divisor and \(Q_{j+n}\) is
composite; there are infinitely many such \(n\). This also proves the
last assertion of the theorem.

% Markdown AST block 423: Para
\emph{Case 2: no such \(j\).} By (7), \(Q_j\le1\) for only finitely many
\(j\), so
\[
\gcd(Q_j,ab)>1
\]
for all large \(j\). A prime
\(p\mid\gcd(Q_j,ab)\) satisfies \(p\le ab\), so once \(Q_j>ab\) we have
\(Q_j=pv\) with \(v>1\). \qedhere
\end{proof}

% Markdown AST block 424: Para
Case 2 does not use recurrence, but the conclusion of Case 1, that
\(Q_j\) divides infinitely many later terms, does; in general
\(Q_j\mid Q_n\) need not occur. This is why the recurrence hypothesis is
part of the statement.

% Markdown AST block 425: Corollary 51 (refutation form).
\begin{samepage}
\begin{corollary}[refutation form]\leavevmode\label{thm:51}
If all terms \(Q_n\) with
\(n\ge n_0\) are prime, then the carry word has no recurrent tail.
\label{thm-end:51}
\end{corollary}
\end{samepage}

% Markdown AST block 426: Proof.
\begin{proof}\leavevmode\label{proof:51}
Suppose \(c[j_0:]\) is recurrent. By (7) choose
\(j\ge\max(j_0,n_0)\) with \(P=Q_j>a\). Then \(P\) is a prime larger
than \(a>b\), so \(\gcd(P,ab)=1\), and Theorem 50 with \(M=P\) gives
arbitrarily large \(n\) with \(P\mid Q_n\) and \(Q_n>P\), hence
composite, a contradiction. \qedhere
\end{proof}

% Markdown AST block 427: Para
The condition \(P>a\) serves only to guarantee \(\gcd(P,ab)=1\).
Theorem 46 excludes recurrent carry tails on an eventually-all-prime
orbit; it does not exclude orbits all of whose carry tails are
non-recurrent (Section B.1 shows that such words exist as words), and
this remaining case is open (Section 11).

% Markdown AST block 428: B.4 The integer copying condition for the family \((6t+3)/(6t+1)\)
\Needspace{6\baselineskip}
\subsection{\texorpdfstring{The integer copying condition for the family \((6t+3)/(6t+1)\)}{The integer copying condition for the family (6t+3)/(6t+1)}}\label{sec:B.4}

% Markdown AST block 429: Para
Let
\[
\begin{aligned}
&t\ge1,\qquad a=6t+3,\qquad b=6t+1,\\
&\delta=t\bmod2,\qquad H=3t-1-\delta
\end{aligned}
\]
(an odd integer), and
\[
\begin{aligned}
T_t(Q)&=Q+2\Bigl\lfloor\frac{Q+3t+1}{b}\Bigr\rfloor,\\
F_t(z)&=\Bigl\lceil\frac{az-\delta}{b}\Bigr\rceil,\\
e&=bF_t(z)-az+\delta .
\end{aligned}
\tag{32}\label{eq:32}
\]
Appendix C shows that a tail of a true orbit all of whose terms are
coprime to \(6\) follows the map \(T_t\); here we study \(T_t\) as an
integer map.

% Markdown AST block 430: Lemma 52 (conjugation and digits).
\begin{samepage}
\begin{lemma}[conjugation and digits]\leavevmode\label{thm:52}
For an odd integer \(Q=2z+H\):
(a) \(T_t(Q)=2F_t(z)+H\), and \(T_t\) preserves parity; (b)
\[
0\le e\le b-1\qquad \text{and}\qquad e\equiv\delta-az\pmod b;
\]
(c) the carry
\(c=bT_t(Q)-aQ\) equals \(2(e-(3t-1))\); (d)
\[
F_t(s+bv)=F_t(s)+av
\]
for all integers \(s,v\).
\label{thm-end:52}
\end{lemma}
\end{samepage}

% Markdown AST block 431: Proof.
\begin{proof}\leavevmode\label{proof:52}
(a)
\[
Q+3t+1=2z+(b-1)-\delta,
\]
so
\[
\lfloor(Q+3t+1)/b\rfloor=1+\lfloor(2z-1-\delta)/b\rfloor.
\]
With
\(a=b+2\),
\[
\begin{aligned}
&(az-\delta)/b=z+(2z-\delta)/b,\qquad\text{and}\\
&\lceil n/b\rceil=\lfloor(n-1)/b\rfloor+1
\end{aligned}
\]
for integers \(n\) gives
\[
F_t(z)=z+1+\lfloor(2z-1-\delta)/b\rfloor.
\]
Hence
\[
2F_t(z)+H=2z+H+2+2\lfloor(2z-1-\delta)/b\rfloor=T_t(Q).
\]
Parity is
preserved since \(T_t(Q)-Q\) is even.

% Markdown AST block 432: Para
\textbf{(b)} is the definition of the
ceiling.

% Markdown AST block 433: Para
\textbf{(c)}
\[
c=2(bF_t(z)-az)+(b-a)H=2(e-\delta)-2H=2(e-(3t-1)).
\]

% Markdown AST block 434: Para
\textbf{(d)}
\[
\lceil(a(s+bv)-\delta)/b\rceil=\lceil(as-\delta)/b+av\rceil.\qedhere
\]
\end{proof}

% Markdown AST block 435: Proposition 53 (digit words and residues).
\begin{samepage}
\begin{proposition}[digit words and residues]\leavevmode\label{thm:53}
Let \(L\ge1\). The map
sending an integer \(z_0\) to the first \(L\) digits
\(e_0\cdots e_{L-1}\) of its \(F_t\)-orbit induces a bijection
\[
\mathbb Z/b^L\to\{0,\ldots,b-1\}^L.
\]
Consequently, for odd
\(Q_0,Q_0'\) and their \(T_t\)-orbits with carry words \(c,c'\),
\[
c[0:L]=c'[0:L]\iff Q_0\equiv Q_0'\pmod{b^L},
\]
and in particular, along one orbit,
\[
c[k:k+L]=c[0:L]\iff Q_k\equiv Q_0\pmod{b^L}.
\]
\label{thm-end:53}
\end{proposition}
\end{samepage}

% Markdown AST block 436: Proof.
\begin{proof}\leavevmode\label{proof:53}
Induction on \(L\). For \(L=1\),
\[
e_0\equiv\delta-az_0\pmod b
\]
with \(\gcd(a,b)=1\) (\(a=b+2\), \(b\)
odd), so \(z_0\bmod b\mapsto e_0\) is a bijection.

% Markdown AST block 437: Para
For the step, if
\(z_0'=z_0+b^{L+1}v\), then \(e_0=e_0'\) and by Lemma 52(d)
\[
z_1'=F_t(z_0')=z_1+ab^Lv\equiv z_1\pmod{b^L},
\]
so the remaining
\(L\) digits agree by induction (the map is well defined).

% Markdown AST block 438: Para
Conversely,
if the digit words agree, then \(e_0\) determines \(z_0\bmod b\); write
\[
\begin{aligned}
&z_0=s+bq,\qquad z_0'=s+bq',\qquad\text{so}\\
&z_1=F_t(s)+aq,\qquad z_1'=F_t(s)+aq';
\end{aligned}
\]
by induction
\[
z_1\equiv z_1'\pmod{b^L},
\]
and since \(a\) is a unit
modulo \(b^L\), \(q\equiv q'\pmod{b^L}\), that is,
\[
z_0\equiv z_0'\pmod{b^{L+1}}
\]
(injectivity).

% Markdown AST block 439: Para
Both sets have
\(b^{L+1}\) elements, so the map is a bijection.

% Markdown AST block 440: Para
For carries, Lemma
52(c) makes \(c_n\leftrightarrow e_n\) bijective, and since \(b\) is
odd, \(Q=2z+H\) gives
\[
Q\equiv Q'\pmod{b^L}\iff z\equiv z'\pmod{b^L}.\qedhere
\]
\end{proof}

% Markdown AST block 441: Proposition 54 (\(P\)-stage congruence certificate).
\begin{samepage}
\begin{proposition}[\(P\)-stage congruence certificate]\leavevmode\label{thm:54}
Let \(P>a\)
be prime, \(Q_n=T_t^n(P)\), and let positive integers
\(k_1,\ldots,k_P\) with \(s_0=0\), \(s_i=s_{i-1}+k_i\) satisfy
\[
Q_{k_i}\equiv P\pmod{b^{s_{i-1}}}\qquad(i=1,\ldots,P).
\tag{33}\label{eq:33}
\]
Then \(P\mid Q_n\) for some \(0<n\le s_P\), and \(Q_n>P\), so \(Q_n\) is
composite.
\label{thm-end:54}
\end{proposition}
\end{samepage}

% Markdown AST block 442: Proof.
\begin{proof}\leavevmode\label{proof:54}
By Proposition 53 with \(L=s_{i-1}\) and \(k=k_i\), (33)
gives
\[
c[k_i:k_i+s_{i-1}]=c[0:s_{i-1}],
\]
which is (31) with
\(t_i=k_i\) and \(L=0\).

% Markdown AST block 443: Para
Apply Lemma 49 on \(X=\mathbb Z/P\) with
\[
f_c(x)=b^{-1}(ax+c)
\]
(permutations, since \(P>a>b\) gives
\(\gcd(P,ab)=1\)) to \(x=0\): some
\(n=k_{i+1}+\cdots+k_j\le s_P\)
has
\[
Q_n\equiv G_n(0)=0.
\]

% Markdown AST block 444: Para
If \(Q\ge3t\) then
\[
\begin{aligned}
&\lfloor(Q+3t+1)/b\rfloor\ge1\qquad\text{and}\\
&T_t(Q)\ge Q+2;
\end{aligned}
\]
as \(T_t\) is
nondecreasing and \(P>a>3t\), the orbit is strictly increasing and
\(Q_n>P\). \qedhere
\end{proof}

% Markdown AST block 445: Para
\emph{Application to true orbits.} Let \(\lfloor\xi\rfloor=P\). If all terms
\(\lfloor\xi r^n\rfloor\), \(0\le n\le s_P\), are coprime to \(6\), then
by Proposition 55 (Appendix C) they coincide with the \(T_t\)-orbit of
\(P\), and Proposition 54 gives a composite term at some
\(n\le s_P\). Otherwise some term \(\ge P>3\) is divisible by \(2\) or
\(3\) and is composite. In either case a composite term occurs at a time
\(\le s_P\), and no interval-consistency check is needed. The existence
of such certificates for all \(t\) and \(P\) is not proved.

% Markdown AST block 446: Para
\emph{Example.} For \(9/7\) (\(t=1\)) and \(\xi=11\), the carries of the true
orbit \(\lfloor11(9/7)^n\rfloor\) satisfy \(c_0=c_{43}=-1\), so the copy condition (31) holds with
\(t_1=1\), \(t_2=43\), and \(L=0\); step (1) of the proof of
Lemma 49 then gives \(G_{44}=G_1\circ G_{43}\) on
\(\mathbb Z/11\), and one checks directly that the residue returns:
\[
11\mid Q_{44}=697829=11\cdot63439.
\]
(Lemma 49 itself would require
eleven copy times; here the return is verified directly.)

% Markdown AST block 447: Para
This certificate lives on the true orbit
rather than on the \(T_1\)-orbit (since \(\lfloor99/7\rfloor=14\ne T_1(11)=15\)), and \(P=11>a=9\) is admissible for Lemma 49 itself because
\(\gcd(11,63)=1\).

% Markdown AST block 448: Para
\emph{Remark (quantitative form).} Let \(R(L)\) be the first positive return
time of the prefix \(w[0:L]\). If \(R(L)\le KL\) for all \(L\ge1\), then
copies with \(L=1\) can be chosen with
\[
t_k=R(S_{k-1}+1)\le
K(S_{k-1}+1),
\]
so
\[
S_k+1\le(K+1)^k,
\]
and a return modulo \(M\) to the
initial state occurs by time \((K+1)^M-1\). No such \(K\) is proved for
the family, and this bound is not used.

% Markdown AST block 449: Appendix C. Non-branching families and the exact integer halting problem
\Needspace{6\baselineskip}\section{Non-branching families and the exact integer halting problem}\label{app:C}

% Markdown AST block 450: Para
This appendix concerns the family \(r_t=(6t+3)/(6t+1)\), \(t\ge1\),
that is \(9/7,15/13,21/19,\ldots\) (the members \(j=3t-2\) of the family
\((2j+7)/(2j+5)\) of Corollary 40). It proves that a tail avoiding the
primes \(2\) and \(3\) follows the integer map \(T_t\) of (32), that the
all-coefficient statement for a fixed \(t\) is equivalent to the halting
of an exact integer program, and it constructs further families with a
unique prime successor. None of this proves halting for the family, and
no all-coefficient theorem for any \(r_t\) is asserted.

% Markdown AST block 451: C.1 The non-branching lemma
\Needspace{6\baselineskip}
\subsection{The non-branching lemma}\label{sec:C.1}

% Markdown AST block 452: Proposition 55 (non-branching).
\begin{samepage}
\begin{proposition}[non-branching]\leavevmode\label{thm:55}
Let
\[
\begin{aligned}
&t\ge1,\qquad a=6t+3,\qquad b=6t+1,\\
&x>0,\qquad Q=\lfloor x\rfloor,\qquad Q'=\lfloor rx\rfloor.
\end{aligned}
\]
If \(Q\) is odd and \(\gcd(Q',6)=1\), then \(Q'=T_t(Q)\).
\label{thm-end:55}
\end{proposition}
\end{samepage}

% Markdown AST block 453: Proof.
\begin{proof}\leavevmode\label{proof:55}
The carry \(c=bQ'-aQ\) satisfies
\[
-6t\le c\le6t+2
\]
by (6).
Since \(a,b,Q,Q'\) are odd, \(c=2s\) is even with
\[
-3t\le s\le3t+1.
\]

% Markdown AST block 454: Para
As \(3\mid a\), \(3\nmid b\), \(3\nmid Q'\), we have
\[
c\equiv bQ'\not\equiv0\pmod3,
\]
so \(3\nmid s\); this excludes
\(s=-3t\), leaving
\[
-3t+1\le s\le3t+1,
\]
an interval of exactly
\(b=6t+1\) integers.

% Markdown AST block 455: Para
Modulo \(b\), \(a\equiv2\), so
\(2s=c\equiv-2Q\) and
\[
s\equiv-Q\pmod b;
\]
the interval contains
exactly one integer in each residue class, so \(s\) is unique.

% Markdown AST block 456: Para
With
\(m=\lfloor(Q+3t+1)/b\rfloor\),
\[
0\le Q+3t+1-bm\le6t
\]
gives
\[
Q-3t+1\le bm\le Q+3t+1,
\]
so \(s=bm-Q\) lies in the interval and is
\(\equiv-Q\). Hence
\[
Q'=(aQ+2s)/b=(a-2)Q/b+2m=Q+2m=T_t(Q).\qedhere
\]
\end{proof}

% Markdown AST block 457: Para
Only the oddness of \(Q\) and \(\gcd(Q',6)=1\) are used; \(3\nmid Q\)
is not needed. The conclusion is uniqueness of the candidate: if
\(T_t(Q)\) is divisible by \(3\), no admissible successor exists.

% Markdown AST block 458: C.2 Conjugation to a ceiling map and digit restriction
\Needspace{6\baselineskip}
\subsection{Conjugation to a ceiling map and digit restriction}\label{sec:C.2}

% Markdown AST block 459: Para
By Lemma 52, the map \(T_t\) on odd integers \(Q=2z+H\) is conjugate to
\[
F_t(z)=\lceil(az-\delta)/b\rceil,
\]
with digits
\[
e=bF_t(z)-az+\delta\in\{0,\ldots,b-1\}
\]
and carries
\[
c=2(e-(3t-1)).
\]
For \(t=2u\) even, \(\delta=0\), \(H=6u-1\), and
\(F_t(z)=\lceil az/b\rceil\) is the ordinary ceiling map.

% Markdown AST block 460: Lemma 56 (digit restriction).
\begin{samepage}
\begin{lemma}[digit restriction]\leavevmode\label{thm:56}
(a)
\[
Q'=T_t(Q)\equiv2e-1\pmod3;
\]
hence \(3\mid Q'\) if and only if \(e\equiv2\pmod3\). (b) If \(Q\) and
\(Q'\) are primes larger than \(a\), then
\[
e\in D_t=\{0\le e\le6t:\ \gcd(e-(3t-1),(6t+3)(6t+1))=1\}.
\]
(c)
\[
d_t=|D_t|\le4t+1<b.
\]
\label{thm-end:56}
\end{lemma}
\end{samepage}

% Markdown AST block 461: Proof.
\begin{proof}\leavevmode\label{proof:56}
(a) \(e\equiv F_t(z)+\delta\pmod3\) since \(b\equiv1\),
\(a\equiv0\); and \(H\equiv-1-\delta\), so
\[
Q'=2F_t(z)+H\equiv2(e-\delta)-1-\delta\equiv2e-1\pmod3.
\]

% Markdown AST block 462: Para
\textbf{(b)} By
Lemma 29, \(\gcd(c,ab)=1\), and \(ab\) is odd, so
\[
\gcd(e-(3t-1),ab)=1.
\]

% Markdown AST block 463: Para
\textbf{(c)} Since \(3t-1\equiv2\pmod3\), every
\(e\equiv2\pmod3\) is excluded, and there are \(2t\) such \(e\) in
\([0,6t]\); so
\[
d_t\le6t+1-2t.\qedhere
\]
\end{proof}

% Markdown AST block 464: Table 9.
% Caption is rendered at the start of the immediately following table.

% Markdown AST block 465: Table 9
\begingroup
\small
\setlength{\tabcolsep}{4pt}
\renewcommand{\arraystretch}{1.15}
\setlength{\LTcapwidth}{\textwidth}
\begin{longtable}{@{}>{\raggedright\arraybackslash}p{\dimexpr0.08\dimexpr\linewidth-16pt\relax} >{\raggedright\arraybackslash}p{\dimexpr0.12\dimexpr\linewidth-16pt\relax} >{\raggedright\arraybackslash}p{\dimexpr0.8\dimexpr\linewidth-16pt\relax}@{}}
\caption{Admissible digit sets for \(t\le3\).}\label{tab:9}\\
\toprule
\(t\) & \(a/b\) & \(D_t\) \\
\midrule
\endfirsthead
\toprule
\(t\) & \(a/b\) & \(D_t\) \\
\midrule
\endhead
\bottomrule
\endlastfoot
1 & \(9/7\) & \(\{0,\allowbreak{}1,\allowbreak{}3,\allowbreak{}4,\allowbreak{}6\}\) \\*
2 & \(15/13\) & \(\{1,\allowbreak{}3,\allowbreak{}4,\allowbreak{}6,\allowbreak{}7,\allowbreak{}9,\allowbreak{}12\}\) \\*
3 & \(21/19\) & \(\{0,\allowbreak{}3,\allowbreak{}4,\allowbreak{}6,\allowbreak{}7,\allowbreak{}9,\allowbreak{}10,\allowbreak{}12,\allowbreak{}13,\allowbreak{}16,\allowbreak{}18\}\) \\
\end{longtable}
\endgroup

% Markdown AST block 466: Para
\emph{Remark (real consistency).} For a true orbit \(x_n=2z_n+H+\theta_n\),
so
\[
x_n/2-z_n\in[H/2,(H+1)/2).
\]
From \(bz_{n+1}=az_n+e_n-\delta\),
\[
r^{-(n+1)}z_{n+1}-r^{-n}z_n=r^{-(n+1)}(e_n-\delta)/b
\]
with
\(|e_n-\delta|\le b\), so
\[
\Lambda=\lim r^{-n}z_n
\]
exists and
\[
z_n=\Lambda r^n-\sum_{j\ge0}(e_{n+j}-\delta)b^j/a^{j+1};
\]
for
\(z_0>b/2\), \(\Lambda>0\), and for a true orbit \(\xi=2\Lambda\) is
determined. Membership of all digits in \(D_t\) implies neither the
interval condition of width \(1/2\) nor the primality of the terms.

% Markdown AST block 467: Remark.
\begin{remark}
The iterates of the ceiling map \(x\mapsto\lceil(p/q)x\rceil\)
on integers are the object of a conjecture on equidistribution of
residues modulo \(q^k\) in \hyperref[ref-aev26]{AEV26, Conjecture 1.6}; that
paper concerns integer seeds and states conjectures, not theorems, and
nothing here depends on it.
\end{remark}

% Markdown AST block 468: C.3 The exact integer halting program
\Needspace{6\baselineskip}
\subsection{The exact integer halting program}\label{sec:C.3}

% Markdown AST block 469: Para
Fix \(t\ge1\) and an initial prime \(P>a\). Starting from
\((Q,L,U,B)=(P,0,1,1)\), repeat:
\[
\begin{aligned}
Q^{+}&=T_t(Q),\quad c=bQ^{+}-aQ,\quad B^{+}=bB,\\
L^{+}&=\max(0,aL-cB),\\
U^{+}&=\min(B^{+},aU-cB),
\end{aligned}
\]
and \emph{halt} if \(Q^{+}\) is composite or \(L^{+}\ge U^{+}\); otherwise
replace \((Q,L,U,B)\) by \((Q^{+},L^{+},U^{+},B^{+})\) and continue.
Primality is decided exactly. Reaching a bound on time, memory, or
iterations is recorded as undetermined and is not a halt.

% Markdown AST block 470: Proposition 57 (exactness).
\begin{samepage}
\begin{proposition}[exactness]\leavevmode\label{thm:57}
After \(n\) accepted steps, \(B=b^n\) and
\([L/B,U/B)\) is exactly the set of fractional parts \(\theta_n\) of
those \(\xi\in[P,P+1)\) whose orbit has integer parts \(P=Q_0,\ldots,Q_n\).
In particular \(L^{+}\ge U^{+}\) means that the finite sequence
\(Q_0,\ldots,Q_{n+1}\) is realized by no real \(\xi\).
\label{thm-end:57}
\end{proposition}
\end{samepage}

% Markdown AST block 471: Proof.
\begin{proof}\leavevmode\label{proof:57}
The map
\[
\theta\mapsto\theta'=(a\theta-c)/b
\]
is an
increasing affine bijection, so the image of \([L/B,U/B)\) is
\[
[(aL-cB)/(bB),(aU-cB)/(bB)),
\]
whose intersection with \([0,1)\) is
\[
[L^{+}/B^{+},U^{+}/B^{+}).
\]

% Markdown AST block 472: Para
Initially \(\theta_0\) ranges over
\([0,1)\). Inductively, \(\theta_{n+1}\) lies in the image intersected
with \([0,1)\) if and only if the \(\theta_n\) producing it lies in the
previous set. \qedhere
\end{proof}

% Markdown AST block 473: C.4 Equivalence with the all-coefficient statement
\Needspace{6\baselineskip}
\subsection{Equivalence with the all-coefficient statement}\label{sec:C.4}

% Markdown AST block 474: Theorem 58 (halting equivalence).
\begin{samepage}
\begin{theorem}[halting equivalence]\leavevmode\label{thm:58}
For fixed \(t\ge1\) the following
are equivalent.

% Markdown AST block 475: OrderedList
\begin{enumerate}
\def\labelenumi{(\Alph{enumi})}
\item
  For every \(\xi>0\), the sequence \(\lfloor\xi r_t^n\rfloor\)
  contains infinitely many composite terms.
\item
  The program of Section C.3 halts for every initial prime \(P>a\).
\end{enumerate}
\label{thm-end:58}
\end{theorem}
\end{samepage}

% Markdown AST block 476: Proof.
\begin{proof}\leavevmode\label{proof:58}
(B)\(\Rightarrow\)(A). If (A) fails, there are \(\xi\) and \(n_0\) with
\(Q_n\) prime for all \(n\ge n_0\). By (7) choose \(n_1\ge n_0\) with
\(Q_{n_1}>a\) and replace \(\xi\) by \(\xi r^{n_1}\); then
\(P=\lfloor\xi\rfloor>a\) is prime and all terms are primes \(>a>3\),
coprime to \(6\). By Proposition 55,
\[
Q_{n+1}=T_t(Q_n)
\]
for all
\(n\), so the program's \(Q\) coincides with the true terms and never
becomes composite, and the true \(\theta_n\) lies in every interval by
Proposition 57, so no interval is empty. The program does not halt,
contradicting (B).

% Markdown AST block 477: Para
(A)\(\Rightarrow\)(B). Suppose the program does not halt for some \(P>a\). Then all
\(Q_n=T_t^n(P)\) are prime (a non-composite integer \(>a\ge9\) is
prime) and all intervals are nonempty, so by Proposition 57 there is,
for each \(n\), some \(\xi_n\in[P,P+1)\) with
\[
\lfloor\xi_nr^j\rfloor=Q_j
\]
for \(0\le j\le n\).

% Markdown AST block 478: Para
Let
\[
\begin{aligned}
&I_j=[Q_jr^{-j},(Q_j+1)r^{-j}]\qquad \text{and}\\
&J_n=I_0\cap\cdots\cap I_n.
\end{aligned}
\]
Then
\[
\xi_n\in J_n\ne\varnothing,
\]
\(J_{n+1}\subseteq J_n\), and each
\(J_n\) is a closed subset of the compact interval \(I_0=[P,P+1]\); by
Cantor's intersection theorem \(\bigcap_nJ_n\ne\varnothing\), and since
the widths are at most \(r^{-n}\to0\), the intersection is a single
point \(\xi_*>0\). (With half-open intervals the finite intersection
property would not yield a common point; this is why closed intervals
are used and the endpoints are treated separately.)

% Markdown AST block 479: Para
Thus
\[
Q_n\le\xi_*r^n\le Q_n+1
\]
for all \(n\).

% Markdown AST block 480: Para
If \(\xi_*r^n=Q_n+1\) for
some \(n\), then
\[
\xi_*=(Q_n+1)b^n/a^n
\]
is rational, \(\xi_*=u/v\) in
lowest terms; \(\xi_*r^m\in\mathbb Z\) requires \(vb^m\mid ua^m\), hence
\(v\mid a^m\) and \(b^m\mid u\), which holds for only finitely many
\(m\) as \(u\ne0\).

% Markdown AST block 481: Para
So for all large \(n\),
\[
Q_n\le\xi_*r^n<Q_n+1,
\]
that is,
\[
Q_n=\lfloor\xi_*r^n\rfloor
\]
is
prime, contradicting (A). \qedhere
\end{proof}

% Markdown AST block 482: Para
Whether (B) holds for every \(t\) is open. The orbit of \(\xi_*\) may
differ from the \(T_t\)-orbit at finitely many indices; agreement of the
tails suffices.

% Markdown AST block 483: C.5 Removing prime forks for general bases
\Needspace{6\baselineskip}
\subsection{Removing prime forks for general bases}\label{sec:C.5}

% Markdown AST block 484: Para
Let \(a=b+d\), \(\gcd(a,b)=1\), \(1<a/b<2\) (so \(d<b\)). For
\(x\in[Q,Q+1)\) the possible values of \(\lfloor rx\rfloor\) are the
integers \(m\) for which \([m,m+1)\) meets \([aQ/b,a(Q+1)/b)\).
The latter interval has length \(a/b<2\), so the candidates consist of
at most three consecutive integers and include at most two odd integers,
necessarily of the form \(L,L+2\) when both occur.

% Markdown AST block 485: Lemma 59 (two odd candidates).
\begin{samepage}
\begin{lemma}[two odd candidates]\leavevmode\label{thm:59}
If \(Q\) is odd and both \(L\) and
\(L+2\) are candidates, then with \(C=L+1\),
\[
\begin{aligned}
aQ&=bC-k,\quad 1\le k<d,\quad k\equiv a\pmod2,\\
bL&=a(Q-1)+(d+k),\\
b(L+2)&=a(Q+1)-(d-k).
\end{aligned}
\]
\label{thm-end:59}
\end{lemma}
\end{samepage}

% Markdown AST block 486: Proof.
\begin{proof}\leavevmode\label{proof:59}
\(L\) is a candidate iff \(aQ/b<L+1\) and \(L<a(Q+1)/b\);
\(L+2\) is a candidate iff \(aQ/b<L+3\) and \(L+2<a(Q+1)/b\).

% Markdown AST block 487: Para
Together,
\(aQ<bC\) and \(bC+b<aQ+a\), that is,
\[
bC-d<aQ<bC,
\]
so
\[
k=bC-aQ\in\{1,\ldots,d-1\}.
\]

% Markdown AST block 488: Para
As \(L\) is odd, \(C\) and \(bC\) are
even, and \(aQ\equiv a\pmod2\), so \(k\equiv a\pmod2\).

% Markdown AST block 489: Para
The identities
follow from
\[
\begin{aligned}
&bL=bC-b=aQ+k-b\qquad \text{and}\\
&b(L+2)=bC+b=aQ+k+b.
\end{aligned}\qedhere
\]
\end{proof}

% Markdown AST block 490: Corollary 60 (no prime fork).
\begin{samepage}
\begin{corollary}[no prime fork]\leavevmode\label{thm:60}
If for every \(k\) with \(1\le k<d\)
and \(k\equiv a\pmod2\) one has
\[
\gcd(k,b)>1,\qquad\text{or}\qquad \gcd(d+k,a)>1,\qquad\text{or}\qquad \gcd(d-k,a)>1,
\]
then there is no odd prime
\(Q>a\) with two prime candidates \(L,L+2>a\). Moreover, if
\(L,L+2>3\) are prime then \(3\mid C\), so for a prime \(Q>3\) a fork
requires
\[
3\mid a\iff3\mid k;
\]
values of \(k\) violating this
equivalence are also excluded.
\label{thm-end:60}
\end{corollary}
\end{samepage}

% Markdown AST block 491: Proof.
\begin{proof}\leavevmode\label{proof:60}
If \(p\mid\gcd(k,b)\) then \(p\mid aQ\), \(p\nmid a\),
\(p\mid Q\), and \(Q\) prime forces \(Q=p\le b<a\).

% Markdown AST block 492: Para
If
\(p\mid\gcd(d+k,a)\) then \(p\mid bL\), \(p\nmid b\), \(p\mid L\),
\(L=p\le a\).

% Markdown AST block 493: Para
If \(p\mid\gcd(d-k,a)\), similarly \(p\mid L+2\).

% Markdown AST block 494: Para
For
the last assertion, among \(L,L+1,L+2\) the multiple of \(3\) is
\(C=L+1\), and
\[
aQ=bC-k\equiv-k\pmod3
\]
with \(3\nmid Q\). \qedhere
\end{proof}

% Markdown AST block 495: Para
\emph{Example.} For \(a=462u+77\), \(b=462u+71\) (\(u\ge0\)): \(d=6\),
\(a\) odd, \(k\in\{1,3,5\}\); \(7\mid a\) and \(11\mid a\) give
\(\gcd(d+k,a)>1\) for \(k=1,5\); for \(k=3\), \(3\mid k\) but
\(a\equiv2\pmod3\), so the equivalence fails. The pairs are reduced
(\(\gcd(a,b)\mid6\), \(a\) odd, \(3\nmid a\)) with \(a<2b\). Hence a
prime \(Q>a\) has at most one prime successor in this family; the
surviving candidate need not be the larger one.

% Markdown AST block 496: C.6 Families with a prescribed prime-power difference
\Needspace{6\baselineskip}
\subsection{Families with a prescribed prime-power difference}\label{sec:C.6}

% Markdown AST block 497: Proposition 61 (families for every prime-power difference).
\begin{samepage}
\begin{proposition}[families for every prime-power difference]\leavevmode\label{thm:61}
Let
\(d=p^h\ge2\) and
\[
K_d=\{k:1\le k<d,\ k\equiv d+1\pmod2\}.
\]
For each
\(k\in K_d\) there is a prime \(\ell_k\ne p\) dividing \(d+k\); let
\(P_d\) be the product of the distinct \(\ell_k\). Then
\(\gcd(P_d,2d)=1\), and every \(a>2d\) with
\[
a\equiv0\pmod{P_d}\qquad \text{and}\qquad a\equiv d+1\pmod{2d}
\]
(an arithmetic progression modulo \(2dP_d\)),
with \(b=a-d\), satisfies \(\gcd(a,b)=1\), \(1<a/b<2\), and the
hypothesis of Corollary 60 in the form
\[
\gcd(d+k,a)>1
\]
for all
relevant \(k\).
\label{thm-end:61}
\end{proposition}
\end{samepage}

% Markdown AST block 498: Proof.
\begin{proof}\leavevmode\label{proof:61}
For \(k\in K_d\), \(d+k\equiv2d+1\) is odd and
\(d<d+k<2d\); there is no power of \(p\) in \((d,2d)\), so
\(d+k\ge3\) has a prime factor \(\ell_k\ne p\), necessarily odd.

% Markdown AST block 499: Para
Hence
\(P_d\) is odd and coprime to \(p\), so \(\gcd(P_d,2d)=1\) and the two
congruences are compatible.

% Markdown AST block 500: Para
From \(a\equiv1\pmod d\),
\[
\gcd(a,b)=\gcd(a,d)=1;
\]
\[
a>2d\iff a<2b;
\]
and \(a\equiv d+1\pmod2\)
makes \(k\equiv a\pmod2\) equivalent to \(k\in K_d\), with
\[
\ell_k\mid\gcd(d+k,a).\qedhere
\]
\end{proof}

% Markdown AST block 501: Table 10.
% Caption is rendered at the start of the immediately following table.

% Markdown AST block 502: Table 10
\begingroup
\small
\setlength{\tabcolsep}{4pt}
\renewcommand{\arraystretch}{1.15}
\setlength{\LTcapwidth}{\textwidth}
\begin{longtable}{@{}>{\raggedright\arraybackslash}p{\dimexpr0.08\dimexpr\linewidth-24pt\relax} >{\raggedright\arraybackslash}p{\dimexpr0.3\dimexpr\linewidth-24pt\relax} >{\raggedright\arraybackslash}p{\dimexpr0.3\dimexpr\linewidth-24pt\relax} >{\raggedright\arraybackslash}p{\dimexpr0.32\dimexpr\linewidth-24pt\relax}@{}}
\caption{Families from Proposition 61 for small \(d\).}\label{tab:10}\\
\toprule
\(d\) & \(a\) & \(b\) & Range \\
\midrule
\endfirsthead
\toprule
\(d\) & \(a\) & \(b\) & Range \\
\midrule
\endhead
\bottomrule
\endlastfoot
2 & \(12u+3\) & \(12u+1\) & \(u\ge1\) \\*
3 & \(30u+10\) & \(30u+7\) & \(u\ge0\) \\*
4 & \(280u+245\) & \(280u+241\) & \(u\ge0\) \\*
5 & \(210u+126\) & \(210u+121\) & \(u\ge0\) \\
\end{longtable}
\endgroup

% Markdown AST block 503: Para
For odd \(d\) the numerator \(a\) is even; Lemma 59 handles this through
the parity condition \(k\equiv a\pmod2\). These are families with a
uniform prime-successor rule, not families for which compositeness has
been established.

% Markdown AST block 504: Proposition 62 (ceiling form).
\begin{samepage}
\begin{proposition}[ceiling form]\leavevmode\label{thm:62}
In the setting of Proposition 61,
write \(b=2dh+1\), \(H=2h-1\), and \(Q=2z+H\). The largest odd integer
below \(r(Q+1)\) is \(2\lceil az/b\rceil+H\); hence, if an odd candidate
for the successor of \(Q\) exists, the largest one is
\[
Q'=2\lceil az/b\rceil+H,
\]
and a prime successor \(>a\) of a prime
\(Q>a\) can only be this value. With \(z'=\lceil az/b\rceil\) and
\(e=bz'-az\in[0,b)\),
\[
\begin{aligned}
\gcd(Q,b)&=\gcd(2e+d+1,b),\\
\gcd(Q',a)&=\gcd(2e+2d+1,a),
\end{aligned}
\]
so a tail of primes \(>a\) has all digits in
\[
D_{a,b}=\{0\le e<b:\gcd(2e+d+1,b)=\gcd(2e+2d+1,a)=1\}.
\]
\label{thm-end:62}
\end{proposition}
\end{samepage}

% Markdown AST block 505: Proof.
\begin{proof}\leavevmode\label{proof:62}
Put \(y=(Q+1)/2=z+h\); the largest odd integer below
\(2ry\) is \(2y'-1\) with \(y'=\lceil ry-1/2\rceil\).

% Markdown AST block 506: Para
Since \(b\) is
odd, \(ry\) is not a half-integer, so
\[
y'=\lfloor ry+1/2\rfloor
=\lfloor(2ay+b)/(2b)\rfloor=\lfloor(ay+(b-1)/2)/b\rfloor.
\]

% Markdown AST block 507: Para
With
\((b-1)/2=dh\) and \(a=b+d\),
\[
ay+dh=az+bh+(b-1),
\]
so
\[
\begin{aligned}
&y'=h+\lfloor(az+b-1)/b\rfloor=h+\lceil az/b\rceil\qquad\text{and}\\
&Q'=2y'-1=2\lceil az/b\rceil+H.
\end{aligned}
\]

% Markdown AST block 508: Para
For the gcd identities,
\[
dQ=2dz+b-1-d\equiv-2e-1-d\pmod b
\]
(as \(2e\equiv-2az\equiv-2dz\)),
and \(\gcd(d,b)=1\); similarly \(dQ'=2dz'+b-1-d\) with
\(bz'\equiv e\pmod a\) and \(b\equiv-d\pmod a\) give
\[
dQ'\equiv-(2e+2d+1)\pmod a.\qedhere
\]
\end{proof}

% Markdown AST block 509: Para
A uniform halting theorem for these families, showing that every
sufficiently large initial value eventually violates the digit
restriction or the interval condition, is not proved.

% Markdown AST block 510: Appendix D. Prime-prefix counting, the word sieve, and finite height
\Needspace{6\baselineskip}\section{Prime-prefix counting, the word sieve, and finite height}\label{app:D}

% Markdown AST block 511: Para
This appendix proves Theorem 47. Section D.1 bounds the number of
initial primes whose orbit stays prime for \(N\) steps; Section D.2
describes a sieve over carry words for a general base and the resulting
power-type bound; Sections D.3--D.4 turn the sieve into a complete
treatment of all initial integers up to an explicit height for \(6/5\)
and \(9/7\). None of these results is an all-coefficient theorem.

% Markdown AST block 512: D.1 A counting bound along the family
\Needspace{6\baselineskip}
\subsection{A counting bound along the family}\label{sec:D.1}

% Markdown AST block 513: Para
For \(t\ge1\), \(N\ge1\) and \(X\ge a=6t+3\), let \(S_t(X,N)\) be the
set of primes \(a<P\le X\) for which the program of Section C.3 accepts
\(N\) steps.

% Markdown AST block 514: Proposition 63.
\begin{samepage}
\begin{proposition}\leavevmode\label{thm:63}
\[
|S_t(X,N)|\le d_t^N(\lfloor X/(2b^N)\rfloor+1),
\]
where \(d_t=|D_t|\le4t+1\). If \(N=N(X)\) is the least integer with
\(b^N\ge X\), then
\[
|S_t(X,N(X))|\le d_t\,X^{\alpha_t}
\]
with
\[
\alpha_t=\log d_t/\log b<1.
\]
The same bounds hold for the set of
initial primes for which the program never halts.
\label{thm-end:63}
\end{proposition}
\end{samepage}

% Markdown AST block 515: Proof.
\begin{proof}\leavevmode\label{proof:63}
Each \(P\in S_t(X,N)\) determines a digit word of length
\(N\) with all digits in \(D_t\) (Lemma 56), hence one of at most
\(d_t^N\) words; by Proposition 53 the word determines
\(z_0\bmod b^N\), hence \(P=2z_0+H\) modulo \(2b^N\); and a residue class
modulo \(2b^N\) has at most \(\lfloor X/(2b^N)\rfloor+1\) elements in
\((a,X]\).

% Markdown AST block 516: Para
For \(N=N(X)\), \(2b^N>X\) gives at most one element per
word, and \(b^{N-1}<X\) gives
\[
d_t^N\le d_t\cdot d_t^{\log_bX}
=d_tX^{\alpha_t};
\]
\(\alpha_t<1\) since \(d_t<b\). A non-halting
initial prime lies in \(S_t(X,N)\) for every \(N\). \qedhere
\end{proof}

% Markdown AST block 517: Para
A bound of the form \(CX^\alpha\) with \(\alpha<1\) does not imply that
the set is empty: the tail of a single non-halting orbit has only
\(O(\log X)\) starting points below \(X\).

% Markdown AST block 518: D.2 The word sieve for a general base
\Needspace{6\baselineskip}
\subsection{The word sieve for a general base}\label{sec:D.2}

% Markdown AST block 519: Para
\begin{samepage}\leavevmode\label{layout:carry-intro-start}
Let \(a>b\ge2\) be coprime and let \(c_0\cdots c_{\ell-1}\) be a carry
word, with \(C_0=0\), \(C_{j+1}=aC_j+b^jc_j\), so that
\[
b^jQ_j=a^jQ_0+C_j
\]
by (8).
\label{layout:carry-intro-end}\end{samepage}

% Markdown AST block 520: Lemma 64.
\begin{samepage}
\begin{lemma}\leavevmode\label{thm:64}
(i) \emph{Integrality.} The congruence
\[
a^\ell Q_0\equiv-C_\ell\pmod{b^\ell}
\]
determines \(Q_0\bmod b^\ell\)
and implies
\[
b^j\mid a^jQ_0+C_j
\]
for all \(j\le\ell\); hence for
\(X\le b^\ell\) at most one integer \(Q_0\in[1,X]\) has the given word.
(ii) \emph{Exact interval.} With
\[
\theta_j=(a^j\theta_0-C_j)/b^j,
\]
the
word is realized by some real orbit segment if and only if
\[
[0,1)\cap\bigcap_{j=1}^{\ell}\Bigl[\frac{C_j}{a^j},\frac{C_j+b^j}{a^j}\Bigr)\ne\varnothing .
\]
(iii) \emph{Fixed divisors.} For a prime \(p\nmid ab\), \(p\mid Q_j\) if and
only if \(Q_0\equiv R_j\pmod p\), where \(R_0=0\) and
\[
R_{j+1}=R_j-c_jb^ja^{-j-1}\pmod p.
\]
If
\[
\{R_0,\ldots,R_\ell\}=\mathbb F_p,
\]
then every initial integer has
\(p\mid Q_j\) for some \(j\le\ell\); since there are at most \(\ell+1\)
roots, only primes \(p\le\ell+1\) can cover \(\mathbb F_p\).
\label{thm-end:64}
\end{lemma}
\end{samepage}

% Markdown AST block 521: Proof.
\begin{proof}\leavevmode\label{proof:64}
(i)
\[
C_\ell=a^{\ell-j}C_j+\sum_{i\ge j}a^{\ell-1-i}b^ic_i
\equiv a^{\ell-j}C_j\pmod{b^j},
\]
so \(b^\ell\mid a^\ell Q_0+C_\ell\)
gives
\[
b^j\mid a^{\ell-j}(a^jQ_0+C_j)
\]
and, as \(\gcd(a,b)=1\),
\[
b^j\mid a^jQ_0+C_j.
\]

% Markdown AST block 522: Para
\textbf{(ii)} \(\theta_j\in[0,1)\) if and only if
\[
\theta_0\in[C_j/a^j,(C_j+b^j)/a^j);
\]
the map
\(\theta_0\mapsto\theta_j\) is exact, as in Proposition 57.

% Markdown AST block 523: Para
\textbf{(iii)}
\[
p\mid Q_j\iff a^jQ_0+C_j\equiv0\iff Q_0\equiv-C_ja^{-j},
\]
\[
\text{and}\quad
-(aC_j+b^jc_j)a^{-j-1}=R_j-c_jb^ja^{-j-1}.\qedhere
\]
\end{proof}

% Markdown AST block 524: Para
Let \(A_\ell\) be the number of words of length \(\ell\) such that (1)
every letter lies in \(\mathcal C(a,b)\), (2) the interval in (ii) is
nonempty, and (3) no prime \(p\le\ell+1\) with \(p\nmid2ab\) has
\(\{R_0,\ldots,R_\ell\}=\mathbb F_p\). If \(Q_0,\ldots,Q_\ell\) are
primes larger than \(\max(a,\ell+1)\), then the carry word of the
segment satisfies (1) by Lemma 29, (2) with the true \(\theta_0\), and
(3) because a covering prime would divide some \(Q_j\), forcing
\(Q_j=p\le\ell+1\). Absence of a fixed divisor is a necessary condition
only; no prime-tuple conjecture is used to make it sufficient.

% Markdown AST block 525: Proposition 65 (power-type bound).
\begin{samepage}
\begin{proposition}[power-type bound]\leavevmode\label{thm:65}
Fix \(\ell\), put
\(\lambda=A_\ell^{1/\ell}\) and \(\alpha=\log\lambda/\log b\), and
suppose
\[
A_s\le C\lambda^{s-1}
\]
for \(0\le s<\ell\). Then the number of
primes \(P\) with \(\max(a,\ell+1)<P\le X\) such that some
\(\xi\in[P,P+1)\) produces \(N(X)+1\) prime terms is at most
\(CX^\alpha\), where \(N(X)\) is the least integer with \(b^{N(X)}\ge X\).
\label{thm-end:65}
\end{proposition}
\end{samepage}

% Markdown AST block 526: Proof.
\begin{proof}\leavevmode\label{proof:65}
Write \(N=k\ell+s\) with \(0\le s<\ell\). The carry word of
an all-prime segment of \(N+1\) terms, all larger than
\(\max(a,\ell+1)\), splits into \(k\) blocks of length \(\ell\) and one
of length \(s\), each satisfying (1)--(3) with the fractional part at the
start of the block as initial value; ignoring compatibility between
blocks overcounts, so there are at most
\[
A_\ell^kA_s\le
C\lambda^{N-1}
\]
words.

% Markdown AST block 527: Para
For \(N=N(X)\), \(b^{N-1}<X\) gives
\[
\lambda^{N-1}=b^{\alpha(N-1)}<X^\alpha,
\]
and by Lemma 64(i) each word
has at most one initial integer in \([1,X]\), distinct \(P\) giving
distinct words. \qedhere
\end{proof}

% Markdown AST block 528: Para
The hypothesis \(A_s\le C\lambda^{s-1}\) is equivalent to the integer
inequalities \(A_\ell\le C^\ell\) (\(s=0\)) and
\(A_s^{\ell}\le C^{\ell}A_\ell^{s-1}\) (\(1\le s<\ell\)). Table 11
gives the complete counts for the two bases, computed as described in
Section D.5; the exponents are certified by
\[
18035747<7^9=40353607\qquad \text{and}\qquad 3955763<5^{10}=9765625,
\]
that is,
\(\alpha<1/2\) for \(9/7\) and \(\alpha<1/4\) for \(6/5\), and the
constants \(C=13\) and \(C=10\) are the least integers satisfying the
inequalities. For \(6/5\) the alphabet is \(\mathcal C(6,5)=\{-1,1\}\),
and a prime successor requires \(Q\equiv1\pmod5\) (then \(c=-1\)) or
\(Q\equiv4\pmod5\) (then \(c=1\)), since \(c\equiv-Q\pmod5\); for
\(9/7\) it is \(\mathcal C(9,7)=\{-4,-2,2,4,8\}\).

% Markdown AST block 529: Table 11.
% Caption is rendered at the start of the immediately following table.

% Markdown AST block 530: Table 11
\begingroup
\small
\setlength{\tabcolsep}{4pt}
\renewcommand{\arraystretch}{1.15}
\setlength{\LTcapwidth}{\textwidth}
\begin{longtable}{@{}>{\raggedright\arraybackslash}p{\dimexpr0.08\dimexpr\linewidth-32pt\relax} >{\raggedright\arraybackslash}p{\dimexpr0.06\dimexpr\linewidth-32pt\relax} >{\raggedright\arraybackslash}p{\dimexpr0.65\dimexpr\linewidth-32pt\relax} >{\raggedright\arraybackslash}p{\dimexpr0.06\dimexpr\linewidth-32pt\relax} >{\raggedright\arraybackslash}p{\dimexpr0.15\dimexpr\linewidth-32pt\relax}@{}}
\caption{The sieve counts \(A_s\) for \(0\le s\le\ell\).}\label{tab:11}\\
\toprule
Base & \(\ell\) & \(A_0,\ldots,A_\ell\) & \(C\) & Range \\
\midrule
\endfirsthead
\toprule
Base & \(\ell\) & \(A_0,\ldots,A_\ell\) & \(C\) & Range \\
\midrule
\endhead
\bottomrule
\endlastfoot
\(9/7\) & 18 & 1,\allowbreak{} 5,\allowbreak{} 19,\allowbreak{} 67,\allowbreak{} 199,\allowbreak{} 530,\allowbreak{} 1315,\allowbreak{} 3130,\allowbreak{} 7250,\allowbreak{} 16511,\allowbreak{} 37155,\allowbreak{} 82942,\allowbreak{} 183838,\allowbreak{} 404586,\allowbreak{} 883163,\allowbreak{} 1909362,\allowbreak{} 4083475,\allowbreak{} 8631875,\allowbreak{} 18035747 & 13 & \(P>19\) \\*
\(6/5\) & 40 & 1,\allowbreak{} 2,\allowbreak{} 4,\allowbreak{} 8,\allowbreak{} 15,\allowbreak{} 28,\allowbreak{} 50,\allowbreak{} 79,\allowbreak{} 124,\allowbreak{} 191,\allowbreak{} 285,\allowbreak{} 417,\allowbreak{} 601,\allowbreak{} 864,\allowbreak{} 1233,\allowbreak{} 1757,\allowbreak{} 2493,\allowbreak{} 3522,\allowbreak{} 4950,\allowbreak{} 6965,\allowbreak{} 9717,\allowbreak{} 13582,\allowbreak{} 18786,\allowbreak{} 26085,\allowbreak{} 35937,\allowbreak{} 49661,\allowbreak{} 67971,\allowbreak{} 93166,\allowbreak{} 126734,\allowbreak{} 172576,\allowbreak{} 233228,\allowbreak{} 315696,\allowbreak{} 423604,\allowbreak{} 568869,\allowbreak{} 757642,\allowbreak{} 1010562,\allowbreak{} 1335866,\allowbreak{} 1766366,\allowbreak{} 2315828,\allowbreak{} 3038864,\allowbreak{} 3955763 & 10 & \(P>41\) \\
\end{longtable}
\endgroup

% Markdown AST block 531: D.3 Recovering the initial integer
\Needspace{6\baselineskip}
\subsection{Recovering the initial integer}\label{sec:D.3}

% Markdown AST block 532: Lemma 66.
\begin{samepage}
\begin{lemma}\leavevmode\label{thm:66}
Let \(N\ge2\) and let \(c\) be a word of length \(N\).
(a) \emph{(Seed.)} Put
\[
s(c)=(-C_N(a^N)^{-1})\bmod b^N\in[0,b^N).
\]
Every
prime \(P\le b^N\) whose orbit has carry word \(c\) equals \(s(c)\); in
particular, if \(s(c)\) is even or \(s(c)\le\max(a,N+1)\), no such
prime exceeds \(\max(a,N+1)\). (b) \emph{(Interval.)} The word is realized by
a real orbit segment with \(\lfloor\xi\rfloor=s(c)\) if and only if the
interval of Lemma 64(ii) is nonempty. (c) \emph{(Roots.)} If a prime
\(p\le N+1\), \(p\nmid2ab\), covers \(\mathbb F_p\), or if
\[
s(c)\bmod p\in\{R_0,\ldots,R_N\},
\]
then every orbit with initial
integer \(P=s(c)>N+1\) and carry word \(c\) has a term \(Q_j\) with
\(p\mid Q_j\), \(Q_j\ge P>p\), hence composite.
\label{thm-end:66}
\end{lemma}
\end{samepage}

% Markdown AST block 533: Proof.
\begin{proof}\leavevmode\label{proof:66}
(a) By Lemma 64(i), \(P\equiv s(c)\pmod{b^N}\); since
\(0<P\le b^N\) and \(b^N\) is composite, \(P=s(c)\) (and \(s(c)=0\)
would give \(P=b^N\), not prime).

% Markdown AST block 534: Para
\textbf{(b)} is Lemma 64(ii).

% Markdown AST block 535: Para
\textbf{(c)} is Lemma
64(iii) together with \(Q_j\ge Q_0=P\). \qedhere
\end{proof}

% Markdown AST block 536: D.4 Proof of Theorem 47
\Needspace{6\baselineskip}
\subsection{Proof of Theorem 47}\label{sec:D.4}

% Markdown AST block 537: Proof of Theorem 47.
\begin{proof}[Proof of Theorem 47]\leavevmode\label{proof:47}
(a) Let \(6<P=\lfloor\xi\rfloor\le5^{40}\) and
suppose that none of \(Q_0,\ldots,Q_{40}\) is composite; since all are
\(\ge P\ge7\), all are prime. Put \(N=40\). If \(P\le41=\max(6,N+1)\),
then \(P\in\{7,11,\ldots,41\}\), and each of these ten primes is
excluded individually (Table 13).

% Markdown AST block 538: Para
If \(P\ge43\), then all terms are
primes larger than \(\max(a,N+1)\), so the carry word \(c\) of length
\(40\) satisfies (1)--(3) of Section D.2 with \(a/b=6/5\), and by Lemma
66(a) \(P=s(c)\), which is odd and larger than \(41\), and by Lemma
66(c)
\[
s(c)\bmod p\notin\{R_j\}
\]
for every prime
\(p\in\{7,11,\ldots,41\}\).

% Markdown AST block 539: Para
The complete enumeration of all
\(A_{40}=3955763\) words satisfying (1)--(3) (Table 12) finds
\(1978668\) with even seed and \(1977095\) rejected by the root test,
none with seed \(\le41\), and no survivor. This contradiction proves
(a).

% Markdown AST block 540: OrderedList
\begin{enumerate}
\def\labelenumi{(\alph{enumi})}
\setcounter{enumi}{1}
\item
  Let \(9<P\le7^{18}\) and suppose none of \(Q_0,\ldots,Q_{18}\) is
  composite. With \(N=18\), the primes \(P\le19=\max(9,N+1)\), namely
  \(11,13,17,19\), are excluded individually (Table 13).

  For \(P\ge23\)
  the same sieve with \(a/b=9/7\), primes \(p\in\{5,11,13,17,19\}\), and
  \(A_{18}=18035747\) words finds \(9018851\) even seeds and \(9012304\)
  root rejections, none with seed \(\le19\), and \(4592\) surviving
  candidates \(P\).

  For each survivor, all terms are primes, hence coprime
  to \(6\), so by Proposition 55 the orbit follows
  \[
  T(Q)=Q+2\lfloor(Q+4)/7\rfloor
  \]
  (the map \(T_1\)), and the carries of
  this \(T\)-orbit are those of the sieve word: the carry is determined by
  \[
  c\equiv-2Q\pmod7,
  \]
  and the letters \(-4,-2,2,4,8\) have distinct
  residues \(3,5,2,4,1\) modulo \(7\).

  Each of the \(4592\) candidates
  has a proper divisor at some time \(j\le7\): there is an integer
  \(1<d<T^j(P)\) with \(d\mid T^j(P)\), and the recorded witnesses use
  divisors at most \(9923\), each checked as an exact equality
  \[
  T^j(P)=d\cdot e.
  \]
  Hence \(Q_j\) is composite for some \(j\le7\le18\),
  a contradiction. \qedhere
\end{enumerate}
\end{proof}

% Markdown AST block 541: Table 12.
% Caption is rendered at the start of the immediately following table.

% Markdown AST block 542: Table 12
\begingroup
\small
\setlength{\tabcolsep}{2.5pt}
\renewcommand{\arraystretch}{1.15}
\setlength{\LTcapwidth}{\textwidth}
\begin{longtable}{@{} r r r r r r r r@{}}
\caption{The finite-height sieve. ``Words'' is \(A_N\); ``even'' and
``roots'' count words whose seed is even or is rejected by Lemma 66(c);
``small'' counts seeds \(\le\max(a,N+1)\).}\label{tab:12}\\
\toprule
\shortstack[r]{Base} & \shortstack[r]{\(N\)} & \shortstack[r]{Height\\\(b^N\)} & \shortstack[r]{Words} & \shortstack[r]{Even\\seed} & \shortstack[r]{Root\\rejection} & \shortstack[r]{Small\\seed} & \shortstack[r]{Survivors} \\
\midrule
\endfirsthead
\toprule
\shortstack[r]{Base} & \shortstack[r]{\(N\)} & \shortstack[r]{Height\\\(b^N\)} & \shortstack[r]{Words} & \shortstack[r]{Even\\seed} & \shortstack[r]{Root\\rejection} & \shortstack[r]{Small\\seed} & \shortstack[r]{Survivors} \\
\midrule
\endhead
\bottomrule
\endlastfoot
\(6/5\) & 40 & \(5^{40}\) & 3955763 & 1978668 & 1977095 & 0 & 0 \\*
\(9/7\) & 18 & \(7^{18}\) & 18035747 & 9018851 & 9012304 & 0 & 4592 \\
\end{longtable}
\endgroup

% Markdown AST block 543: Para
For the \(4592\) survivors the maximal time at which a divisor is used
is \(7\). The distribution of these times over the candidates depends on
the order and range of the trial divisions and is not reported; the
statement used in the proof is the existence of a proper divisor at some
time \(j\le7\) for every candidate.

% Markdown AST block 544: Table 13.
% Caption is rendered at the start of the immediately following table.

% Markdown AST block 545: Table 13
\begingroup
\small
\setlength{\tabcolsep}{4pt}
\renewcommand{\arraystretch}{1.15}
\setlength{\LTcapwidth}{\textwidth}
\begin{longtable}{@{}>{\raggedright\arraybackslash}p{\dimexpr0.1\dimexpr\linewidth-24pt\relax} >{\raggedright\arraybackslash}p{\dimexpr0.08\dimexpr\linewidth-24pt\relax} >{\raggedright\arraybackslash}p{\dimexpr0.67\dimexpr\linewidth-24pt\relax} >{\raggedright\arraybackslash}p{\dimexpr0.15\dimexpr\linewidth-24pt\relax}@{}}
\caption{Small initial primes. For \(x\in[P,P+1)\) the successor
\(\lfloor rx\rfloor\) has at most two candidate values for \(6/5\)
and at most three for \(9/7\) (Section C.5); the rows list every
candidate where it matters. For \(6/5\),
``\(Q\equiv s\pmod5\)'' with \(s\notin\{1,4\}\) means that no candidate
successor of \(Q\) is coprime to \(30\) (Section D.2), so every
candidate, being at least \(8\), is composite; where a unique candidate
coprime to \(30\) exists, it is named. ``Empty interval'' means that the
exact interval of Section C.3 for the all-prime \(T\)-orbit becomes empty
at the indicated step, so that orbit is realized by no \(\xi\), while
any deviation from it is a term divisible by \(2\) or \(3\), hence
composite.}\label{tab:13}\\
\toprule
Base & \(P\) & Witness & Composite by time \\
\midrule
\endfirsthead
\toprule
Base & \(P\) & Witness & Composite by time \\
\midrule
\endhead
\bottomrule
\endlastfoot
\(6/5\) & 7 & \(7\equiv2\pmod5\) & 1 \\
\(6/5\) & 11 & the only candidate coprime to \(30\) is \(13\); \(13\equiv3\pmod5\) & 2 \\
\(6/5\) & 13 & \(13\equiv3\pmod5\) & 1 \\
\(6/5\) & 17 & \(17\equiv2\pmod5\) & 1 \\
\(6/5\) & 19 & the only candidate coprime to \(30\) is \(23\); \(23\equiv3\pmod5\) & 2 \\
\(6/5\) & 23 & \(23\equiv3\pmod5\) & 1 \\
\(6/5\) & 29 & both candidates \(34\) and \(35=5\cdot7\) are composite & 1 \\
\(6/5\) & 31 & the only candidate coprime to \(30\) is \(37\); \(37\equiv2\pmod5\) & 2 \\
\(6/5\) & 37 & \(37\equiv2\pmod5\) & 1 \\
\(6/5\) & 41 & both candidates \(49=7\cdot7\) and \(50\) are composite & 1 \\
\(9/7\) & 11 & both candidates \(14\) and \(15=3\cdot5\) are composite & 1 \\
\(9/7\) & 13 & empty interval after step 4 (\(L=2450>U=2401\),\allowbreak{} \(B=7^4\)) & 4 \\
\(9/7\) & 17 & empty interval after step 3 (\(L=350>U=343\),\allowbreak{} \(B=7^3\)) & 3 \\
\(9/7\) & 19 & both candidates \(24\) and \(25=5\cdot5\) are composite & 1 \\
\end{longtable}
\endgroup

% Markdown AST block 546: Para
Theorem 47 says nothing about \(\lfloor\xi\rfloor>5^{40}\) or
\(\lfloor\xi\rfloor>7^{18}\). For an initial integer \(s+kb^N\) with the
same first \(N\) carries as \(s\), the terms are
\[
Q_j(s)+ka^jb^{N-j},
\]
and a divisor found for \(k=0\) need not divide
the terms for \(k\ne0\).

% Markdown AST block 547: D.5 Verification of the finite-height computation
\Needspace{6\baselineskip}
\subsection{Verification of the finite-height computation}\label{sec:D.5}

% Markdown AST block 548: Para
The counts of Tables 11--13 and the list of \(4592\) survivors were
produced by implementation A of Section 9.3, that is, by the two
same-lineage constructions supplied with the research notes, which track
the exact interval in different coordinates (the current fractional
interval with the seed lifted digit by digit, and the initial interval
with common denominator \(a^N\) with the seed solved at the leaves) and
agree on all depth counts, all classification counts, and the survivor
list; on six small instances
\[
(a,b,N)\in\{(3,2,5),(4,3,5),(5,3,5),
(6,5,6),(7,5,5),(9,7,5)\}
\]
with \(9,11,37,50,54,530\) words, they also
agree with a Python oracle enumerating all words by direct product and
exact rational arithmetic. The sieve counts of Table 11 were recomputed
in full. The finite-height computations and their reproduction are described in
the \href{supplement/COMPUTATION.md}{computation guide}.

% Markdown AST block 549: Para
Independently of these, implementation B of Section 9.3 (a different
programmer, code written from the specification of Lemmas 64--66 and not
from the supplied programs) re-enumerated the sieve for both bases,
recomputed the seeds, the root rejections and the survivor list, checked
every survivor's divisor equality and every small-prime witness, and
recomputed the constants of Table 11 as least integers from its own
counts. All quantities compared element by element (the count arrays,
the classification counts, the survivor set, the \(4592\) witness
equalities, the fourteen small-prime witnesses, the six small
instances, and the constants) agree. The comparison is described in
the \href{supplement/COMPUTATION.md}{computation guide}, with reference data in \href{checks/reference.json}{checks/reference.json}.
Agreement is evidence that both implementations computed the same finite
objects; the passage from these objects to Theorem 47 is the argument of
Sections D.2--D.4.

% Markdown AST block 550: Appendix E. Long prime prefixes and return-word obstructions
\Needspace{6\baselineskip}\section{Long prime prefixes and return-word obstructions}\label{app:E}

% Markdown AST block 551: E.1 Arbitrarily long prime prefixes along the family
\subsection{Arbitrarily long prime prefixes along the family}\label{sec:E.1}

% Markdown AST block 552: Lemma 67 (transplanting a prime progression).
\begin{samepage}
\begin{lemma}[transplanting a prime progression]\leavevmode\label{thm:67}
Let \(L\ge3\), let
\(P,P+D,\ldots,P+(L-1)D\) be primes larger than \(3\), and let
\[
P>48L^2D^2.
\]
Let \(b\) be an integer with \(b\equiv1\pmod6\) and
\(|b-2P/D|\le3\), and put \(a=b+2\), \(t=(b-1)/6\), \(\xi=P+1/2\). Then
\(t\ge1\), \(P>a\), and
\[
\lfloor\xi(a/b)^j\rfloor=P+jD
\]
for
\(0\le j<L\).
\label{thm-end:67}
\end{lemma}
\end{samepage}

% Markdown AST block 553: Proof.
\begin{proof}\leavevmode\label{proof:67}
\emph{\(6\mid D\).} If \(D\) were odd, one of \(P,P+D\) would be
even; if \(3\nmid D\), the three terms \(P,P+D,P+2D\) would cover all
residues modulo \(3\).

% Markdown AST block 554: Para
\emph{Existence of \(b\).} Every real number is
within \(3\) of an integer \(\equiv1\pmod6\). Since
\[
2P/D>96L^2D\ge96\cdot9\cdot6,
\]
we get
\[
b\ge2P/D-3\ge P/D\ge7,
\]
so
\(t\ge1\).

% Markdown AST block 555: Para
\emph{First estimate.} Write \(b=2P/D+\varepsilon\) with
\(|\varepsilon|\le3\); then
\[
(2P+1)/b-D=(1-D\varepsilon)/b
\]
and \(|1-D\varepsilon|\le4D\), so
\[
|(2P+1)/b-D|\le4D^2/P.
\]

% Markdown AST block 556: Para
\emph{Binomial expansion.} With \(u=2/b\),
\[
(1+u)^j=1+ju+R_j
\]
where
\[
0\le R_j\le\sum_{k\ge2}(ju)^k/k!\le\tfrac12(ju)^2e^{ju};
\]
here
\[
ju\le2L/b\le2LD/P<1,
\]
so \(e^{ju}<3\) and
\[
R_j\le\tfrac32(ju)^2=6j^2/b^2\le6j^2D^2/P^2.
\]
Therefore
\[
\begin{aligned}
&\Bigl|\bigl(P+\tfrac12\bigr)(1+u)^j-\bigl(P+jD+\tfrac12\bigr)\Bigr|\\
&\le j\Bigl|\frac{2P+1}{b}-D\Bigr|+\bigl(P+\tfrac12\bigr)R_j\\
&\le\frac{4jD^2}{P}+\frac{8j^2D^2}{P}\\
&\le\frac{12L^2D^2}{P}<\frac14,
\end{aligned}
\]
using \((P+1/2)/P\le4/3\) and \(4j+8j^2\le12L^2\) for \(j<L\). Hence
\[
\xi(a/b)^j\in(P+jD+1/4,P+jD+3/4)
\]
and the integer part is \(P+jD\).

% Markdown AST block 557: Para
\emph{\(P>a\).}
\[
a=b+2\le2P/D+5\le P/3+5<P
\]
since \(D\ge6\). \qedhere
\end{proof}

% Markdown AST block 558: Theorem 68 (arbitrarily long prime prefixes).
\begin{samepage}
\begin{theorem}[arbitrarily long prime prefixes]\leavevmode\label{thm:68}
For every \(L\ge3\)
and every \(T\) there are \(t\ge T\) and a rational \(\xi>0\) such that
\(\lfloor\xi((6t+3)/(6t+1))^j\rfloor\) is prime for \(0\le j<L\) and
the first term exceeds \(6t+3\). Consequently there is no bound, uniform
in \(t\) and in the initial prime, on the number of steps after which
the program of Section C.3 halts.
\label{thm-end:68}
\end{theorem}
\end{samepage}

% Markdown AST block 559: Proof.
\begin{proof}\leavevmode\label{proof:68}
This uses as an external input the theorem on narrow
arithmetic progressions in the primes: for \(k\ge2\) and \(\delta>0\),
every subset \(A\) of the primes in \([1,Y]\) with
\[
|A|\ge\delta Y/\log Y
\]
contains, for all sufficiently large \(Y\), a
nontrivial \(k\)-term arithmetic progression whose common difference
\(D\) satisfies
\[
|D|\le C_{k,\delta}(\log Y)^{(k-1)2^{k-2}}
\]
\hyperref[ref-shao15]{SHAO15, Theorem 1.3}; the qualitative form, with an
exponent depending on \(k\) and \(\delta\) only, is
\hyperref[ref-tz14]{TZ14, Theorem 5} specialized to arithmetic progressions.
(In {[}TZ14{]} and {[}SHAO15{]} the letter \(L\) denotes the exponent; here
\(L\) is the length \(k\).)

% Markdown AST block 560: Para
Take \(k=L\) and
\[
A=\{p\text{ prime}:Y/2<p\le Y\};
\]
by the prime number theorem
\[
|A|\ge\tfrac13Y/\log Y
\]
for large \(Y\), so \(\delta=1/3\) is
admissible.

% Markdown AST block 561: Para
The progression obtained lies in \((Y/2,Y]\), so its least
term \(P\) satisfies \(P>Y/2\), all terms exceed \(3\) for large \(Y\),
and after reversing the order if necessary \(D>0\) with
\[
D\le C(\log Y)^{(L-1)2^{L-2}}.
\]

% Markdown AST block 562: Para
Hence
\[
P/(48L^2D^2)\to\infty
\]
as
\(Y\to\infty\), and Lemma 67 applies for large \(Y\), giving
\[
t=(b-1)/6\ge(P/D-1)/6\to\infty.\qedhere
\]
\end{proof}

% Markdown AST block 563: Para
The statement is
\[
\forall L\,\forall T\,\exists t\ge T\,\exists\xi;
\]
it does not assert an infinite tail of primes for any base, and it does
not contradict any halting statement for a fixed \(t\). The implied
constants and the threshold for \(Y\) in the external theorem are not
explicit, so no explicit \(t\) is obtained this way. A narrow-progression
theorem is essential: a progression with \(D\) comparable to \(P\) would
not satisfy the hypothesis of Lemma 67.

% Markdown AST block 564: Para
\emph{Example (nine prime terms).} With
\[
\begin{aligned}
&t=1654617,\qquad a=9927705,\qquad b=9927703,\\
&\xi=2084817839/2,\qquad P=1042408919,\qquad D=210,
\end{aligned}
\]
one
has
\[
\lfloor\xi(a/b)^j\rfloor=P+210j
\]
for \(0\le j\le8\), and these
nine integers are prime, while
\[
\lfloor\xi(a/b)^9\rfloor=1042410809=11\cdot94764619.
\]
The hypothesis
\[
P>48\cdot81\cdot210^2=171460800
\]
holds, and \(b\) is the integer
\(\equiv1\pmod6\) nearest to \(2P/D\). The integer parts were computed
in exact rational arithmetic and the primes checked by trial division.
This is not a counterexample to anything.

% Markdown AST block 565: E.2 Return words as an obstruction
\Needspace{6\baselineskip}
\subsection{Return words as an obstruction}\label{sec:E.2}

% Markdown AST block 566: Para
The following lemma is of the same type as Theorem 42: it constructs a
right-directed pseudo-orbit with non-eventually-periodic carries and
applies Lemma 38.

% Markdown AST block 567: Lemma 69 (return-word obstruction).
\begin{samepage}
\begin{lemma}[return-word obstruction]\leavevmode\label{thm:69}
Let \(\gcd(m,ab)=1\),
\(q\in(\mathbb Z/m)^\times\), and let \(U,V\) be carry words such that\\
\textbf{(1)} following
\[
bz\equiv aq'+c\pmod m
\]
from \(q\) along each word, all residues are units and the last residue is \(q\);\\
\textbf{(2)} there is a closed interval \(I\subset(0,1)\) such that the backward maps
\[
\theta\mapsto(b\theta+c)/a
\]
applied along each word from a point of \(I\) keep all values in \((0,1)\)
and map \(I\) into \(I\);\\
\textbf{(3)} \(UV\ne VU\);\\
\textbf{(4)} every letter satisfies
\[
\gcd(c,ab)=1\qquad\text{and}\qquad c\equiv b-a\pmod2.
\]
Then for every \(M\) whose prime factors divide \(2abm\) there is a
right-directed \(M\)-pseudo-orbit with non-eventually-periodic carries,
and no execution of \(\mathfrak M_0\) (or of an extension by
pseudo-orbit-sound operations) with all \(M_t\mid M\) succeeds.
\label{thm-end:69}
\end{lemma}
\end{samepage}

% Markdown AST block 568: Proof.
\begin{proof}\leavevmode\label{proof:69}
\emph{(i) Prime powers of \(m\).} Let \(\tilde m\) be the least
common multiple of \(m\) and the part of \(M\) composed of primes of
\(m\). For a word \(w\) of length \(\ell\), the residue map
\[
\pi_w:x\mapsto(a/b)^\ell x+C(w)/b^\ell
\]
modulo \(\tilde m\) is an
affine permutation (\(p\nmid ab\)). Since \(\pi_U(q)\equiv q\pmod m\),
\(\pi_U\) permutes the finite fibre
\[
F=\{x\in\mathbb Z/\tilde m:x\equiv q\pmod m\},
\]
and likewise
\(\pi_V\); so for any \(\hat x\in F\) there are \(s,t\ge1\) with
\[
\pi_U^s(\hat x)=\hat x\qquad \text{and}\qquad \pi_V^t(\hat x)=\hat x.
\]
Intermediate
residues reduce modulo \(m\) to the units of (1), hence are units
modulo \(\tilde m\).

% Markdown AST block 569: Para
\emph{(ii) Non-commutation.} If \(U^sV^t=V^tU^s\), then
\(U^s\) and \(V^t\) are powers of a common primitive word \(z\); since
the primitive root of \(w^k\) is that of \(w\), \(U\) and \(V\) are
both powers of \(z\) and commute, contradicting (3).

% Markdown AST block 570: Para
\emph{(iii) Fractional
parts.} \(U,V\) map \(I\) into \(I\), hence so do \(U^s,V^t\), with all
intermediate values in \((0,1)\). Concatenate the blocks
\(X_0=U^sV^t\) and \(X_1=V^tU^s\) (of equal length and distinct)
according to a binary sequence that is not eventually periodic; the
nested closed intervals of the proof of Theorem 42 (each backward block
map is a contraction with slope \((b/a)^{|X_i|}\)) define \(\theta_n\)
for all \(n\ge0\) with
\[
a\theta_n-b\theta_{n+1}=c_n
\]
exactly, and
block decoding shows that the carries are not eventually periodic.

% Markdown AST block 571: Para
\emph{(iv) Primes dividing \(2ab\).} For \(p\mid a\), \(p\mid b\), and
\(p=2\) with \(a,b\) odd, construct unit residues modulo \(p^e\) as in
the proof of Theorem 42, using (4). Combine all residues by the Chinese
remainder theorem to a pseudo-orbit modulo \(\tilde m\cdot(\text{the
rest of }M)\) and reduce to \(M\). Lemma 38 gives the conclusion. \qedhere
\end{proof}

% Markdown AST block 572: Corollary 70 (explicit witnesses).
\begin{samepage}
\begin{corollary}[explicit witnesses]\leavevmode\label{thm:70}
(a) For \(9/7\), with
\[
m=715=5\cdot11\cdot13,\qquad q=713,\qquad I=[1/5,3/10],
\]
and the words
\(U,V\) of Table 14, no execution of \(\mathfrak M_0\) whose moduli
involve only the primes \(S_{13}=\{2,3,5,7,11,13\}\) succeeds; in
particular none with modulus dividing a power of \(4290\). (b) For
\(6/5\), with
\[
m=17017=7\cdot11\cdot13\cdot17,\qquad q=12854,\qquad I=[9/10,49/50],
\]
and the words \(U,V\) of Table 15, no execution
whose moduli involve only the primes
\(S_{17}=\{2,3,5,7,11,13,17\}\) succeeds.
\label{thm-end:70}
\end{corollary}
\end{samepage}

% Markdown AST block 573: Proof.
\begin{proof}\leavevmode\label{proof:70}
The hypotheses of Lemma 69 are verified in exact arithmetic:
the residue sequences of Tables 14--15 return to \(q\) through units;
the backward images of the endpoints of \(I\) (sufficient, since the
maps are increasing and affine) stay in \((0,1)\) and land in \(I\);
\(UV\ne VU\) since the words have different lengths and neither is a
prefix of the other's period structure (for \(9/7\), \(U\) and \(V\)
differ at position \(2\); for \(6/5\), at position \(19\)); and all
letters lie in \(\mathcal C(9,7)=\{-4,-2,2,4,8\}\), respectively
\(\mathcal C(6,5)=\{-1,1\}\). \qedhere
\end{proof}

% Markdown AST block 574: Table 14.
% Caption is rendered at the start of the immediately following table.

% Markdown AST block 575: Table 14
\begingroup
\small
\setlength{\tabcolsep}{4pt}
\renewcommand{\arraystretch}{1.15}
\setlength{\LTcapwidth}{\textwidth}
\begin{longtable}{@{}>{\raggedright\arraybackslash}p{\dimexpr0.16\dimexpr\linewidth-16pt\relax} >{\raggedright\arraybackslash}p{\dimexpr0.38\dimexpr\linewidth-16pt\relax} >{\raggedright\arraybackslash}p{\dimexpr0.46\dimexpr\linewidth-16pt\relax}@{}}
\caption{The \(9/7\) witness (\(m=715\), \(q=713\)). Residues
are listed from \(q\) to the return.}\label{tab:14}\\
\toprule
Word & Letters & Residues \\
\midrule
\endfirsthead
\toprule
Word & Letters & Residues \\
\midrule
\endhead
\bottomrule
\endlastfoot
\(U\) (length 18) & \(2,\allowbreak{}-2,\allowbreak{}2,\allowbreak{}-2,\allowbreak{}4,\allowbreak{}-4,\allowbreak{}2,\allowbreak{}2,\allowbreak{}2,\allowbreak{}2,\allowbreak{}-4,\allowbreak{}4,\allowbreak{}-2,\allowbreak{}4,\allowbreak{}4,\allowbreak{}-4,\allowbreak{}2,\allowbreak{}4\) & 713,\allowbreak{} 202,\allowbreak{} 668,\allowbreak{} 42,\allowbreak{} 258,\allowbreak{} 128,\allowbreak{} 164,\allowbreak{} 109,\allowbreak{} 549,\allowbreak{} 604,\allowbreak{} 164,\allowbreak{} 6,\allowbreak{} 519,\allowbreak{} 667,\allowbreak{} 41,\allowbreak{} 564,\allowbreak{} 316,\allowbreak{} 713,\allowbreak{} 713 \\*
\(V\) (length 17) & \(2,\allowbreak{}-2,\allowbreak{}-4,\allowbreak{}8,\allowbreak{}-4,\allowbreak{}4,\allowbreak{}-2,\allowbreak{}4,\allowbreak{}-4,\allowbreak{}4,\allowbreak{}2,\allowbreak{}2,\allowbreak{}2,\allowbreak{}-4,\allowbreak{}2,\allowbreak{}4,\allowbreak{}4\) & 713,\allowbreak{} 202,\allowbreak{} 668,\allowbreak{} 654,\allowbreak{} 127,\allowbreak{} 367,\allowbreak{} 166,\allowbreak{} 111,\allowbreak{} 654,\allowbreak{} 636,\allowbreak{} 614,\allowbreak{} 279,\allowbreak{} 359,\allowbreak{} 564,\allowbreak{} 316,\allowbreak{} 713,\allowbreak{} 713,\allowbreak{} 713 \\
\end{longtable}
\endgroup

% Markdown AST block 576: Table 15.
% Caption is rendered at the start of the immediately following table.

% Markdown AST block 577: Table 15
\begingroup
\small
\setlength{\tabcolsep}{4pt}
\renewcommand{\arraystretch}{1.15}
\setlength{\LTcapwidth}{\textwidth}
\begin{longtable}{@{}>{\raggedright\arraybackslash}p{\dimexpr0.16\dimexpr\linewidth-16pt\relax} >{\raggedright\arraybackslash}p{\dimexpr0.27\dimexpr\linewidth-16pt\relax} >{\raggedright\arraybackslash}p{\dimexpr0.57\dimexpr\linewidth-16pt\relax}@{}}
\caption{The \(6/5\) witness (\(m=17017\), \(q=12854\)); \(+\)
and \(-\) denote the carries \(+1\) and \(-1\).}\label{tab:15}\\
\toprule
Word & Letters & Residues \\
\midrule
\endfirsthead
\toprule
Word & Letters & Residues \\
\midrule
\endhead
\bottomrule
\endlastfoot
\(U\) (length 33) & \(+^{15}-+^{3}-^{3}+^{9}-^{2}\) & 12854,\allowbreak{} 15425,\allowbreak{} 8300,\allowbreak{} 16767,\allowbreak{} 6507,\allowbreak{} 11212,\allowbreak{} 16858,\allowbreak{} 13423,\allowbreak{} 9301,\allowbreak{} 7758,\allowbreak{} 2503,\allowbreak{} 13214,\allowbreak{} 15857,\allowbreak{} 5415,\allowbreak{} 13305,\allowbreak{} 5756,\allowbreak{} 6907,\allowbreak{} 11692,\allowbreak{} 417,\allowbreak{} 3904,\allowbreak{} 8088,\allowbreak{} 6302,\allowbreak{} 14369,\allowbreak{} 226,\allowbreak{} 13885,\allowbreak{} 6452,\allowbreak{} 11146,\allowbreak{} 9972,\allowbreak{} 15370,\allowbreak{} 8234,\allowbreak{} 9881,\allowbreak{} 8454,\allowbreak{} 13548,\allowbreak{} 12854 \\*
\(V\) (length 39) & \(+^{15}-+^{12}-^{2}+^{3}-+^{3}-^{2}\) & 12854,\allowbreak{} 15425,\allowbreak{} 8300,\allowbreak{} 16767,\allowbreak{} 6507,\allowbreak{} 11212,\allowbreak{} 16858,\allowbreak{} 13423,\allowbreak{} 9301,\allowbreak{} 7758,\allowbreak{} 2503,\allowbreak{} 13214,\allowbreak{} 15857,\allowbreak{} 5415,\allowbreak{} 13305,\allowbreak{} 5756,\allowbreak{} 6907,\allowbreak{} 11692,\allowbreak{} 417,\allowbreak{} 3904,\allowbreak{} 4685,\allowbreak{} 12429,\allowbreak{} 14915,\allowbreak{} 7688,\allowbreak{} 2419,\allowbreak{} 2903,\allowbreak{} 13694,\allowbreak{} 16433,\allowbreak{} 12913,\allowbreak{} 12092,\allowbreak{} 4300,\allowbreak{} 11967,\allowbreak{} 747,\allowbreak{} 4300,\allowbreak{} 15370,\allowbreak{} 8234,\allowbreak{} 9881,\allowbreak{} 8454,\allowbreak{} 13548,\allowbreak{} 12854 \\
\end{longtable}
\endgroup

% Markdown AST block 578: Para
For \(6/5\) a shorter witness modulo \(1001\) (root \(380\),
\(I=[4/5,19/20]\), two words of length \(27\)) also exists; the
witness modulo \(17017\) suffices for (b). Examples of the lifting
orders in step (i): for \(9/7\) modulo \(715\cdot55\) one may take
\((s,t)=(10,20)\), and modulo \(715\cdot169\), \((26,52)\); for
\(6/5\) modulo \(17017\cdot77\), \((14,14)\), and modulo
\(17017\cdot289\), \((136,136)\).

% Markdown AST block 579: Para
As in Section 8, these are statements about the non-termination of the
procedure. The residues \(q_n\) form a sequence of residue classes, not
the integer parts of a single orbit. By Lemma 41 every finite window of
the witness is realized by a true orbit with arbitrarily large
coefficient, in the form \(\forall N\,\exists\xi_N\); this does not give
a single \(\xi\) realizing the whole witness, and the realizing integers
are not claimed to be prime, only coprime to the specified primes.
Unlike Theorem 58, no initial integer is fixed here, so the compactness
argument does not apply.

% Markdown AST block 580: E.3 Eliminating one return-word language with a further prime
\Needspace{6\baselineskip}
\subsection{Eliminating one return-word language with a further prime}\label{sec:E.3}

% Markdown AST block 581: Proposition 71.
\begin{samepage}
\begin{proposition}\leavevmode\label{thm:71}
Let \(a=9\), \(b=7\), and let \(U,V\) be the words
of Table 14. Regard \(U\) and \(V\) as single edges and form the graph
on the residues \(q\in\{1,\ldots,16\}\) modulo \(17\) with an edge
\(q\xrightarrow{w}z\) whenever following \(w\) from \(q\) modulo \(17\)
keeps all residues, including the last, nonzero. The graph has exactly
the twelve edges
\[
\begin{aligned}
&4\xrightarrow{U}16,\ 4\xrightarrow{V}16,\ 5\xrightarrow{U}1,\
\\
&6\xrightarrow{V}4,\ 8\xrightarrow{V}9,\ 9\xrightarrow{U}9,\
\\
&10\xrightarrow{V}14,\ 11\xrightarrow{U}13,\ 11\xrightarrow{V}8,\
\\
&14\xrightarrow{U}2,\ 15\xrightarrow{U}4,\ 16\xrightarrow{V}12,
\end{aligned}
\]
and with the ranks
\[
\operatorname{rank}(1,\ldots,16)=(0,0,0,2,1,3,0,1,0,2,2,0,0,1,3,1)
\]
every edge other than the self-loop \(9\xrightarrow{U}9\) strictly
decreases the rank. Consequently there is no true orbit with ratio
\(9/7\) all of whose terms are prime and whose carry sequence is an
infinite concatenation of \(U\) and \(V\).
\label{thm-end:71}
\end{proposition}
\end{samepage}

% Markdown AST block 582: Proof.
\begin{proof}\leavevmode\label{proof:71}
The edge list and the rank inequalities are finite
verifications. For an all-prime orbit, \(17\nmid Q_n\) once \(Q_n>17\),
so the residues \(Q_n\bmod17\) at block boundaries form an infinite walk
in the graph. Since the rank strictly decreases on every edge except the
self-loop, after finitely many blocks only \(9\xrightarrow{U}9\) is
used, and the carry sequence is eventually \(U^\infty\), periodic with
period \(18\), contradicting Lemma 4. \qedhere
\end{proof}

% Markdown AST block 583: Para
The edge table depends only on \(a\bmod17\) and \(b\bmod17\), so the
same twelve edges arise for every pair \(a\equiv9\), \(b\equiv7\pmod{17}\);
but \(U\) and \(V\) are return words only for \(9/7\) (for
\(a=111\), \(b=109\), the word \(U\) from \(713\) modulo \(715\) does
not return and meets a non-unit), so nothing is asserted for other
members of the family. Proposition 71 eliminates the language generated
by \(U\) and \(V\); it does not cover all carry sequences of \(9/7\),
and it does not imply that adding the prime \(17\) yields a certificate
for \(9/7\).

% Markdown AST block 584: E.4 Explicit two-carry certificates, prime selection, and the whole-word interval test
\Needspace{6\baselineskip}
\subsection{Explicit two-carry certificates, prime selection, and the whole-word interval test}\label{sec:E.4}

% Markdown AST block 585: Para
\emph{Instances of Theorem 42.} For \(a/b=101/2\) with \(R=15\) (auxiliary
primes \(3,5\)):
\[
\begin{aligned}
&\Delta=99,\qquad m=1,\qquad h=30,\\
&c=69,\qquad \gcd(69,202)=1;
\end{aligned}
\]
the least admissible \(L\) is \(0\), with
\[
\begin{aligned}
&x_0=23/33,\qquad x_1=3379/3399,\\
&I=[23/33,3379/3399],\quad\text{and}\quad U=(69),\quad V=(99,69).
\end{aligned}
\]
Hence no execution of \(\mathfrak M_0\) whose
moduli involve only the primes \(\{2,3,5,101\}\) succeeds.

% Markdown AST block 586: Para
For
\(a/b=101/100\) with \(R=21\) (primes \(3,7\)):
\[
\begin{aligned}
&\Delta=1,\qquad m=1,\qquad h=42,\\
&c=-41,\qquad \gcd(41,10100)=1;
\end{aligned}
\]
the threshold is
\(-c/b=41/100\), and the least \(L\) with \(x_L>41/100\), where
\[
x_L=1-42\cdot100^L/(101^{L+1}-100^{L+1}),
\]
is \(L=54\), so
\(I=[x_{54},x_{55}]\) with the words \(U=1^{54}(-41)\) and
\(V=1^{55}(-41)\). This example shows that negative carries are handled
exactly.

% Markdown AST block 587: Para
For \(5/2\) with \(R=1\):
\[
\begin{aligned}
&\Delta=3,\qquad m=1,\qquad h=2,\qquad c=1,\qquad L=0,\\
&U=(1),\qquad V=(3,1),
\end{aligned}
\]
which is the witness of
Corollary 43 for this base.

% Markdown AST block 588: Lemma 72 (primes producing a common nonzero translation).
\begin{samepage}
\begin{lemma}[primes producing a common nonzero translation]\leavevmode\label{thm:72}
Let
\(W=\{w_1,\ldots,w_s\}\) be nonempty words with lengths \(\ell_i\) and
\(C_i=C(w_i)\) as in (8), \(g=\gcd(\ell_1,\ldots,\ell_s)\), and
\[
D=\gcd\bigl(a^g-b^g,\ C_ib^{\ell_1}-C_1b^{\ell_i}\ (2\le i\le s)\bigr).
\]
For a prime \(p\nmid ab\), the residue maps
\[
F_i(q)=(a^{\ell_i}q+C_i)/b^{\ell_i}
\]
are all congruent modulo \(p\)
to one and the same translation \(q\mapsto q+\tau\) with \(\tau\ne0\) if
and only if \(p\mid D\) and \(p\nmid C_1\). In that case, no infinite
concatenation of words of \(W\) survives the lift by \(p\): the residues
at block boundaries are \(q,q+\tau,q+2\tau,\ldots\) and vanish within
\(p\) blocks.
\label{thm-end:72}
\end{lemma}
\end{samepage}

% Markdown AST block 589: Proof.
\begin{proof}\leavevmode\label{proof:72}
The slope of \(F_i\) is \((a/b)^{\ell_i}\); all slopes are
\(1\) if and only if the order of \(a/b\) modulo \(p\) divides every
\(\ell_i\), hence divides \(g\), that is, \(p\mid a^g-b^g\).

% Markdown AST block 590: Para
The
translations \(C_i/b^{\ell_i}\) coincide if and only if
\[
C_ib^{\ell_1}\equiv C_1b^{\ell_i}\pmod p,
\]
and the common translation
is nonzero if and only if \(p\nmid C_1\) (\(b\) being a unit). \qedhere
\end{proof}

% Markdown AST block 591: Para
This eliminates only concatenations of \(W\); it removes a component only
when all return words of the component lie in \(W\). For the two-carry
words \(U=\Delta^Lc\), \(V=\Delta^{L+1}c\) of Theorem 42, \(g=1\) and
\(D\mid a-b\); if \(p\mid\Delta\) and \(p\nmid c\), both words act as
\(q\mapsto q+c/b\). For \(101/2\): \(D=99\), \(C_1=69\); the prime
\(3\) divides \(C_1\) and does not work, while \(11\) gives the
translation
\[
\tau=69\cdot2^{-1}\equiv7\pmod{11},
\]
so the two-word
branching that survives \(3\) and \(5\) is eliminated by adding \(11\).
For \(5/2\): \(D=3\), and \(3\) plays the same role. When \(\Delta\in
\{1,2\}\), \(D\mid2\) and no odd prime is obtained this way; this is a
limitation of this particular mechanism.

% Markdown AST block 592: Proposition 73 (whole-word interval test).
\begin{samepage}
\begin{proposition}[whole-word interval test]\leavevmode\label{thm:73}
For a macro edge with
starting cell \([j/K,(j+1)/K)\) and word \(c_0\cdots c_{\ell-1}\),
define
\[
\begin{aligned}
(L_0,U_0,B_0)&=(j,j+1,K)\qquad\text{and}\\
B_{i+1}&=bB_i,\\
L_{i+1}&=\max(0,aL_i-c_iB_i),\\
U_{i+1}&=\min(B_{i+1},aU_i-c_iB_i).
\end{aligned}
\]
Then \([L_i/B_i,U_i/B_i)\) is
exactly the set of fractional parts after \(i\) letters of the real
orbit segments starting in the cell with these carries. Hence rejecting
a macro edge when \(L_i\ge U_i\) for some \(i\) (or when the final
interval misses the target cell) is pseudo-orbit-sound: every
pseudo-orbit, and in particular every true orbit, passing through the
edge has exact fractional parts witnessing nonemptiness. The operation
is not walk-sound, since (E3) expresses only the compatibility of two
adjacent cells, and a path of length \(\ge2\) in \(G(M,K)\) need not
admit real fractional parts.
\label{thm-end:73}
\end{proposition}
\end{samepage}

% Markdown AST block 593: Proof.
\begin{proof}\leavevmode\label{proof:73}
The map
\[
\theta\mapsto(a\theta-c)/b
\]
is an increasing
affine bijection, so the image of \([L/B,U/B)\) is
\[
[(aL-cB)/(bB),(aU-cB)/(bB)),
\]
and intersecting with \([0,1)\) gives
the update; exactness follows by induction as in Proposition 57. \qedhere
\end{proof}

% Markdown AST block 594: Para
Whether adjoining this test changes the class of Theorem 37 is not
determined here; an exploratory computation with it is not adopted,
and the certificates of Section 5 do not use it. The witnesses of
Theorem 42 and Lemma 69 pass this test at every intermediate time,
which distinguishes them from spurious local paths.

% BEGIN PUBLIC ADDITION acknowledgments
\section*{Acknowledgments and disclosure of AI involvement}

This project began with the author's interest in Mills' constant, the least real number \(A>1\) for which \(\lfloor A^{3^n}\rfloor\) is prime for every positive integer \(n\). After learning that this least constant is irrational \cite{SA25}, the author asked \href{https://help.openai.com/en/articles/20001354-gpt-5-6}{GPT-6 Pro in ChatGPT} (powered by GPT-6 Astra) to prove the stronger assertion that every Mills number is irrational. The author estimates having issued roughly five hundred prompts urging the system to continue and not give up. That exploration did not produce the requested proof. The author subsequently asked the AI to identify mathematically meaningful results suitable for a paper among the material obtained, and to write that paper. The present work grew out of that request; it does not resolve the original question about all Mills numbers.

The author thanks GPT-6 Astra, accessed as GPT-6 Pro in ChatGPT and also used in Codex, as the principal AI system in this work. The author also thanks Claude Fable 5.1 for limited assistance through Claude Code with some tasks during preparation of the paper when Codex's weekly usage limit was reached. After the initial conjecture, all mathematical development, identification of results suitable for publication, proof attempts, computation programs, exposition, and Lean formalization were performed by AI systems. The author supplied no mathematically substantive instructions or advice for these tasks. The author's requests for formalization and for an automated workflow in Codex, in which drafting agents and critic agents iterated before typesetting, specified tasks and a procedure without contributing mathematical content. These exchanges were automated critiques, not independent review by human mathematicians.

The author's sole mathematical contribution to this process was the initial conjecture that every Mills number, not only the least one, is irrational. Beyond that conjecture, the author made no mathematical contribution and supplied no ideas for proofs, mathematical arguments, calculations, or mathematical assessment of the results. The author's subsequent involvement consisted of nonmathematical requests to continue working, identify publishable material, write the paper, and produce Lean proofs, together with the decision to release the work. The author has not fully verified the mathematical arguments of the paper, but has personally confirmed that the Lean proofs are accepted by Lean; their coverage and trust assumptions are specified in Section 9.4. The author's belief that the results are likely to be correct and mathematically worthwhile is not a report of a complete verification of the arguments. Readers should assess the arguments, reproducible computations, and formal artifacts themselves.

The author expects that research in which a human poses an initial question and AI carries out all subsequent mathematical or technical work will become increasingly common in mathematics and other disciplines, including cases in which the initiating human neither provides substantive guidance nor fully understands the resulting work. This is the author's personal expectation, not a conclusion established by the mathematical results here. This account makes the provenance of the work, the contributions of the AI systems, and the limits of the author's contribution explicit.
% END PUBLIC ADDITION acknowledgments

% Markdown AST block 595: References
\begin{thebibliography}{SHAO15}

% Markdown AST block 596: reference AFS08
\bibitem[AFS08]{AFS08}\phantomsection\label{ref-afs08}
Shigeki Akiyama, Christiane Frougny, and Jacques Sakarovitch,
\emph{Powers of rationals modulo 1 and rational base number systems},
Israel Journal of Mathematics \textbf{168} (2008), 53--91.
\href{https://doi.org/10.1007/s11856-008-1056-4}{Published article}.

% Markdown AST block 597: reference AD04
\bibitem[AD04]{AD04}\phantomsection\label{ref-ad04}
Giedrius Alkauskas and Artūras Dubickas, \emph{Prime and composite
numbers as integer parts of powers}, Acta Mathematica Hungarica
\textbf{105} (2004), no. 3, 249--256.
\href{https://doi.org/10.1023/B:AMHU.0000049291.38284.2c}{Published article}.

% Markdown AST block 598: reference AEV26
\bibitem[AEV26]{AEV26}\phantomsection\label{ref-aev26}
Mélodie Andrieu, Shalom Eliahou, and Léo Vivion, \emph{A
Normality Conjecture on Rational Base Number Systems}, arXiv:2510.11723
{[}math.NT{]}, version 2, April 7, 2026 (version 1: October 6, 2025).
\href{https://arxiv.org/abs/2510.11723v2}{Preprint}.

% Markdown AST block 599: reference CW13
\bibitem[CW13]{CW13}\phantomsection\label{ref-cw13}
Fintan Costello and Paul Watts, \emph{A short note on
Jacobsthal's function}, arXiv:1306.1064 {[}math.NT{]}, version 1, June 5,
2013. \href{https://arxiv.org/abs/1306.1064v1}{Preprint}.

% Markdown AST block 600: reference Dub06
\bibitem[Dub06]{Dub06}\phantomsection\label{ref-dub06}
Artūras Dubickas, \emph{Truncatable primes and unavoidable sets of
divisors}, Acta Mathematica Universitatis Ostraviensis \textbf{14}
(2006), 21--25.
\href{https://klevas.mif.vu.lt/~dubickas/files/dvifai/cekko1.pdf}{Author's copy}.

% Markdown AST block 600: reference Dub08
\bibitem[Dub08]{Dub08}\phantomsection\label{ref-dub08}
Artūras Dubickas, \emph{On the powers of 3/2 and other rational
numbers}, Mathematische Nachrichten \textbf{281} (2008), no. 7, 951--958.
\href{https://doi.org/10.1002/mana.200510651}{Published article}.

% Markdown AST block 600: reference Dub09
\bibitem[Dub09]{Dub09}\phantomsection\label{ref-dub09}
Artūras Dubickas, \emph{Prime and composite integers close to powers of
a number}, Monatshefte für Mathematik \textbf{158} (2009), no. 3,
271--284. \href{https://doi.org/10.1007/s00605-008-0042-6}{Published article}.

% Markdown AST block 601: reference DJ22
\bibitem[DJ22]{DJ22}\phantomsection\label{ref-dj22}
Artūras Dubickas and Lukas Jonuška, \emph{Divisibility of integers
obtained from truncated periodic sequences}, International Journal of
Number Theory \textbf{18} (2022), no. 1, 165--174.
\href{https://doi.org/10.1142/S1793042122500129}{Published article}.

% Markdown AST block 602: reference DN
\bibitem[DN]{DN}\phantomsection\label{ref-dn}
Artūras Dubickas and Aivaras Novikas, \emph{Integer parts of powers of
rational numbers}, Mathematische Zeitschrift \textbf{251} (2005), no. 3,
635--648. \href{https://doi.org/10.1007/s00209-005-0827-4}{Published article}.

% Markdown AST block 603: reference FLP
\bibitem[FLP]{FLP}\phantomsection\label{ref-flp}
Leopold Flatto, Jeffrey C. Lagarias, and Andrew D. Pollington,
\emph{On the range of fractional parts \(\{\xi(p/q)^n\}\)},
Acta Arithmetica \textbf{70} (1995), no. 2, 125--147.
\href{https://doi.org/10.4064/aa-70-2-125-147}{Published article}.

% Markdown AST block 604: reference FS
\bibitem[FS]{FS}\phantomsection\label{ref-fs}
W. Forman and H. N. Shapiro, \emph{An arithmetic property of certain
rational powers}, Communications on Pure and Applied Mathematics
\textbf{20} (1967), no. 3, 561--573.
\href{https://doi.org/10.1002/cpa.3160200305}{Published article}.

% Markdown AST block 605: reference IW78
\bibitem[IW78]{IW78}\phantomsection\label{ref-iw78}
Henryk Iwaniec, \emph{On the problem of Jacobsthal}, Demonstratio
Mathematica \textbf{11} (1978), no. 1, 225--231.
\href{https://doi.org/10.1515/dema-1978-0121}{Published article}.

% Markdown AST block 606: reference M68
\bibitem[M68]{M68}\phantomsection\label{ref-m68}
K. Mahler, \emph{An unsolved problem on the powers of 3/2},
Journal of the Australian Mathematical Society \textbf{8} (1968), no. 2,
313--321. \href{https://doi.org/10.1017/S1446788700005371}{Published article}.

% Markdown AST block 607: reference N
\bibitem[N]{N}\phantomsection\label{ref-n}
Aivaras Novikas, \emph{Composite numbers in the sequences of integers},
doctoral dissertation, Vilnius University, 2012.
\href{https://epublications.vu.lt/object/elaba:2089125/index.html}{University repository}.

% Markdown AST block 608: reference SHAO15
\bibitem[SHAO15]{SHAO15}\phantomsection\label{ref-shao15}
Xuancheng Shao, \emph{Narrow arithmetic progressions in the
primes}, arXiv:1509.04955 {[}math.NT{]}, version 1, September 16, 2015;
published in International Mathematics Research Notices \textbf{2017}, no. 2,
391--428. \href{https://arxiv.org/abs/1509.04955v1}{Preprint},
\href{https://doi.org/10.1093/imrn/rnv393}{Published article}.

% Markdown AST block 609: reference S26
\bibitem[S26]{S26}\phantomsection\label{ref-s26}
Ralf Stephan, \emph{On the composites among \([\xi 7^n]\)}, author-uploaded
preprint dated August 15, 2026.
\href{https://www.researchgate.net/publication/412294027_On_the_composites_among_x7}{Preprint};
\href{https://github.com/rwst/On-Composites}{Lean files}.

% Markdown AST block 610: reference TZ14
\bibitem[TZ14]{TZ14}\phantomsection\label{ref-tz14}
Terence Tao and Tamar Ziegler, \emph{Narrow progressions in the
primes}, arXiv:1409.1327 {[}math.NT{]}, version 2, October 10, 2014;
published in \emph{Analytic Number Theory} (Springer, 2015), 357--379.
\href{https://arxiv.org/abs/1409.1327v2}{Preprint},
\href{https://doi.org/10.1007/978-3-319-22240-0_22}{Published chapter}.

% Markdown AST block 611: reference T
\bibitem[T]{T}\phantomsection\label{ref-t}
Robert Tarjan, \emph{Depth-First Search and Linear Graph Algorithms},
SIAM Journal on Computing \textbf{1} (1972), no. 2, 146--160.
\href{https://doi.org/10.1137/0201010}{Published article}.

% BEGIN PUBLIC ADDITION reference-sa25
\bibitem[SA25]{SA25}\phantomsection\label{ref-sa25}
Kota Saito, \emph{Mills' constant is irrational},
Mathematika \textbf{71} (2025), no. 3, e70027.
\href{https://doi.org/10.1112/mtk.70027}{Published article},
\href{https://arxiv.org/abs/2404.19461}{Preprint}.
% END PUBLIC ADDITION reference-sa25

\end{thebibliography}
\end{document}
